% GENERATED FILE -- DO NOT EDIT.
%
% Produced by harness/assemble_contract.py from module.yaml and claims/*.tex.
% Edit the claim files and regenerate; `--check` fails if this file is stale.
\documentclass[11pt,a4paper]{article}
\pdftrailerid{b0b708a3755eccb572cb7bec24764a7a}
%% --- 1. Mathematics (order matters: physics after braket wins \ket, \tr) ---
\usepackage{amsmath}
\usepackage{amssymb}
\usepackage{amsfonts}
\usepackage{amsthm}
\usepackage{mathtools}
\usepackage{dsfont}
\usepackage{braket}
\usepackage{physics}
%% --- 2. Page furniture ------------------------------------------------------
\usepackage{xcolor}
\usepackage{array}    % >{...} column prefixes in the wide tables of some claims
\usepackage{enumitem}
\usepackage{graphicx}
\usepackage{booktabs}
\usepackage{longtable}
\usepackage{xltabular}              % front- and back-matter tables across pages
\usepackage{needspace}
\usepackage{xurl}
\usepackage{tcolorbox}
\tcbuselibrary{skins,breakable}
%% --- 3. Links and navigation (hyperref late; bookmark, refcount after it) ---
\usepackage{xr-hyper}   % before hyperref: the contract's labels
\usepackage{hyperref}
\usepackage{bookmark}
\usepackage{refcount}
%% --- 4. Colours: every distinction also has a shape (rule, caps, size, weight)
\definecolor{nnGrey}{gray}{0.40}      % non-normative text
\definecolor{metaGrey}{gray}{0.30}    % the generated meta line
\definecolor{techGrey}{gray}{0.45}    % heading of a technical claim
\definecolor{mainRule}{RGB}{0,110,140}
\definecolor{reqAccent}{RGB}{165,70,0}
\definecolor{linkBlue}{RGB}{0,62,150}
\definecolor{extLink}{RGB}{120,40,140}  % links into another PDF
\hypersetup{colorlinks=true, linkcolor=linkBlue, citecolor=linkBlue,
  urlcolor=linkBlue, filecolor=extLink, bookmarksnumbered=true,
  bookmarksopen=true, bookmarksopenlevel=1, pdfdisplaydoctitle=true}
% A link into another PDF (\href to a .pdf, \hyperref{file}..., xr-hyper) is a
% relative URI, <file>#nameddest=<destination>, in the file-link colour:
% browser viewers (Edge, Chrome) ignore hyperref's default GoToR action.
\makeatletter
\renewcommand*{\hyper@linkfile}[3]{% anchor text, file name, destination
  \begingroup
    \ifx\\#3\\\def\CQ@target{#2}\else\def\CQ@target{#2\hyper@hash nameddest=#3}\fi
    \let\@urlcolor\@filecolor
    \hyper@linkurl{#1}{\CQ@target}%
  \endgroup}
\makeatother
%% --- 5. The mapped Unicode set (harness/contract_navigation.py, UNICODE_TEX)
\DeclareUnicodeCharacter{03B1}{\ensuremath{\alpha}}
\DeclareUnicodeCharacter{03B2}{\ensuremath{\beta}}
\DeclareUnicodeCharacter{03B3}{\ensuremath{\gamma}}
\DeclareUnicodeCharacter{03B4}{\ensuremath{\delta}}
\DeclareUnicodeCharacter{03B5}{\ensuremath{\varepsilon}}
\DeclareUnicodeCharacter{03F5}{\ensuremath{\epsilon}}
\DeclareUnicodeCharacter{03B6}{\ensuremath{\zeta}}
\DeclareUnicodeCharacter{03B7}{\ensuremath{\eta}}
\DeclareUnicodeCharacter{03B8}{\ensuremath{\theta}}
\DeclareUnicodeCharacter{03D1}{\ensuremath{\vartheta}}
\DeclareUnicodeCharacter{03B9}{\ensuremath{\iota}}
\DeclareUnicodeCharacter{03BA}{\ensuremath{\kappa}}
\DeclareUnicodeCharacter{03BB}{\ensuremath{\lambda}}
\DeclareUnicodeCharacter{03BC}{\ensuremath{\mu}}
\DeclareUnicodeCharacter{03BD}{\ensuremath{\nu}}
\DeclareUnicodeCharacter{03BE}{\ensuremath{\xi}}
\DeclareUnicodeCharacter{03C0}{\ensuremath{\pi}}
\DeclareUnicodeCharacter{03C1}{\ensuremath{\rho}}
\DeclareUnicodeCharacter{03F1}{\ensuremath{\varrho}}
\DeclareUnicodeCharacter{03C3}{\ensuremath{\sigma}}
\DeclareUnicodeCharacter{03C2}{\ensuremath{\varsigma}}
\DeclareUnicodeCharacter{03C4}{\ensuremath{\tau}}
\DeclareUnicodeCharacter{03C5}{\ensuremath{\upsilon}}
\DeclareUnicodeCharacter{03C6}{\ensuremath{\varphi}}
\DeclareUnicodeCharacter{03D5}{\ensuremath{\phi}}
\DeclareUnicodeCharacter{03C7}{\ensuremath{\chi}}
\DeclareUnicodeCharacter{03C8}{\ensuremath{\psi}}
\DeclareUnicodeCharacter{03C9}{\ensuremath{\omega}}
\DeclareUnicodeCharacter{0393}{\ensuremath{\Gamma}}
\DeclareUnicodeCharacter{0394}{\ensuremath{\Delta}}
\DeclareUnicodeCharacter{0398}{\ensuremath{\Theta}}
\DeclareUnicodeCharacter{039B}{\ensuremath{\Lambda}}
\DeclareUnicodeCharacter{039E}{\ensuremath{\Xi}}
\DeclareUnicodeCharacter{03A0}{\ensuremath{\Pi}}
\DeclareUnicodeCharacter{03A3}{\ensuremath{\Sigma}}
\DeclareUnicodeCharacter{03A5}{\ensuremath{\Upsilon}}
\DeclareUnicodeCharacter{03A6}{\ensuremath{\Phi}}
\DeclareUnicodeCharacter{03A8}{\ensuremath{\Psi}}
\DeclareUnicodeCharacter{03A9}{\ensuremath{\Omega}}
\DeclareUnicodeCharacter{00E1}{\'a}
\DeclareUnicodeCharacter{00E0}{\`a}
\DeclareUnicodeCharacter{00E2}{\^a}
\DeclareUnicodeCharacter{00E4}{\"a}
\DeclareUnicodeCharacter{00E5}{\aa{}}
\DeclareUnicodeCharacter{00E9}{\'e}
\DeclareUnicodeCharacter{00E8}{\`e}
\DeclareUnicodeCharacter{00EA}{\^e}
\DeclareUnicodeCharacter{00EB}{\"e}
\DeclareUnicodeCharacter{00ED}{\'{\i}}
\DeclareUnicodeCharacter{00EF}{\"{\i}}
\DeclareUnicodeCharacter{00F3}{\'o}
\DeclareUnicodeCharacter{00F2}{\`o}
\DeclareUnicodeCharacter{00F4}{\^o}
\DeclareUnicodeCharacter{00F6}{\"o}
\DeclareUnicodeCharacter{00F8}{\o{}}
\DeclareUnicodeCharacter{00FA}{\'u}
\DeclareUnicodeCharacter{00F9}{\`u}
\DeclareUnicodeCharacter{00FB}{\^u}
\DeclareUnicodeCharacter{00FC}{\"u}
\DeclareUnicodeCharacter{00F1}{\~n}
\DeclareUnicodeCharacter{00E7}{\c{c}}
\DeclareUnicodeCharacter{00DF}{\ss{}}
\DeclareUnicodeCharacter{00C4}{\"A}
\DeclareUnicodeCharacter{00C5}{\AA{}}
\DeclareUnicodeCharacter{00C9}{\'E}
\DeclareUnicodeCharacter{00D6}{\"O}
\DeclareUnicodeCharacter{00DC}{\"U}
\DeclareUnicodeCharacter{2013}{\textendash{}}
\DeclareUnicodeCharacter{2014}{\textemdash{}}
\DeclareUnicodeCharacter{2018}{`}
\DeclareUnicodeCharacter{2019}{'}
\DeclareUnicodeCharacter{201C}{``}
\DeclareUnicodeCharacter{201D}{''}
\DeclareUnicodeCharacter{2026}{\ldots{}}
\DeclareUnicodeCharacter{00A0}{~}
\DeclareUnicodeCharacter{00A7}{\S{}}
\DeclareUnicodeCharacter{2264}{\ensuremath{\le}}
\DeclareUnicodeCharacter{2265}{\ensuremath{\ge}}
\DeclareUnicodeCharacter{2260}{\ensuremath{\neq}}
\DeclareUnicodeCharacter{2248}{\ensuremath{\approx}}
\DeclareUnicodeCharacter{2261}{\ensuremath{\equiv}}
\DeclareUnicodeCharacter{223C}{\ensuremath{\sim}}
\DeclareUnicodeCharacter{2243}{\ensuremath{\simeq}}
\DeclareUnicodeCharacter{2245}{\ensuremath{\cong}}
\DeclareUnicodeCharacter{221D}{\ensuremath{\propto}}
\DeclareUnicodeCharacter{226A}{\ensuremath{\ll}}
\DeclareUnicodeCharacter{226B}{\ensuremath{\gg}}
\DeclareUnicodeCharacter{2192}{\ensuremath{\to}}
\DeclareUnicodeCharacter{2190}{\ensuremath{\leftarrow}}
\DeclareUnicodeCharacter{2194}{\ensuremath{\leftrightarrow}}
\DeclareUnicodeCharacter{21D2}{\ensuremath{\Rightarrow}}
\DeclareUnicodeCharacter{21D0}{\ensuremath{\Leftarrow}}
\DeclareUnicodeCharacter{21D4}{\ensuremath{\Leftrightarrow}}
\DeclareUnicodeCharacter{21A6}{\ensuremath{\mapsto}}
\DeclareUnicodeCharacter{2208}{\ensuremath{\in}}
\DeclareUnicodeCharacter{2209}{\ensuremath{\notin}}
\DeclareUnicodeCharacter{2286}{\ensuremath{\subseteq}}
\DeclareUnicodeCharacter{2282}{\ensuremath{\subset}}
\DeclareUnicodeCharacter{2287}{\ensuremath{\supseteq}}
\DeclareUnicodeCharacter{2283}{\ensuremath{\supset}}
\DeclareUnicodeCharacter{2229}{\ensuremath{\cap}}
\DeclareUnicodeCharacter{222A}{\ensuremath{\cup}}
\DeclareUnicodeCharacter{2216}{\ensuremath{\setminus}}
\DeclareUnicodeCharacter{2205}{\ensuremath{\emptyset}}
\DeclareUnicodeCharacter{2200}{\ensuremath{\forall}}
\DeclareUnicodeCharacter{2203}{\ensuremath{\exists}}
\DeclareUnicodeCharacter{2115}{\ensuremath{\mathbb{N}}}
\DeclareUnicodeCharacter{2124}{\ensuremath{\mathbb{Z}}}
\DeclareUnicodeCharacter{211A}{\ensuremath{\mathbb{Q}}}
\DeclareUnicodeCharacter{211D}{\ensuremath{\mathbb{R}}}
\DeclareUnicodeCharacter{2102}{\ensuremath{\mathbb{C}}}
\DeclareUnicodeCharacter{00B7}{\ensuremath{\cdot}}
\DeclareUnicodeCharacter{00D7}{\ensuremath{\times}}
\DeclareUnicodeCharacter{2212}{\ensuremath{-}}
\DeclareUnicodeCharacter{00B1}{\ensuremath{\pm}}
\DeclareUnicodeCharacter{2213}{\ensuremath{\mp}}
\DeclareUnicodeCharacter{221A}{\ensuremath{\surd}}
\DeclareUnicodeCharacter{2211}{\ensuremath{\sum}}
\DeclareUnicodeCharacter{220F}{\ensuremath{\prod}}
\DeclareUnicodeCharacter{222B}{\ensuremath{\int}}
\DeclareUnicodeCharacter{221E}{\ensuremath{\infty}}
\DeclareUnicodeCharacter{2202}{\ensuremath{\partial}}
\DeclareUnicodeCharacter{2207}{\ensuremath{\nabla}}
\DeclareUnicodeCharacter{2297}{\ensuremath{\otimes}}
\DeclareUnicodeCharacter{2295}{\ensuremath{\oplus}}
\DeclareUnicodeCharacter{2218}{\ensuremath{\circ}}
\DeclareUnicodeCharacter{2308}{\ensuremath{\lceil}}
\DeclareUnicodeCharacter{2309}{\ensuremath{\rceil}}
\DeclareUnicodeCharacter{230A}{\ensuremath{\lfloor}}
\DeclareUnicodeCharacter{230B}{\ensuremath{\rfloor}}
\DeclareUnicodeCharacter{27E8}{\ensuremath{\langle}}
\DeclareUnicodeCharacter{27E9}{\ensuremath{\rangle}}
\DeclareUnicodeCharacter{2016}{\ensuremath{\|}}
\DeclareUnicodeCharacter{2113}{\ensuremath{\ell}}
\DeclareUnicodeCharacter{00B0}{\ensuremath{^{\circ}}}
\DeclareUnicodeCharacter{00B2}{\ensuremath{^{2}}}
\DeclareUnicodeCharacter{00B3}{\ensuremath{^{3}}}
\DeclareUnicodeCharacter{00B9}{\ensuremath{^{1}}}
\DeclareUnicodeCharacter{2080}{\ensuremath{_{0}}}
\DeclareUnicodeCharacter{2081}{\ensuremath{_{1}}}
\DeclareUnicodeCharacter{2082}{\ensuremath{_{2}}}
\DeclareUnicodeCharacter{207F}{\ensuremath{^{n}}}
\DeclareUnicodeCharacter{2096}{\ensuremath{_{k}}}
%% --- 6. Numbering: one counter per section -----------------------
\numberwithin{equation}{section}
\setcounter{tocdepth}{1}
%% --- 7. The project macro set, then the claim environments -----------------
% spec/macros.tex, read relative to this module's directory (latexmk -cd).
% =============================================================================
%  spec/macros.tex -- the one macro set every generated contract inputs.
%
%  Two parts: part A, the macros spec/notation.md catalogues, and part B,
%  the macros local to one module, grouped by module.
%  harness/assemble_contract.py emits % =============================================================================
%  spec/macros.tex -- the one macro set every generated contract inputs.
%
%  Two parts: part A, the macros spec/notation.md catalogues, and part B,
%  the macros local to one module, grouped by module.
%  harness/assemble_contract.py emits % =============================================================================
%  spec/macros.tex -- the one macro set every generated contract inputs.
%
%  Two parts: part A, the macros spec/notation.md catalogues, and part B,
%  the macros local to one module, grouped by module.
%  harness/assemble_contract.py emits \input{../../macros} for it, so the
%  file is read from spec/modules/<id>/.
%
%  RULES
%    * A macro that spec/notation.md catalogues is a verbatim copy of it
%      (part A). Changing one is a notation change: bump spec/notation.md.
%    * Part B holds the module-local macros. A macro defined by two modules
%      with the same expansion appears once, under the first module (in the
%      order below) that defines it. Two macros with the same name and
%      DIFFERENT expansions are both kept under a module prefix, and the
%      claims write the prefixed names:
%        \Chain  -> \redChain (reductions, \mathfrak{C}: the reduction chain)
%                   \pcChain  (protocol-compiler, \mathfrak{R}: the two-segment
%                             chain; \Chsrc and \Chrec are built on it)
%        \Nmsm   -> \redNmsm  (reductions, \mathcal{N}: a POVM measurement
%                             channel)
%                   \pcNmsm   (\mathcal{N}_{\mathsf A}: Alice's measurement
%                             channel, written by
%                             transmitter/thm:certificate-sufficiency,
%                             transmitter/rem:xreg and
%                             reductions/lem:source-transfer)
%    * Deliberately omitted: a \provenance{tag}{details} macro. Provenance
%      is a header field of every claim file (spec/claim-format.md section
%      3); a claim body that calls such a macro fails to compile, which is
%      the intended signal.
%
%  REQUIRED PACKAGES, loaded by the assembler's preamble before this file:
%    amsmath, amssymb, amsfonts, mathtools, dsfont, braket, physics, xcolor
%    (braket before physics, so that physics wins for \ket, \bra, \tr, ...).
%
%  This file contains only definitions: no \usepackage, no document body.
% =============================================================================


% =============================================================================
%  PART A -- notation macros, the ones spec/notation.md catalogues, each a
%  verbatim copy of its catalogue row; grouped under numbered headings.
% =============================================================================

% -----------------------------------------------------------------------------
% 1. State spaces and operator sets (a "system label" -> spaces of operators)
% -----------------------------------------------------------------------------
\newcommand{\Hs}[1]{\mathcal{H}(#1)}        % Hilbert space of a system label
\newcommand{\Lin}[1]{\mathcal{L}(#1)}       % bounded linear operators on H(#1)
\newcommand{\Linn}[2]{\mathcal{L}(#1,#2)}   % linear maps H(#1) -> H(#2)
\newcommand{\Pos}[1]{\mathrm{Pos}(#1)}      % positive semidefinite operators on H(#1)
\newcommand{\Density}{\mathrm{D}}           % density operators: \Density(A) = D(A)
\newcommand{\Seq}[1]{S_{=}(#1)}             % normalized states (trace = 1)
\newcommand{\Sleq}[1]{S_{\leq}(#1)}         % subnormalized states (0 <= trace <= 1)
\newcommand{\id}{\mathds{1}}                % identity operator
\newcommand{\idon}[1]{\mathds{1}_{#1}}      % identity on a named system, e.g. \idon{C}

% DEPRECATED: applies \mathcal to both arguments itself, so it takes bare
% letters, not system labels. Prefer \Linn; do not use in new contracts.
\providecommand{\Lop}[2]{\mathcal{L}(\mathcal{#1}, \mathcal{#2})}


% -----------------------------------------------------------------------------
% 2. Composite systems
% -----------------------------------------------------------------------------
\newcommand{\dsum}{\oplus}                  % direct sum of system labels
% Tensor product: use \otimes directly.
% n-fold direct sum of an operator space: \Lin{A}^{\oplus n}


% -----------------------------------------------------------------------------
% 3. Events, conditioning and classical registers (for STATES)
%    Convention:  \wedge  -> subnormalized "and-event" component, e.g. \rho_{A\wedge x}
%                 |        -> normalized conditional state,        e.g. \rho_{A|x}
% (Write these directly in subscripts; no extra macros needed.)
% -----------------------------------------------------------------------------


% -----------------------------------------------------------------------------
% 4. Photon-number subspaces / direct-sum blocks
%    Convention:  operators carry the restriction "on top", in brackets;
%                 states carry it as a \wedge subscript (see Sec. 3).
% -----------------------------------------------------------------------------
\newcommand{\Pnum}[1]{P_{#1}}               % photon-number projector (system from context)
\newcommand{\Pnumon}[2]{P_{#1}^{#2}}        % with explicit system superscript, e.g. P_n^C
\newcommand{\blk}[2]{#1^{(#2)}}             % operator block, e.g. \blk{M}{\leqN} = M^{(<=N)}
\newcommand{\leqN}{\leq N}                  % use as \blk{M}{\leqN}
\newcommand{\gtN}{> N}                      % use as \blk{M}{\gtN}
\newcommand{\eqn}{= n}                      % use as \blk{M}{\eqn}


% -----------------------------------------------------------------------------
% 5. System labels used in the protocol
% -----------------------------------------------------------------------------
\newcommand{\Chan}{C}                       % quantum-channel (Eve) input to Bob
\newcommand{\Prb}{P}                        % probe input
\newcommand{\Flag}{{\mathrm{Flg}}}          % flag system (EXISTS ONLY AFTER squashing)
\newcommand{\Loss}{{\mathrm{Lss}}}          % loss system
\newcommand{\Vc}{{\mathrm{Vcm}}}            % vacuum system
\newcommand{\Nout}{N_{\mathrm{out}}}
\newcommand{\Nin}{N_{\mathrm{in}}}


% -----------------------------------------------------------------------------
% 6. Alphabets
% -----------------------------------------------------------------------------
\newcommand{\Yset}{\mathcal{Y}}             % Bob's (coarse-grained) detection alphabet
\newcommand{\Yraw}{\mathcal{Y}_{\mathrm{raw}}}% raw click patterns {0,1}^4
\newcommand{\Yscset}{\mathcal{Y}_{\singlec}}% single-click outcomes
\newcommand{\Ymcset}{\mathcal{Y}_{\mc}}     % multi-click outcomes
\newcommand{\Tset}{\mathcal{T}_{\mathrm{set}}}% private setting set (the probe settings in the optical instance)
\newcommand{\Xset}{\mathcal{X}}             % Alice's alphabet
\newcommand{\Dset}{\{H,V,+,-\}}             % polarization / detector labels
\newcommand{\pass}{\top}                    % pass sifting


% -----------------------------------------------------------------------------
% 7. Mode / outcome abbreviations
% -----------------------------------------------------------------------------
\newcommand{\on}{{\mathrm{on}}}
\newcommand{\off}{{\mathrm{off}}}
\newcommand{\dc}{{\mathrm{dc}}}             % dark count
\newcommand{\sq}{\mathrm{sq}}               % squashed
\newcommand{\mc}{\mathrm{mc}}               % multi-click
\newcommand{\nc}{\mathrm{nc}}               % no-click
\newcommand{\singlec}{\mathrm{sc}}          % single-click
\newcommand{\vac}{\mathrm{vac}}             % vacuum
\newcommand{\kout}{{k_\mathrm{out}}}
\newcommand{\kin}{{k_\mathrm{in}}}


% -----------------------------------------------------------------------------
% 8. Maps, channels and miscellaneous operators
% -----------------------------------------------------------------------------
\newcommand{\Prob}{\mathrm{Prob}}           % probability (use \Prob[\cdot] or \Prob(\cdot))
\newcommand{\lmin}{\lambda_{\min}}          % minimum eigenvalue
\newcommand{\Wflag}{W_{\Flag}}              % flag-space weight
\newcommand{\transpose}{\mathsf{T}}
\newcommand{\innerp}[2]{\left\langle #1 \vert #2 \right\rangle}
\renewcommand{\vec}[1]{\boldsymbol{#1}}     % NB: overrides the LaTeX arrow accent


% -----------------------------------------------------------------------------
% 9. Key-rate numerics
% -----------------------------------------------------------------------------
\newcommand{\cG}{\mathcal{G}}               % protocol map
\newcommand{\cZ}{\mathcal{Z}}               % pinching map
\newcommand{\Amap}{\mathcal{A}_{\mathrm{ann}}} % announcement + measurement map
\newcommand{\Amapnb}{\mathcal{A}_{\mathrm{ann}|\mathrm{nb}}}
\newcommand{\simpl}{}                       % empty suffix for simplified block maps
\newcommand{\acSet}{\mathcal{S}_{\mathrm{acc}}} % acceptance set
\newcommand{\testE}{\mathcal{E}_{\mathrm{test}}} % test-statistics map
\newcommand{\ecLeak}{\mathrm{leak}^{\mathrm{EC}}_{\mathrm{obs}}}
\newcommand{\pgen}{p_{\mathrm{gen}}}
\newcommand{\askeyr}{R_{\infty}}             % asymptotic key rate
\newcommand{\ppass}{p_{\mathrm{pass}}}       % sifting pass probability


% -----------------------------------------------------------------------------
% 10. Finite-size analysis / Renyi entropies
% -----------------------------------------------------------------------------

% --- Round types: test vs. key generation (register T, alphabet \Fset) --------
\newcommand{\tst}{\mathrm{test}}            % test round
\newcommand{\gen}{\mathrm{gen}}             % key-generation round
\newcommand{\Fset}{\mathcal{T}}             % round-type alphabet {test, gen} (not \Tset)
\newcommand{\Chat}{\widehat{C}}             % acceptance-test register
\newcommand{\Chatset}{\widehat{\mathcal{C}}}% its alphabet

% --- Basis / key extraction functions ----------------------------------------
\newcommand{\bA}{b_{A}}                     % Alice's basis-extraction function
\newcommand{\bB}{b_{B}}                     % Bob's basis-extraction function
\newcommand{\kB}{k_{B}}                     % Bob's key-extraction function

% --- Maps --------------------------------------------------------------------
\newcommand{\Mprot}{\mathcal{M}}            % single-round protocol map
\newcommand{\Phimap}{\boldsymbol{\Phi}}     % test-statistics map
\newcommand{\Pdeph}{\mathcal{P}}            % dephasing (pinching) map on \Blreg

% --- Isometries --------------------------------------------------------------
\newcommand{\Vmsm}{V_{\mathrm{msm}}}        % measurement isometry
\newcommand{\Vann}{V_{\mathrm{ann}}}        % announcement isometry
\newcommand{\VI}{V_{I}}                     % Eve's copy of the announcements
\newcommand{\VChat}{V_{\Chat}}              % acceptance-data isometry
\newcommand{\VR}{V_{R}}                     % sifting + key map -> raw key R
\newcommand{\VZ}{V_{Z}}                     % raw key -> key register Z
\newcommand{\Wprot}{W}                      % full protocol isometry

% --- Renyi entropies ---------------------------------------------------------
% SECURITY-CRITICAL: see spec/notation.md §4. Never change an arrow to
% make an inequality work out; never convert silently between the sandwiched
% (tilde) and Petz (overline) families.
\newcommand{\Hsandup}{\widetilde{H}^{\uparrow}_{\alpha}}   % sandwiched, up
\newcommand{\Hsanddown}{\widetilde{H}^{\downarrow}_{\alpha}}% sandwiched, down
\newcommand{\Hpetzup}{\overline{H}^{\uparrow}_{\alpha}}    % Petz, up
\newcommand{\Hpetzdown}{\overline{H}^{\downarrow}_{\alpha}}% Petz, down
\newcommand{\hup}{h^{\uparrow}_{\alpha}}                   % single-round rate
\newcommand{\hqkd}{h^{\uparrow,\mathrm{QKD}}_{\alpha}}     % single-round QKD rate
\newcommand{\hlow}{\underline{h}_{\alpha}}                 % lower bound on \hqkd
% --- spec/notation.md 3.5 and 4.5 -------------------------------------------
% \hlb  is the first-order linearization lower bound on \hlow: the headline
%       key rate (the `first-order` tier of foundations/def:tiers). NOT \hlow
%       itself.
% \tnum is the instance's one declared numerical tolerance. Its
%       near-collision with the declared marginal \tau_A is disclosed in
%       spec/notation.md 3.5; bare \tau is reserved and denotes neither.
\newcommand{\hlb}{\underline{h}^{\mathrm{FO}}_{\alpha}}    % first-order lower bound on \hlow
\newcommand{\tnum}{\tau_{\mathrm{num}}}                    % declared numerical tolerance


% =============================================================================
%  PART B -- module-local macros, grouped by module. Deprecated aliases,
%  defined so that a text written with them compiles, used in no new contract: \Lop,
%  \testE, \Pdeph, \pass, \ecLeak (spec/notation.md section 3.9).
% =============================================================================

% --- meat-core -----------------------------------------------------------------
\newcommand{\SecSet}{\mathbb{S}}            % sector index set of Q
\newcommand{\Qsys}{Q}                       % declared sector-carrying measured system
\newcommand{\Vset}{\mathcal{V}}             % declared joint outcome alphabet of the Q-POVM
\newcommand{\CAset}{\mathcal{C}_{A}}        % Alice announcement alphabet
\newcommand{\CBset}{\mathcal{C}_{B}}        % Bob announcement alphabet
\newcommand{\Party}{\mathfrak{P}}           % declared party partition
\newcommand{\Oset}[1]{\mathcal{O}_{#1}}     % readout alphabet of party #1
\newcommand{\rd}[1]{\pi_{#1}}               % readout map of party #1
\newcommand{\Sadm}{\mathcal{S}_{A}}         % admissible Alice subspace
\newcommand{\nomarg}{\mathsf{none}}         % "no trusted marginal declared"
\newcommand{\Zset}{\mathcal{Z}_{\perp}}     % raw-key alphabet incl. \perp
\newcommand{\Zval}{\mathcal{Z}}             % key alphabet without \perp
\newcommand{\Etil}{\widetilde{E}}           % single-round purifying register
\newcommand{\Ehat}{\widehat{E}}             % Eve's n-round register
\newcommand{\OmAT}{\Omega_{\mathrm{AT}}}    % acceptance event
\newcommand{\epa}{\varepsilon_{\mathrm{PA}}}
\newcommand{\eev}{\varepsilon_{\mathrm{EV}}}
\newcommand{\lEC}{\lambda_{\mathrm{EC}}}
\newcommand{\Decl}{\mathfrak{D}}            % protocol declaration
\newcommand{\CList}{\mathfrak{L}}           % declared constraint list
\newcommand{\PinS}{\mathcal{P}_{\SecSet}}   % sector pinching on Q
\newcommand{\Feas}{\mathcal{W}}             % feasible set of block states
\newcommand{\Chkey}{\Chatset_{\mathrm{key}}}% key-live acceptance values
\newcommand{\Kann}{\mathcal{K}}             % announcement-class alphabet
\newcommand{\hold}{\mathsf{h}}              % declared key holder (A or B)
\newcommand{\sfD}{D}                        % classical KL divergence (base 2)
\newcommand{\Dtwo}{D}                       % quantum relative entropy (base 2)
\newcommand{\Dsand}{\widetilde{D}_{\alpha}} % sandwiched Renyi divergence
\newcommand{\idmap}{\mathrm{id}}            % identity channel
\newcommand{\Aev}{\mathcal{A}}              % Eve's per-round attack channel
\newcommand{\AttSet}{\mathfrak{A}}          % licensed attack-channel set
\newcommand{\Chacc}{\Chatset_{\mathrm{acc}}}% acceptance support of \acSet
\newcommand{\epsstar}{\varepsilon^{\star}}      % distance to an exact POVM
\newcommand{\dfH}{\Delta_{\mathrm{herm}}}       % Hermiticity defect
\newcommand{\dfP}{\Delta_{\mathrm{pos}}}        % negative part
\newcommand{\dfC}{\Delta_{\mathrm{cpl}}}        % completeness defect
\newcommand{\lmoq}{g}                       % acceptance-side gradient block
\newcommand{\lmoJ}{G}                       % feasible-side gradient block
\newcommand{\bacc}{\beta_{\mathrm{acc}}}    % certified acceptance-side minimum
\newcommand{\Itset}{\mathfrak{I}}           % finite index set of iterates (C22)

% --- reductions ----------------------------------------------------------------
\newcommand{\Bsys}{B}                        % physical (pre-squashing) system
\newcommand{\Tcl}[1]{\mathcal{T}(#1)}        % trace-class operators on H(#1)
\newcommand{\Ins}{\mathcal{I}}               % generalized-instrument branch map
\newcommand{\redNmsm}{\mathcal{N}}              % POVM measurement channel
\newcommand{\Eprot}{\mathcal{E}}             % n-round protocol channel
\newcommand{\Grnd}{\mathsf{G}}               % one round map of a protocol
\newcommand{\Kfin}{\mathsf{K}}               % end-of-protocol processing
\newcommand{\Crit}{\mathsf{sec}}             % value of a security criterion
\newcommand{\Pflag}{\Pi_{\Flag}}             % flag-sector projector on Q
\newcommand{\Pleq}{\Pi_{\leqN}}              % preserved-block projector
\newcommand{\Pgt}{\Pi_{\gtN}}                % above-cutoff projector
\newcommand{\Sph}{S^{\mathrm{ph}}}           % family of physical inputs
\newcommand{\Srch}{S^{\mathrm{rch}}}         % reachable squashed states
\newcommand{\Ssq}{S^{\sq}}                   % squashed admissible set
\newcommand{\Ffss}{F^{\mathrm{FSS}}}         % flag-state-squasher target
\newcommand{\Fwp}{F^{\mathrm{WP}}}           % multi-bound WPFSS target
\newcommand{\Mwp}{M^{\mathrm{WP}}}           % renormalized above-cutoff POVM
\newcommand{\Mdec}{{M'}^{\mathrm{ph}}}       % declared physical family on B
\newcommand{\Ppvm}{P^{\Flag}}                % flag PVM
\newcommand{\Lmp}{{\Lambda_{\mathrm{mp}}}}   % measure-and-prepare channel
\newcommand{\Lwp}{{\Lambda^{\mathrm{WP}}}}   % multi-bound WPFSS channel
\newcommand{\Lfss}{{\Lambda^{\mathrm{FSS}}}} % flag-state-squasher channel
\newcommand{\lmax}{\lambda_{\max}}           % maximum eigenvalue
\newcommand{\Lnoise}{\widetilde{\Lambda}}    % flag-noise channel
\newcommand{\cg}{\varpi}                     % declared coarse-graining map
\newcommand{\MCK}{\mathsf{MC}}               % multiclick raw-pattern set
\newcommand{\redChain}{\mathfrak{C}}            % reduction chain
\newcommand{\Lnk}{\ell}                      % a single typed link
\newcommand{\dsq}{\preceq}                      % the squashing relation
\newcommand{\opn}[1]{\left\lVert #1\right\rVert}           % operator (spectral) norm
\newcommand{\mxn}[1]{\left\lVert #1\right\rVert_{\max}}    % entrywise max modulus
\newcommand{\Fld}{\mathbb{F}}                   % exact ordered substrate field
\newcommand{\Fldi}{\mathbb{F}_{i}}              % its Gaussian extension F + iF
\newcommand{\lat}[1]{\big\lfloor #1\big\rfloor_{\kappa}} % round DOWN onto the lattice

% --- protocol-compiler ---------------------------------------------------------
\newcommand{\Spec}{\mathfrak{S}}             % protocol specification (compiler input)
\newcommand{\pcChain}{\mathfrak{R}}            % reduction chain (two typed segments)
\newcommand{\Chsrc}{\pcChain_{\mathrm{src}}}   % its Alice-side source-map segment
\newcommand{\Chrec}{\pcChain_{\mathrm{rec}}}   % its receiver-side squashing segment
\newcommand{\rvirt}{\rho^{\mathrm{v}}}       % virtual (PRE-source-map) signal state
\newcommand{\Fph}{F^{\mathrm{ph}}}           % physical (pre-reduction) POVM element
\newcommand{\Vph}{\mathcal{V}^{\mathrm{ph}}} % physical outcome alphabet
\newcommand{\Tpub}{\mathcal{T}_{\mathrm{pub}}}% public-choice alphabet
\newcommand{\Uset}{\mathcal{U}}              % residual outcome alphabet
\newcommand{\idxm}{\mathsf{idx}}             % label-to-index map
\newcommand{\Asig}{A'}                       % transmitted (signal) system -- the REAL one
\newcommand{\Avirt}{A''}                     % virtual transmitted system (source descriptor of Xi_1)
\newcommand{\Abar}{\bar{A}}                  % Alice's measured register
\newcommand{\Ahat}{\hat{A}}                  % Alice's shield register
\newcommand{\pcNmsm}{\mathcal{N}_{\mathsf A}}  % Alice's measurement channel (see the RULES block)
\newcommand{\Phip}{\Phi^{+}}                 % maximally entangled vector label
\newcommand{\Smap}{\mathfrak{s}}             % the emitted sector map (fraktur: a compile-time object)

% --- receiver ------------------------------------------------------------------
\newcommand{\Modes}{\mathbb{M}}              % declared ordered mode set
\newcommand{\Inp}{\mathbb{I}}                % declared input modes
\newcommand{\InpC}{\mathbb{I}_{\Chan}}       % channel input modes
\newcommand{\InpS}{\mathbb{I}_{\Prb}}        % probe (source) input modes
\newcommand{\Env}{\mathbb{E}}                % dilation (environment) modes
\newcommand{\Port}[1]{\mathbb{P}_{#1}}       % mode set of detector port #1
\newcommand{\Emd}[1]{E_{#1}}                 % coordinate projector on modes #1
\newcommand{\Gmat}{G}                        % instrument (transfer) matrix
\newcommand{\Fock}{\mathfrak{F}}             % Fock / second-quantization functor
\newcommand{\Symn}[1]{\mathrm{Sym}^{#1}}     % symmetric tensor power
\newcommand{\normord}[1]{{:}\mspace{2mu}#1\mspace{2mu}{:}}   % normal ordering
\newcommand{\rk}{\operatorname{rank}}        % rank (physics.sty owns \rank)
\newcommand{\SCK}{\mathsf{SC}}               % single-click pattern set
\newcommand{\Recv}{\mathfrak{R}}             % receiver instrument tuple
\newcommand{\Cat}{\mathfrak{C}}              % component catalogue
\newcommand{\Grph}{\mathfrak{E}}             % declared element graph
\newcommand{\defic}{\mathbb{L}}              % loss deficiency operator
\newcommand{\Pvac}[1]{\Pi_{\vac}(#1)}        % joint no-click element on ports #1
\newcommand{\pmc}{p_{\mc}}                   % uniform high-photon mc weight
\newcommand{\Bnd}{\beta}                     % the scalar bound function beta(n)
\newcommand{\Rot}{R}                         % planar rotation matrix
\newcommand{\Wplate}{W}                      % retarder (wave-plate) matrix
\newcommand{\Bsp}{B}                         % beamsplitter core matrix

% --- statistics ----------------------------------------------------------------
\newcommand{\honest}{p^{\mathrm{hon}}}          % honest per-round distribution
\newcommand{\ecomp}{\varepsilon_{\mathrm{comp}}}% completeness budget
\newcommand{\simplex}[1]{\Delta(#1)}            % probability simplex
\newcommand{\Chatcon}{\Chatset_{\mathrm{con}}}  % constrained outcomes
\newcommand{\Ireg}[3]{\mathcal{I}_{#1}(#2,#3)}  % regularized incomplete beta
\newcommand{\Ncount}[1]{\mathsf{N}_{#1}}        % per-outcome count
\newcommand{\qlo}{\mathbf{q}^{-}}               % box lower bounds
\newcommand{\qhi}{\mathbf{q}^{+}}               % box upper bounds
 for it, so the
%  file is read from spec/modules/<id>/.
%
%  RULES
%    * A macro that spec/notation.md catalogues is a verbatim copy of it
%      (part A). Changing one is a notation change: bump spec/notation.md.
%    * Part B holds the module-local macros. A macro defined by two modules
%      with the same expansion appears once, under the first module (in the
%      order below) that defines it. Two macros with the same name and
%      DIFFERENT expansions are both kept under a module prefix, and the
%      claims write the prefixed names:
%        \Chain  -> \redChain (reductions, \mathfrak{C}: the reduction chain)
%                   \pcChain  (protocol-compiler, \mathfrak{R}: the two-segment
%                             chain; \Chsrc and \Chrec are built on it)
%        \Nmsm   -> \redNmsm  (reductions, \mathcal{N}: a POVM measurement
%                             channel)
%                   \pcNmsm   (\mathcal{N}_{\mathsf A}: Alice's measurement
%                             channel, written by
%                             transmitter/thm:certificate-sufficiency,
%                             transmitter/rem:xreg and
%                             reductions/lem:source-transfer)
%    * Deliberately omitted: a \provenance{tag}{details} macro. Provenance
%      is a header field of every claim file (spec/claim-format.md section
%      3); a claim body that calls such a macro fails to compile, which is
%      the intended signal.
%
%  REQUIRED PACKAGES, loaded by the assembler's preamble before this file:
%    amsmath, amssymb, amsfonts, mathtools, dsfont, braket, physics, xcolor
%    (braket before physics, so that physics wins for \ket, \bra, \tr, ...).
%
%  This file contains only definitions: no \usepackage, no document body.
% =============================================================================


% =============================================================================
%  PART A -- notation macros, the ones spec/notation.md catalogues, each a
%  verbatim copy of its catalogue row; grouped under numbered headings.
% =============================================================================

% -----------------------------------------------------------------------------
% 1. State spaces and operator sets (a "system label" -> spaces of operators)
% -----------------------------------------------------------------------------
\newcommand{\Hs}[1]{\mathcal{H}(#1)}        % Hilbert space of a system label
\newcommand{\Lin}[1]{\mathcal{L}(#1)}       % bounded linear operators on H(#1)
\newcommand{\Linn}[2]{\mathcal{L}(#1,#2)}   % linear maps H(#1) -> H(#2)
\newcommand{\Pos}[1]{\mathrm{Pos}(#1)}      % positive semidefinite operators on H(#1)
\newcommand{\Density}{\mathrm{D}}           % density operators: \Density(A) = D(A)
\newcommand{\Seq}[1]{S_{=}(#1)}             % normalized states (trace = 1)
\newcommand{\Sleq}[1]{S_{\leq}(#1)}         % subnormalized states (0 <= trace <= 1)
\newcommand{\id}{\mathds{1}}                % identity operator
\newcommand{\idon}[1]{\mathds{1}_{#1}}      % identity on a named system, e.g. \idon{C}

% DEPRECATED: applies \mathcal to both arguments itself, so it takes bare
% letters, not system labels. Prefer \Linn; do not use in new contracts.
\providecommand{\Lop}[2]{\mathcal{L}(\mathcal{#1}, \mathcal{#2})}


% -----------------------------------------------------------------------------
% 2. Composite systems
% -----------------------------------------------------------------------------
\newcommand{\dsum}{\oplus}                  % direct sum of system labels
% Tensor product: use \otimes directly.
% n-fold direct sum of an operator space: \Lin{A}^{\oplus n}


% -----------------------------------------------------------------------------
% 3. Events, conditioning and classical registers (for STATES)
%    Convention:  \wedge  -> subnormalized "and-event" component, e.g. \rho_{A\wedge x}
%                 |        -> normalized conditional state,        e.g. \rho_{A|x}
% (Write these directly in subscripts; no extra macros needed.)
% -----------------------------------------------------------------------------


% -----------------------------------------------------------------------------
% 4. Photon-number subspaces / direct-sum blocks
%    Convention:  operators carry the restriction "on top", in brackets;
%                 states carry it as a \wedge subscript (see Sec. 3).
% -----------------------------------------------------------------------------
\newcommand{\Pnum}[1]{P_{#1}}               % photon-number projector (system from context)
\newcommand{\Pnumon}[2]{P_{#1}^{#2}}        % with explicit system superscript, e.g. P_n^C
\newcommand{\blk}[2]{#1^{(#2)}}             % operator block, e.g. \blk{M}{\leqN} = M^{(<=N)}
\newcommand{\leqN}{\leq N}                  % use as \blk{M}{\leqN}
\newcommand{\gtN}{> N}                      % use as \blk{M}{\gtN}
\newcommand{\eqn}{= n}                      % use as \blk{M}{\eqn}


% -----------------------------------------------------------------------------
% 5. System labels used in the protocol
% -----------------------------------------------------------------------------
\newcommand{\Chan}{C}                       % quantum-channel (Eve) input to Bob
\newcommand{\Prb}{P}                        % probe input
\newcommand{\Flag}{{\mathrm{Flg}}}          % flag system (EXISTS ONLY AFTER squashing)
\newcommand{\Loss}{{\mathrm{Lss}}}          % loss system
\newcommand{\Vc}{{\mathrm{Vcm}}}            % vacuum system
\newcommand{\Nout}{N_{\mathrm{out}}}
\newcommand{\Nin}{N_{\mathrm{in}}}


% -----------------------------------------------------------------------------
% 6. Alphabets
% -----------------------------------------------------------------------------
\newcommand{\Yset}{\mathcal{Y}}             % Bob's (coarse-grained) detection alphabet
\newcommand{\Yraw}{\mathcal{Y}_{\mathrm{raw}}}% raw click patterns {0,1}^4
\newcommand{\Yscset}{\mathcal{Y}_{\singlec}}% single-click outcomes
\newcommand{\Ymcset}{\mathcal{Y}_{\mc}}     % multi-click outcomes
\newcommand{\Tset}{\mathcal{T}_{\mathrm{set}}}% private setting set (the probe settings in the optical instance)
\newcommand{\Xset}{\mathcal{X}}             % Alice's alphabet
\newcommand{\Dset}{\{H,V,+,-\}}             % polarization / detector labels
\newcommand{\pass}{\top}                    % pass sifting


% -----------------------------------------------------------------------------
% 7. Mode / outcome abbreviations
% -----------------------------------------------------------------------------
\newcommand{\on}{{\mathrm{on}}}
\newcommand{\off}{{\mathrm{off}}}
\newcommand{\dc}{{\mathrm{dc}}}             % dark count
\newcommand{\sq}{\mathrm{sq}}               % squashed
\newcommand{\mc}{\mathrm{mc}}               % multi-click
\newcommand{\nc}{\mathrm{nc}}               % no-click
\newcommand{\singlec}{\mathrm{sc}}          % single-click
\newcommand{\vac}{\mathrm{vac}}             % vacuum
\newcommand{\kout}{{k_\mathrm{out}}}
\newcommand{\kin}{{k_\mathrm{in}}}


% -----------------------------------------------------------------------------
% 8. Maps, channels and miscellaneous operators
% -----------------------------------------------------------------------------
\newcommand{\Prob}{\mathrm{Prob}}           % probability (use \Prob[\cdot] or \Prob(\cdot))
\newcommand{\lmin}{\lambda_{\min}}          % minimum eigenvalue
\newcommand{\Wflag}{W_{\Flag}}              % flag-space weight
\newcommand{\transpose}{\mathsf{T}}
\newcommand{\innerp}[2]{\left\langle #1 \vert #2 \right\rangle}
\renewcommand{\vec}[1]{\boldsymbol{#1}}     % NB: overrides the LaTeX arrow accent


% -----------------------------------------------------------------------------
% 9. Key-rate numerics
% -----------------------------------------------------------------------------
\newcommand{\cG}{\mathcal{G}}               % protocol map
\newcommand{\cZ}{\mathcal{Z}}               % pinching map
\newcommand{\Amap}{\mathcal{A}_{\mathrm{ann}}} % announcement + measurement map
\newcommand{\Amapnb}{\mathcal{A}_{\mathrm{ann}|\mathrm{nb}}}
\newcommand{\simpl}{}                       % empty suffix for simplified block maps
\newcommand{\acSet}{\mathcal{S}_{\mathrm{acc}}} % acceptance set
\newcommand{\testE}{\mathcal{E}_{\mathrm{test}}} % test-statistics map
\newcommand{\ecLeak}{\mathrm{leak}^{\mathrm{EC}}_{\mathrm{obs}}}
\newcommand{\pgen}{p_{\mathrm{gen}}}
\newcommand{\askeyr}{R_{\infty}}             % asymptotic key rate
\newcommand{\ppass}{p_{\mathrm{pass}}}       % sifting pass probability


% -----------------------------------------------------------------------------
% 10. Finite-size analysis / Renyi entropies
% -----------------------------------------------------------------------------

% --- Round types: test vs. key generation (register T, alphabet \Fset) --------
\newcommand{\tst}{\mathrm{test}}            % test round
\newcommand{\gen}{\mathrm{gen}}             % key-generation round
\newcommand{\Fset}{\mathcal{T}}             % round-type alphabet {test, gen} (not \Tset)
\newcommand{\Chat}{\widehat{C}}             % acceptance-test register
\newcommand{\Chatset}{\widehat{\mathcal{C}}}% its alphabet

% --- Basis / key extraction functions ----------------------------------------
\newcommand{\bA}{b_{A}}                     % Alice's basis-extraction function
\newcommand{\bB}{b_{B}}                     % Bob's basis-extraction function
\newcommand{\kB}{k_{B}}                     % Bob's key-extraction function

% --- Maps --------------------------------------------------------------------
\newcommand{\Mprot}{\mathcal{M}}            % single-round protocol map
\newcommand{\Phimap}{\boldsymbol{\Phi}}     % test-statistics map
\newcommand{\Pdeph}{\mathcal{P}}            % dephasing (pinching) map on \Blreg

% --- Isometries --------------------------------------------------------------
\newcommand{\Vmsm}{V_{\mathrm{msm}}}        % measurement isometry
\newcommand{\Vann}{V_{\mathrm{ann}}}        % announcement isometry
\newcommand{\VI}{V_{I}}                     % Eve's copy of the announcements
\newcommand{\VChat}{V_{\Chat}}              % acceptance-data isometry
\newcommand{\VR}{V_{R}}                     % sifting + key map -> raw key R
\newcommand{\VZ}{V_{Z}}                     % raw key -> key register Z
\newcommand{\Wprot}{W}                      % full protocol isometry

% --- Renyi entropies ---------------------------------------------------------
% SECURITY-CRITICAL: see spec/notation.md §4. Never change an arrow to
% make an inequality work out; never convert silently between the sandwiched
% (tilde) and Petz (overline) families.
\newcommand{\Hsandup}{\widetilde{H}^{\uparrow}_{\alpha}}   % sandwiched, up
\newcommand{\Hsanddown}{\widetilde{H}^{\downarrow}_{\alpha}}% sandwiched, down
\newcommand{\Hpetzup}{\overline{H}^{\uparrow}_{\alpha}}    % Petz, up
\newcommand{\Hpetzdown}{\overline{H}^{\downarrow}_{\alpha}}% Petz, down
\newcommand{\hup}{h^{\uparrow}_{\alpha}}                   % single-round rate
\newcommand{\hqkd}{h^{\uparrow,\mathrm{QKD}}_{\alpha}}     % single-round QKD rate
\newcommand{\hlow}{\underline{h}_{\alpha}}                 % lower bound on \hqkd
% --- spec/notation.md 3.5 and 4.5 -------------------------------------------
% \hlb  is the first-order linearization lower bound on \hlow: the headline
%       key rate (the `first-order` tier of foundations/def:tiers). NOT \hlow
%       itself.
% \tnum is the instance's one declared numerical tolerance. Its
%       near-collision with the declared marginal \tau_A is disclosed in
%       spec/notation.md 3.5; bare \tau is reserved and denotes neither.
\newcommand{\hlb}{\underline{h}^{\mathrm{FO}}_{\alpha}}    % first-order lower bound on \hlow
\newcommand{\tnum}{\tau_{\mathrm{num}}}                    % declared numerical tolerance


% =============================================================================
%  PART B -- module-local macros, grouped by module. Deprecated aliases,
%  defined so that a text written with them compiles, used in no new contract: \Lop,
%  \testE, \Pdeph, \pass, \ecLeak (spec/notation.md section 3.9).
% =============================================================================

% --- meat-core -----------------------------------------------------------------
\newcommand{\SecSet}{\mathbb{S}}            % sector index set of Q
\newcommand{\Qsys}{Q}                       % declared sector-carrying measured system
\newcommand{\Vset}{\mathcal{V}}             % declared joint outcome alphabet of the Q-POVM
\newcommand{\CAset}{\mathcal{C}_{A}}        % Alice announcement alphabet
\newcommand{\CBset}{\mathcal{C}_{B}}        % Bob announcement alphabet
\newcommand{\Party}{\mathfrak{P}}           % declared party partition
\newcommand{\Oset}[1]{\mathcal{O}_{#1}}     % readout alphabet of party #1
\newcommand{\rd}[1]{\pi_{#1}}               % readout map of party #1
\newcommand{\Sadm}{\mathcal{S}_{A}}         % admissible Alice subspace
\newcommand{\nomarg}{\mathsf{none}}         % "no trusted marginal declared"
\newcommand{\Zset}{\mathcal{Z}_{\perp}}     % raw-key alphabet incl. \perp
\newcommand{\Zval}{\mathcal{Z}}             % key alphabet without \perp
\newcommand{\Etil}{\widetilde{E}}           % single-round purifying register
\newcommand{\Ehat}{\widehat{E}}             % Eve's n-round register
\newcommand{\OmAT}{\Omega_{\mathrm{AT}}}    % acceptance event
\newcommand{\epa}{\varepsilon_{\mathrm{PA}}}
\newcommand{\eev}{\varepsilon_{\mathrm{EV}}}
\newcommand{\lEC}{\lambda_{\mathrm{EC}}}
\newcommand{\Decl}{\mathfrak{D}}            % protocol declaration
\newcommand{\CList}{\mathfrak{L}}           % declared constraint list
\newcommand{\PinS}{\mathcal{P}_{\SecSet}}   % sector pinching on Q
\newcommand{\Feas}{\mathcal{W}}             % feasible set of block states
\newcommand{\Chkey}{\Chatset_{\mathrm{key}}}% key-live acceptance values
\newcommand{\Kann}{\mathcal{K}}             % announcement-class alphabet
\newcommand{\hold}{\mathsf{h}}              % declared key holder (A or B)
\newcommand{\sfD}{D}                        % classical KL divergence (base 2)
\newcommand{\Dtwo}{D}                       % quantum relative entropy (base 2)
\newcommand{\Dsand}{\widetilde{D}_{\alpha}} % sandwiched Renyi divergence
\newcommand{\idmap}{\mathrm{id}}            % identity channel
\newcommand{\Aev}{\mathcal{A}}              % Eve's per-round attack channel
\newcommand{\AttSet}{\mathfrak{A}}          % licensed attack-channel set
\newcommand{\Chacc}{\Chatset_{\mathrm{acc}}}% acceptance support of \acSet
\newcommand{\epsstar}{\varepsilon^{\star}}      % distance to an exact POVM
\newcommand{\dfH}{\Delta_{\mathrm{herm}}}       % Hermiticity defect
\newcommand{\dfP}{\Delta_{\mathrm{pos}}}        % negative part
\newcommand{\dfC}{\Delta_{\mathrm{cpl}}}        % completeness defect
\newcommand{\lmoq}{g}                       % acceptance-side gradient block
\newcommand{\lmoJ}{G}                       % feasible-side gradient block
\newcommand{\bacc}{\beta_{\mathrm{acc}}}    % certified acceptance-side minimum
\newcommand{\Itset}{\mathfrak{I}}           % finite index set of iterates (C22)

% --- reductions ----------------------------------------------------------------
\newcommand{\Bsys}{B}                        % physical (pre-squashing) system
\newcommand{\Tcl}[1]{\mathcal{T}(#1)}        % trace-class operators on H(#1)
\newcommand{\Ins}{\mathcal{I}}               % generalized-instrument branch map
\newcommand{\redNmsm}{\mathcal{N}}              % POVM measurement channel
\newcommand{\Eprot}{\mathcal{E}}             % n-round protocol channel
\newcommand{\Grnd}{\mathsf{G}}               % one round map of a protocol
\newcommand{\Kfin}{\mathsf{K}}               % end-of-protocol processing
\newcommand{\Crit}{\mathsf{sec}}             % value of a security criterion
\newcommand{\Pflag}{\Pi_{\Flag}}             % flag-sector projector on Q
\newcommand{\Pleq}{\Pi_{\leqN}}              % preserved-block projector
\newcommand{\Pgt}{\Pi_{\gtN}}                % above-cutoff projector
\newcommand{\Sph}{S^{\mathrm{ph}}}           % family of physical inputs
\newcommand{\Srch}{S^{\mathrm{rch}}}         % reachable squashed states
\newcommand{\Ssq}{S^{\sq}}                   % squashed admissible set
\newcommand{\Ffss}{F^{\mathrm{FSS}}}         % flag-state-squasher target
\newcommand{\Fwp}{F^{\mathrm{WP}}}           % multi-bound WPFSS target
\newcommand{\Mwp}{M^{\mathrm{WP}}}           % renormalized above-cutoff POVM
\newcommand{\Mdec}{{M'}^{\mathrm{ph}}}       % declared physical family on B
\newcommand{\Ppvm}{P^{\Flag}}                % flag PVM
\newcommand{\Lmp}{{\Lambda_{\mathrm{mp}}}}   % measure-and-prepare channel
\newcommand{\Lwp}{{\Lambda^{\mathrm{WP}}}}   % multi-bound WPFSS channel
\newcommand{\Lfss}{{\Lambda^{\mathrm{FSS}}}} % flag-state-squasher channel
\newcommand{\lmax}{\lambda_{\max}}           % maximum eigenvalue
\newcommand{\Lnoise}{\widetilde{\Lambda}}    % flag-noise channel
\newcommand{\cg}{\varpi}                     % declared coarse-graining map
\newcommand{\MCK}{\mathsf{MC}}               % multiclick raw-pattern set
\newcommand{\redChain}{\mathfrak{C}}            % reduction chain
\newcommand{\Lnk}{\ell}                      % a single typed link
\newcommand{\dsq}{\preceq}                      % the squashing relation
\newcommand{\opn}[1]{\left\lVert #1\right\rVert}           % operator (spectral) norm
\newcommand{\mxn}[1]{\left\lVert #1\right\rVert_{\max}}    % entrywise max modulus
\newcommand{\Fld}{\mathbb{F}}                   % exact ordered substrate field
\newcommand{\Fldi}{\mathbb{F}_{i}}              % its Gaussian extension F + iF
\newcommand{\lat}[1]{\big\lfloor #1\big\rfloor_{\kappa}} % round DOWN onto the lattice

% --- protocol-compiler ---------------------------------------------------------
\newcommand{\Spec}{\mathfrak{S}}             % protocol specification (compiler input)
\newcommand{\pcChain}{\mathfrak{R}}            % reduction chain (two typed segments)
\newcommand{\Chsrc}{\pcChain_{\mathrm{src}}}   % its Alice-side source-map segment
\newcommand{\Chrec}{\pcChain_{\mathrm{rec}}}   % its receiver-side squashing segment
\newcommand{\rvirt}{\rho^{\mathrm{v}}}       % virtual (PRE-source-map) signal state
\newcommand{\Fph}{F^{\mathrm{ph}}}           % physical (pre-reduction) POVM element
\newcommand{\Vph}{\mathcal{V}^{\mathrm{ph}}} % physical outcome alphabet
\newcommand{\Tpub}{\mathcal{T}_{\mathrm{pub}}}% public-choice alphabet
\newcommand{\Uset}{\mathcal{U}}              % residual outcome alphabet
\newcommand{\idxm}{\mathsf{idx}}             % label-to-index map
\newcommand{\Asig}{A'}                       % transmitted (signal) system -- the REAL one
\newcommand{\Avirt}{A''}                     % virtual transmitted system (source descriptor of Xi_1)
\newcommand{\Abar}{\bar{A}}                  % Alice's measured register
\newcommand{\Ahat}{\hat{A}}                  % Alice's shield register
\newcommand{\pcNmsm}{\mathcal{N}_{\mathsf A}}  % Alice's measurement channel (see the RULES block)
\newcommand{\Phip}{\Phi^{+}}                 % maximally entangled vector label
\newcommand{\Smap}{\mathfrak{s}}             % the emitted sector map (fraktur: a compile-time object)

% --- receiver ------------------------------------------------------------------
\newcommand{\Modes}{\mathbb{M}}              % declared ordered mode set
\newcommand{\Inp}{\mathbb{I}}                % declared input modes
\newcommand{\InpC}{\mathbb{I}_{\Chan}}       % channel input modes
\newcommand{\InpS}{\mathbb{I}_{\Prb}}        % probe (source) input modes
\newcommand{\Env}{\mathbb{E}}                % dilation (environment) modes
\newcommand{\Port}[1]{\mathbb{P}_{#1}}       % mode set of detector port #1
\newcommand{\Emd}[1]{E_{#1}}                 % coordinate projector on modes #1
\newcommand{\Gmat}{G}                        % instrument (transfer) matrix
\newcommand{\Fock}{\mathfrak{F}}             % Fock / second-quantization functor
\newcommand{\Symn}[1]{\mathrm{Sym}^{#1}}     % symmetric tensor power
\newcommand{\normord}[1]{{:}\mspace{2mu}#1\mspace{2mu}{:}}   % normal ordering
\newcommand{\rk}{\operatorname{rank}}        % rank (physics.sty owns \rank)
\newcommand{\SCK}{\mathsf{SC}}               % single-click pattern set
\newcommand{\Recv}{\mathfrak{R}}             % receiver instrument tuple
\newcommand{\Cat}{\mathfrak{C}}              % component catalogue
\newcommand{\Grph}{\mathfrak{E}}             % declared element graph
\newcommand{\defic}{\mathbb{L}}              % loss deficiency operator
\newcommand{\Pvac}[1]{\Pi_{\vac}(#1)}        % joint no-click element on ports #1
\newcommand{\pmc}{p_{\mc}}                   % uniform high-photon mc weight
\newcommand{\Bnd}{\beta}                     % the scalar bound function beta(n)
\newcommand{\Rot}{R}                         % planar rotation matrix
\newcommand{\Wplate}{W}                      % retarder (wave-plate) matrix
\newcommand{\Bsp}{B}                         % beamsplitter core matrix

% --- statistics ----------------------------------------------------------------
\newcommand{\honest}{p^{\mathrm{hon}}}          % honest per-round distribution
\newcommand{\ecomp}{\varepsilon_{\mathrm{comp}}}% completeness budget
\newcommand{\simplex}[1]{\Delta(#1)}            % probability simplex
\newcommand{\Chatcon}{\Chatset_{\mathrm{con}}}  % constrained outcomes
\newcommand{\Ireg}[3]{\mathcal{I}_{#1}(#2,#3)}  % regularized incomplete beta
\newcommand{\Ncount}[1]{\mathsf{N}_{#1}}        % per-outcome count
\newcommand{\qlo}{\mathbf{q}^{-}}               % box lower bounds
\newcommand{\qhi}{\mathbf{q}^{+}}               % box upper bounds
 for it, so the
%  file is read from spec/modules/<id>/.
%
%  RULES
%    * A macro that spec/notation.md catalogues is a verbatim copy of it
%      (part A). Changing one is a notation change: bump spec/notation.md.
%    * Part B holds the module-local macros. A macro defined by two modules
%      with the same expansion appears once, under the first module (in the
%      order below) that defines it. Two macros with the same name and
%      DIFFERENT expansions are both kept under a module prefix, and the
%      claims write the prefixed names:
%        \Chain  -> \redChain (reductions, \mathfrak{C}: the reduction chain)
%                   \pcChain  (protocol-compiler, \mathfrak{R}: the two-segment
%                             chain; \Chsrc and \Chrec are built on it)
%        \Nmsm   -> \redNmsm  (reductions, \mathcal{N}: a POVM measurement
%                             channel)
%                   \pcNmsm   (\mathcal{N}_{\mathsf A}: Alice's measurement
%                             channel, written by
%                             transmitter/thm:certificate-sufficiency,
%                             transmitter/rem:xreg and
%                             reductions/lem:source-transfer)
%    * Deliberately omitted: a \provenance{tag}{details} macro. Provenance
%      is a header field of every claim file (spec/claim-format.md section
%      3); a claim body that calls such a macro fails to compile, which is
%      the intended signal.
%
%  REQUIRED PACKAGES, loaded by the assembler's preamble before this file:
%    amsmath, amssymb, amsfonts, mathtools, dsfont, braket, physics, xcolor
%    (braket before physics, so that physics wins for \ket, \bra, \tr, ...).
%
%  This file contains only definitions: no \usepackage, no document body.
% =============================================================================


% =============================================================================
%  PART A -- notation macros, the ones spec/notation.md catalogues, each a
%  verbatim copy of its catalogue row; grouped under numbered headings.
% =============================================================================

% -----------------------------------------------------------------------------
% 1. State spaces and operator sets (a "system label" -> spaces of operators)
% -----------------------------------------------------------------------------
\newcommand{\Hs}[1]{\mathcal{H}(#1)}        % Hilbert space of a system label
\newcommand{\Lin}[1]{\mathcal{L}(#1)}       % bounded linear operators on H(#1)
\newcommand{\Linn}[2]{\mathcal{L}(#1,#2)}   % linear maps H(#1) -> H(#2)
\newcommand{\Pos}[1]{\mathrm{Pos}(#1)}      % positive semidefinite operators on H(#1)
\newcommand{\Density}{\mathrm{D}}           % density operators: \Density(A) = D(A)
\newcommand{\Seq}[1]{S_{=}(#1)}             % normalized states (trace = 1)
\newcommand{\Sleq}[1]{S_{\leq}(#1)}         % subnormalized states (0 <= trace <= 1)
\newcommand{\id}{\mathds{1}}                % identity operator
\newcommand{\idon}[1]{\mathds{1}_{#1}}      % identity on a named system, e.g. \idon{C}

% DEPRECATED: applies \mathcal to both arguments itself, so it takes bare
% letters, not system labels. Prefer \Linn; do not use in new contracts.
\providecommand{\Lop}[2]{\mathcal{L}(\mathcal{#1}, \mathcal{#2})}


% -----------------------------------------------------------------------------
% 2. Composite systems
% -----------------------------------------------------------------------------
\newcommand{\dsum}{\oplus}                  % direct sum of system labels
% Tensor product: use \otimes directly.
% n-fold direct sum of an operator space: \Lin{A}^{\oplus n}


% -----------------------------------------------------------------------------
% 3. Events, conditioning and classical registers (for STATES)
%    Convention:  \wedge  -> subnormalized "and-event" component, e.g. \rho_{A\wedge x}
%                 |        -> normalized conditional state,        e.g. \rho_{A|x}
% (Write these directly in subscripts; no extra macros needed.)
% -----------------------------------------------------------------------------


% -----------------------------------------------------------------------------
% 4. Photon-number subspaces / direct-sum blocks
%    Convention:  operators carry the restriction "on top", in brackets;
%                 states carry it as a \wedge subscript (see Sec. 3).
% -----------------------------------------------------------------------------
\newcommand{\Pnum}[1]{P_{#1}}               % photon-number projector (system from context)
\newcommand{\Pnumon}[2]{P_{#1}^{#2}}        % with explicit system superscript, e.g. P_n^C
\newcommand{\blk}[2]{#1^{(#2)}}             % operator block, e.g. \blk{M}{\leqN} = M^{(<=N)}
\newcommand{\leqN}{\leq N}                  % use as \blk{M}{\leqN}
\newcommand{\gtN}{> N}                      % use as \blk{M}{\gtN}
\newcommand{\eqn}{= n}                      % use as \blk{M}{\eqn}


% -----------------------------------------------------------------------------
% 5. System labels used in the protocol
% -----------------------------------------------------------------------------
\newcommand{\Chan}{C}                       % quantum-channel (Eve) input to Bob
\newcommand{\Prb}{P}                        % probe input
\newcommand{\Flag}{{\mathrm{Flg}}}          % flag system (EXISTS ONLY AFTER squashing)
\newcommand{\Loss}{{\mathrm{Lss}}}          % loss system
\newcommand{\Vc}{{\mathrm{Vcm}}}            % vacuum system
\newcommand{\Nout}{N_{\mathrm{out}}}
\newcommand{\Nin}{N_{\mathrm{in}}}


% -----------------------------------------------------------------------------
% 6. Alphabets
% -----------------------------------------------------------------------------
\newcommand{\Yset}{\mathcal{Y}}             % Bob's (coarse-grained) detection alphabet
\newcommand{\Yraw}{\mathcal{Y}_{\mathrm{raw}}}% raw click patterns {0,1}^4
\newcommand{\Yscset}{\mathcal{Y}_{\singlec}}% single-click outcomes
\newcommand{\Ymcset}{\mathcal{Y}_{\mc}}     % multi-click outcomes
\newcommand{\Tset}{\mathcal{T}_{\mathrm{set}}}% private setting set (the probe settings in the optical instance)
\newcommand{\Xset}{\mathcal{X}}             % Alice's alphabet
\newcommand{\Dset}{\{H,V,+,-\}}             % polarization / detector labels
\newcommand{\pass}{\top}                    % pass sifting


% -----------------------------------------------------------------------------
% 7. Mode / outcome abbreviations
% -----------------------------------------------------------------------------
\newcommand{\on}{{\mathrm{on}}}
\newcommand{\off}{{\mathrm{off}}}
\newcommand{\dc}{{\mathrm{dc}}}             % dark count
\newcommand{\sq}{\mathrm{sq}}               % squashed
\newcommand{\mc}{\mathrm{mc}}               % multi-click
\newcommand{\nc}{\mathrm{nc}}               % no-click
\newcommand{\singlec}{\mathrm{sc}}          % single-click
\newcommand{\vac}{\mathrm{vac}}             % vacuum
\newcommand{\kout}{{k_\mathrm{out}}}
\newcommand{\kin}{{k_\mathrm{in}}}


% -----------------------------------------------------------------------------
% 8. Maps, channels and miscellaneous operators
% -----------------------------------------------------------------------------
\newcommand{\Prob}{\mathrm{Prob}}           % probability (use \Prob[\cdot] or \Prob(\cdot))
\newcommand{\lmin}{\lambda_{\min}}          % minimum eigenvalue
\newcommand{\Wflag}{W_{\Flag}}              % flag-space weight
\newcommand{\transpose}{\mathsf{T}}
\newcommand{\innerp}[2]{\left\langle #1 \vert #2 \right\rangle}
\renewcommand{\vec}[1]{\boldsymbol{#1}}     % NB: overrides the LaTeX arrow accent


% -----------------------------------------------------------------------------
% 9. Key-rate numerics
% -----------------------------------------------------------------------------
\newcommand{\cG}{\mathcal{G}}               % protocol map
\newcommand{\cZ}{\mathcal{Z}}               % pinching map
\newcommand{\Amap}{\mathcal{A}_{\mathrm{ann}}} % announcement + measurement map
\newcommand{\Amapnb}{\mathcal{A}_{\mathrm{ann}|\mathrm{nb}}}
\newcommand{\simpl}{}                       % empty suffix for simplified block maps
\newcommand{\acSet}{\mathcal{S}_{\mathrm{acc}}} % acceptance set
\newcommand{\testE}{\mathcal{E}_{\mathrm{test}}} % test-statistics map
\newcommand{\ecLeak}{\mathrm{leak}^{\mathrm{EC}}_{\mathrm{obs}}}
\newcommand{\pgen}{p_{\mathrm{gen}}}
\newcommand{\askeyr}{R_{\infty}}             % asymptotic key rate
\newcommand{\ppass}{p_{\mathrm{pass}}}       % sifting pass probability


% -----------------------------------------------------------------------------
% 10. Finite-size analysis / Renyi entropies
% -----------------------------------------------------------------------------

% --- Round types: test vs. key generation (register T, alphabet \Fset) --------
\newcommand{\tst}{\mathrm{test}}            % test round
\newcommand{\gen}{\mathrm{gen}}             % key-generation round
\newcommand{\Fset}{\mathcal{T}}             % round-type alphabet {test, gen} (not \Tset)
\newcommand{\Chat}{\widehat{C}}             % acceptance-test register
\newcommand{\Chatset}{\widehat{\mathcal{C}}}% its alphabet

% --- Basis / key extraction functions ----------------------------------------
\newcommand{\bA}{b_{A}}                     % Alice's basis-extraction function
\newcommand{\bB}{b_{B}}                     % Bob's basis-extraction function
\newcommand{\kB}{k_{B}}                     % Bob's key-extraction function

% --- Maps --------------------------------------------------------------------
\newcommand{\Mprot}{\mathcal{M}}            % single-round protocol map
\newcommand{\Phimap}{\boldsymbol{\Phi}}     % test-statistics map
\newcommand{\Pdeph}{\mathcal{P}}            % dephasing (pinching) map on \Blreg

% --- Isometries --------------------------------------------------------------
\newcommand{\Vmsm}{V_{\mathrm{msm}}}        % measurement isometry
\newcommand{\Vann}{V_{\mathrm{ann}}}        % announcement isometry
\newcommand{\VI}{V_{I}}                     % Eve's copy of the announcements
\newcommand{\VChat}{V_{\Chat}}              % acceptance-data isometry
\newcommand{\VR}{V_{R}}                     % sifting + key map -> raw key R
\newcommand{\VZ}{V_{Z}}                     % raw key -> key register Z
\newcommand{\Wprot}{W}                      % full protocol isometry

% --- Renyi entropies ---------------------------------------------------------
% SECURITY-CRITICAL: see spec/notation.md §4. Never change an arrow to
% make an inequality work out; never convert silently between the sandwiched
% (tilde) and Petz (overline) families.
\newcommand{\Hsandup}{\widetilde{H}^{\uparrow}_{\alpha}}   % sandwiched, up
\newcommand{\Hsanddown}{\widetilde{H}^{\downarrow}_{\alpha}}% sandwiched, down
\newcommand{\Hpetzup}{\overline{H}^{\uparrow}_{\alpha}}    % Petz, up
\newcommand{\Hpetzdown}{\overline{H}^{\downarrow}_{\alpha}}% Petz, down
\newcommand{\hup}{h^{\uparrow}_{\alpha}}                   % single-round rate
\newcommand{\hqkd}{h^{\uparrow,\mathrm{QKD}}_{\alpha}}     % single-round QKD rate
\newcommand{\hlow}{\underline{h}_{\alpha}}                 % lower bound on \hqkd
% --- spec/notation.md 3.5 and 4.5 -------------------------------------------
% \hlb  is the first-order linearization lower bound on \hlow: the headline
%       key rate (the `first-order` tier of foundations/def:tiers). NOT \hlow
%       itself.
% \tnum is the instance's one declared numerical tolerance. Its
%       near-collision with the declared marginal \tau_A is disclosed in
%       spec/notation.md 3.5; bare \tau is reserved and denotes neither.
\newcommand{\hlb}{\underline{h}^{\mathrm{FO}}_{\alpha}}    % first-order lower bound on \hlow
\newcommand{\tnum}{\tau_{\mathrm{num}}}                    % declared numerical tolerance


% =============================================================================
%  PART B -- module-local macros, grouped by module. Deprecated aliases,
%  defined so that a text written with them compiles, used in no new contract: \Lop,
%  \testE, \Pdeph, \pass, \ecLeak (spec/notation.md section 3.9).
% =============================================================================

% --- meat-core -----------------------------------------------------------------
\newcommand{\SecSet}{\mathbb{S}}            % sector index set of Q
\newcommand{\Qsys}{Q}                       % declared sector-carrying measured system
\newcommand{\Vset}{\mathcal{V}}             % declared joint outcome alphabet of the Q-POVM
\newcommand{\CAset}{\mathcal{C}_{A}}        % Alice announcement alphabet
\newcommand{\CBset}{\mathcal{C}_{B}}        % Bob announcement alphabet
\newcommand{\Party}{\mathfrak{P}}           % declared party partition
\newcommand{\Oset}[1]{\mathcal{O}_{#1}}     % readout alphabet of party #1
\newcommand{\rd}[1]{\pi_{#1}}               % readout map of party #1
\newcommand{\Sadm}{\mathcal{S}_{A}}         % admissible Alice subspace
\newcommand{\nomarg}{\mathsf{none}}         % "no trusted marginal declared"
\newcommand{\Zset}{\mathcal{Z}_{\perp}}     % raw-key alphabet incl. \perp
\newcommand{\Zval}{\mathcal{Z}}             % key alphabet without \perp
\newcommand{\Etil}{\widetilde{E}}           % single-round purifying register
\newcommand{\Ehat}{\widehat{E}}             % Eve's n-round register
\newcommand{\OmAT}{\Omega_{\mathrm{AT}}}    % acceptance event
\newcommand{\epa}{\varepsilon_{\mathrm{PA}}}
\newcommand{\eev}{\varepsilon_{\mathrm{EV}}}
\newcommand{\lEC}{\lambda_{\mathrm{EC}}}
\newcommand{\Decl}{\mathfrak{D}}            % protocol declaration
\newcommand{\CList}{\mathfrak{L}}           % declared constraint list
\newcommand{\PinS}{\mathcal{P}_{\SecSet}}   % sector pinching on Q
\newcommand{\Feas}{\mathcal{W}}             % feasible set of block states
\newcommand{\Chkey}{\Chatset_{\mathrm{key}}}% key-live acceptance values
\newcommand{\Kann}{\mathcal{K}}             % announcement-class alphabet
\newcommand{\hold}{\mathsf{h}}              % declared key holder (A or B)
\newcommand{\sfD}{D}                        % classical KL divergence (base 2)
\newcommand{\Dtwo}{D}                       % quantum relative entropy (base 2)
\newcommand{\Dsand}{\widetilde{D}_{\alpha}} % sandwiched Renyi divergence
\newcommand{\idmap}{\mathrm{id}}            % identity channel
\newcommand{\Aev}{\mathcal{A}}              % Eve's per-round attack channel
\newcommand{\AttSet}{\mathfrak{A}}          % licensed attack-channel set
\newcommand{\Chacc}{\Chatset_{\mathrm{acc}}}% acceptance support of \acSet
\newcommand{\epsstar}{\varepsilon^{\star}}      % distance to an exact POVM
\newcommand{\dfH}{\Delta_{\mathrm{herm}}}       % Hermiticity defect
\newcommand{\dfP}{\Delta_{\mathrm{pos}}}        % negative part
\newcommand{\dfC}{\Delta_{\mathrm{cpl}}}        % completeness defect
\newcommand{\lmoq}{g}                       % acceptance-side gradient block
\newcommand{\lmoJ}{G}                       % feasible-side gradient block
\newcommand{\bacc}{\beta_{\mathrm{acc}}}    % certified acceptance-side minimum
\newcommand{\Itset}{\mathfrak{I}}           % finite index set of iterates (C22)

% --- reductions ----------------------------------------------------------------
\newcommand{\Bsys}{B}                        % physical (pre-squashing) system
\newcommand{\Tcl}[1]{\mathcal{T}(#1)}        % trace-class operators on H(#1)
\newcommand{\Ins}{\mathcal{I}}               % generalized-instrument branch map
\newcommand{\redNmsm}{\mathcal{N}}              % POVM measurement channel
\newcommand{\Eprot}{\mathcal{E}}             % n-round protocol channel
\newcommand{\Grnd}{\mathsf{G}}               % one round map of a protocol
\newcommand{\Kfin}{\mathsf{K}}               % end-of-protocol processing
\newcommand{\Crit}{\mathsf{sec}}             % value of a security criterion
\newcommand{\Pflag}{\Pi_{\Flag}}             % flag-sector projector on Q
\newcommand{\Pleq}{\Pi_{\leqN}}              % preserved-block projector
\newcommand{\Pgt}{\Pi_{\gtN}}                % above-cutoff projector
\newcommand{\Sph}{S^{\mathrm{ph}}}           % family of physical inputs
\newcommand{\Srch}{S^{\mathrm{rch}}}         % reachable squashed states
\newcommand{\Ssq}{S^{\sq}}                   % squashed admissible set
\newcommand{\Ffss}{F^{\mathrm{FSS}}}         % flag-state-squasher target
\newcommand{\Fwp}{F^{\mathrm{WP}}}           % multi-bound WPFSS target
\newcommand{\Mwp}{M^{\mathrm{WP}}}           % renormalized above-cutoff POVM
\newcommand{\Mdec}{{M'}^{\mathrm{ph}}}       % declared physical family on B
\newcommand{\Ppvm}{P^{\Flag}}                % flag PVM
\newcommand{\Lmp}{{\Lambda_{\mathrm{mp}}}}   % measure-and-prepare channel
\newcommand{\Lwp}{{\Lambda^{\mathrm{WP}}}}   % multi-bound WPFSS channel
\newcommand{\Lfss}{{\Lambda^{\mathrm{FSS}}}} % flag-state-squasher channel
\newcommand{\lmax}{\lambda_{\max}}           % maximum eigenvalue
\newcommand{\Lnoise}{\widetilde{\Lambda}}    % flag-noise channel
\newcommand{\cg}{\varpi}                     % declared coarse-graining map
\newcommand{\MCK}{\mathsf{MC}}               % multiclick raw-pattern set
\newcommand{\redChain}{\mathfrak{C}}            % reduction chain
\newcommand{\Lnk}{\ell}                      % a single typed link
\newcommand{\dsq}{\preceq}                      % the squashing relation
\newcommand{\opn}[1]{\left\lVert #1\right\rVert}           % operator (spectral) norm
\newcommand{\mxn}[1]{\left\lVert #1\right\rVert_{\max}}    % entrywise max modulus
\newcommand{\Fld}{\mathbb{F}}                   % exact ordered substrate field
\newcommand{\Fldi}{\mathbb{F}_{i}}              % its Gaussian extension F + iF
\newcommand{\lat}[1]{\big\lfloor #1\big\rfloor_{\kappa}} % round DOWN onto the lattice

% --- protocol-compiler ---------------------------------------------------------
\newcommand{\Spec}{\mathfrak{S}}             % protocol specification (compiler input)
\newcommand{\pcChain}{\mathfrak{R}}            % reduction chain (two typed segments)
\newcommand{\Chsrc}{\pcChain_{\mathrm{src}}}   % its Alice-side source-map segment
\newcommand{\Chrec}{\pcChain_{\mathrm{rec}}}   % its receiver-side squashing segment
\newcommand{\rvirt}{\rho^{\mathrm{v}}}       % virtual (PRE-source-map) signal state
\newcommand{\Fph}{F^{\mathrm{ph}}}           % physical (pre-reduction) POVM element
\newcommand{\Vph}{\mathcal{V}^{\mathrm{ph}}} % physical outcome alphabet
\newcommand{\Tpub}{\mathcal{T}_{\mathrm{pub}}}% public-choice alphabet
\newcommand{\Uset}{\mathcal{U}}              % residual outcome alphabet
\newcommand{\idxm}{\mathsf{idx}}             % label-to-index map
\newcommand{\Asig}{A'}                       % transmitted (signal) system -- the REAL one
\newcommand{\Avirt}{A''}                     % virtual transmitted system (source descriptor of Xi_1)
\newcommand{\Abar}{\bar{A}}                  % Alice's measured register
\newcommand{\Ahat}{\hat{A}}                  % Alice's shield register
\newcommand{\pcNmsm}{\mathcal{N}_{\mathsf A}}  % Alice's measurement channel (see the RULES block)
\newcommand{\Phip}{\Phi^{+}}                 % maximally entangled vector label
\newcommand{\Smap}{\mathfrak{s}}             % the emitted sector map (fraktur: a compile-time object)

% --- receiver ------------------------------------------------------------------
\newcommand{\Modes}{\mathbb{M}}              % declared ordered mode set
\newcommand{\Inp}{\mathbb{I}}                % declared input modes
\newcommand{\InpC}{\mathbb{I}_{\Chan}}       % channel input modes
\newcommand{\InpS}{\mathbb{I}_{\Prb}}        % probe (source) input modes
\newcommand{\Env}{\mathbb{E}}                % dilation (environment) modes
\newcommand{\Port}[1]{\mathbb{P}_{#1}}       % mode set of detector port #1
\newcommand{\Emd}[1]{E_{#1}}                 % coordinate projector on modes #1
\newcommand{\Gmat}{G}                        % instrument (transfer) matrix
\newcommand{\Fock}{\mathfrak{F}}             % Fock / second-quantization functor
\newcommand{\Symn}[1]{\mathrm{Sym}^{#1}}     % symmetric tensor power
\newcommand{\normord}[1]{{:}\mspace{2mu}#1\mspace{2mu}{:}}   % normal ordering
\newcommand{\rk}{\operatorname{rank}}        % rank (physics.sty owns \rank)
\newcommand{\SCK}{\mathsf{SC}}               % single-click pattern set
\newcommand{\Recv}{\mathfrak{R}}             % receiver instrument tuple
\newcommand{\Cat}{\mathfrak{C}}              % component catalogue
\newcommand{\Grph}{\mathfrak{E}}             % declared element graph
\newcommand{\defic}{\mathbb{L}}              % loss deficiency operator
\newcommand{\Pvac}[1]{\Pi_{\vac}(#1)}        % joint no-click element on ports #1
\newcommand{\pmc}{p_{\mc}}                   % uniform high-photon mc weight
\newcommand{\Bnd}{\beta}                     % the scalar bound function beta(n)
\newcommand{\Rot}{R}                         % planar rotation matrix
\newcommand{\Wplate}{W}                      % retarder (wave-plate) matrix
\newcommand{\Bsp}{B}                         % beamsplitter core matrix

% --- statistics ----------------------------------------------------------------
\newcommand{\honest}{p^{\mathrm{hon}}}          % honest per-round distribution
\newcommand{\ecomp}{\varepsilon_{\mathrm{comp}}}% completeness budget
\newcommand{\simplex}[1]{\Delta(#1)}            % probability simplex
\newcommand{\Chatcon}{\Chatset_{\mathrm{con}}}  % constrained outcomes
\newcommand{\Ireg}[3]{\mathcal{I}_{#1}(#2,#3)}  % regularized incomplete beta
\newcommand{\Ncount}[1]{\mathsf{N}_{#1}}        % per-outcome count
\newcommand{\qlo}{\mathbf{q}^{-}}               % box lower bounds
\newcommand{\qhi}{\mathbf{q}^{+}}               % box upper bounds

\makeatletter
\newcommand*{\ClaimHeadFont}{\bfseries}   % switched by prominence
\newcommand*{\ClaimNoteFont}{\mdseries}
\newcommand*{\ClaimShortId}{}
\newcommand*{\SetShortId}[1]{\renewcommand*{\ClaimShortId}{#1}}
\newcommand*{\RestatedName}{}
% Upright bodies throughout: italic is reserved for non-normative text.
% \@endparpenalty\@M keeps the meta line on the statement's page.
\newtheoremstyle{claim}{\topsep}{\topsep}{\normalfont\@endparpenalty\@M}{}%
  {\ClaimHeadFont}{.}{.5em}{\thmname{#1}\thmnumber{ #2}\thmnote{ {\ClaimNoteFont(#3)}}}
\newtheoremstyle{claimass}{\topsep}{\topsep}{\normalfont\@endparpenalty\@M}{}%
  {\ClaimHeadFont}{.}{.5em}{\thmname{#1}\thmnumber{ #2}%
  \ifx\ClaimShortId\@empty\else\ [\textsf{\ClaimShortId}]\fi%
  \thmnote{ {\ClaimNoteFont(#3)}}}
\newtheoremstyle{claimremark}{.5\topsep}{.5\topsep}{\small\@endparpenalty\@M}{}%
  {\itshape}{.}{.5em}{\thmname{#1}\thmnumber{ #2}\thmnote{ {\upshape(#3)}}}
\newtheoremstyle{claimreq}{\topsep}{\topsep}{\normalfont\@endparpenalty\@M}{}%
  {\color{reqAccent}\scshape}{.}{.5em}%
  {\thmname{#1}\thmnumber{ #2}\thmnote{ {\upshape(#3)}}}
\makeatother
\theoremstyle{claim}
\newtheorem{theorem}{Theorem}[section]
\newtheorem{lemma}[theorem]{Lemma}
\newtheorem{proposition}[theorem]{Proposition}
\newtheorem{corollary}[theorem]{Corollary}
\newtheorem{fact}[theorem]{Fact}
\newtheorem{definition}[theorem]{Definition}
\newtheorem*{restatedclaim}{\RestatedName}
\theoremstyle{claimass}
\newtheorem{assumption}[theorem]{Assumption}
\newtheorem*{restatedassumption}{\RestatedName}
\theoremstyle{claimremark}
\newtheorem{remark}[theorem]{Remark}
\newtheorem{example}[theorem]{Example}
\theoremstyle{claimreq}
\newtheorem{requirement}[theorem]{Requirement}
\newtheorem*{restatedrequirement}{\RestatedName}
%% --- 8. Prominence: main results a margin rule and a bold heading ----------
\newtcolorbox{mainclaim}{enhanced jigsaw, breakable, blanker,
  borderline west={1.5pt}{-8pt}{mainRule}, left=0pt, right=0pt, top=0pt, bottom=0pt,
  boxsep=0pt, before skip=\topsep, after skip=\topsep, parbox=false,
  before upper={\renewcommand*{\ClaimNoteFont}{\bfseries}}}
\newenvironment{technicalclaim}{\renewcommand*{\ClaimHeadFont}{\mdseries\color{techGrey}}%
  \renewcommand*{\ClaimNoteFont}{}}{}
%% --- 9. Non-normative text, the meta line, navigation ------------------------
\makeatletter
% Tie what follows to this paragraph: no page break, no list penalty before
% the next theorem head (the \@nobreak route a section heading uses).
\newcommand{\keepwithnext}{\par\nobreak\vspace{3pt}\@nobreaktrue
  \everypar{\if@nobreak\@nobreakfalse\clubpenalty\@M\else\clubpenalty\@clubpenalty\everypar{}\fi}}
\newcommand{\claimsummary}[2]{% #1 claim id (named destination), #2 summary
  \par\addvspace{\bigskipamount}\Needspace{6\baselineskip}%
  {\small\itshape\color{nnGrey}\noindent\hypertarget{claim-#1}{}#2\par}\keepwithnext}
\newcommand{\claimanchor}[1]{% a claim without a summary: the destination alone
  \par\addvspace{\bigskipamount}\Needspace{6\baselineskip}%
  \phantomsection\hypertarget{claim-#1}{}\keepwithnext}
\newcommand{\claimleadin}[1]{{\color{nnGrey}\noindent#1\par}\keepwithnext}
\newcommand{\claimmeta}[1]{\par\nobreak\vspace{2pt}%
  {\footnotesize\color{metaGrey}\raggedright\interlinepenalty\@M\noindent#1\par}%
  \addvspace{\medskipamount}}
% After an assumption: the sentence generated from its verification field; after
% any claim with one: its plain-text reading (user_statement). Both
% non-normative, both grey.
\newcommand{\claimverification}[1]{\par\nobreak\vspace{2pt}%
  {\footnotesize\color{metaGrey}\noindent#1\par}\nobreak}
\newcommand{\claimreading}[1]{\par\nobreak\vspace{2pt}%
  {\footnotesize\color{nnGrey}\noindent\emph{Reading:}\enspace#1\par}\nobreak}
\makeatother
\newcommand{\claimbookmark}[4]{% #1 id, #2 kind word, #3 label, #4 plain title
  \bookmark[dest=claim-#1,level=2]{#2 \getrefnumber{#3} #4}}
\newcommand{\metasep}{\unskip\enspace\textperiodcentered\enspace\ignorespaces}
\newcommand{\metalabel}[1]{\emph{#1}\enspace}
\newcommand{\sectionintro}[1]{{\color{nnGrey}\noindent#1\par}\smallskip}
\newcommand{\sectionimports}[1]{{\footnotesize\color{metaGrey}\raggedright\noindent\metalabel{Imports}#1\par}\smallskip}
\newenvironment{overview}{\par\color{nnGrey}\itshape\noindent\ignorespaces}{\par}
% References: the whole "Def. 2.1" is the link; ids are the fallback text of
% every link into another PDF (a relative URI to a named destination).
\newcommand{\mref}[2]{\hyperref[#2]{#1~\ref*{#2}}}
\newcommand{\aref}[2]{\hyperref[#2]{\textsf{#1}}}
\DeclareUrlCommand{\claimid}{\urlstyle{sf}\def\UrlBreaks{\do\/\do\:\do\-\do\.\do\_}%
  \def\UrlBigBreaks{}\def\UrlNoBreaks{}}
\newcommand{\extref}[2]{\hyperref{#1}{}{claim-#2}{\claimid{#2}}}
\newcommand{\extaref}[3]{\hyperref{#1}{}{claim-#2}{\textsf{#3}}}
\newcommand{\coderef}[2]{\href{#1}{\claimid{#2}}}
% Claim-body commands (spec/claim-format.md section 4): \xref{<full id>} prints
% the kind, title and id of a claim of another contract, linked to its stable
% destination in that contract's PDF (one \DeclareXref per cited claim below);
% \term[gloss]{phrase} marks a defined phrase for the index (the gloss feeds
% the Terms table); \tref[phrase]{text} is a use of a defined term. The
% assembler rewrites both in the body into \termdef (the definition's named
% destination term-<slug>) and \termref / \exttermref (a link to it, in this
% or another contract's PDF); the fallbacks below serve text it copies
% elsewhere, such as the digest. \iface is an interface key.
\makeatletter
\newcommand{\DeclareXref}[4]{% id, kind word, title, contract PDF
  \@namedef{xref@#1}{\hyperref{#4}{}{claim-#1}{#2 ``#3'' ({\small\claimid{#1}})}}}
\newcommand{\xref}[1]{\@ifundefined{xref@#1}%
  {\textbf{[unresolved \nolinkurl{#1}]}\@latex@warning{Unresolved xref #1}}%
  {\@nameuse{xref@#1}}}
\makeatother
\DeclareUrlCommand{\iface}{\urlstyle{tt}}
\newcommand{\term}[2][]{\emph{#2}\index{#2}}
\newcommand{\tref}[2][]{#2}
\newcommand{\termdef}[2]{\hypertarget{#1}{}\emph{#2}\index{#2}}
\newcommand{\termref}[2]{\hyperlink{#1}{#2}}
\newcommand{\exttermref}[3]{\hyperref{#1}{}{#2}{#3}}
\hypersetup{pdftitle={Protocol compiler (protocol-compiler) -- statements digest 0.24.0}}

% The digest: the contract's labels under the prefix C-, its numbers checked.
\externaldocument[C-]{contract}
\theoremstyle{claim}\newtheorem*{digestclaim}{\RestatedName}
\theoremstyle{claimass}\newtheorem*{digestassumption}{\RestatedName}
\theoremstyle{claimreq}\newtheorem*{digestrequirement}{\RestatedName}
\makeatletter
% #1 label, #2 the number the assembler computed: a local label carrying it, so
% that a reference between two digest items stays in the digest; the build
% stops when the contract's own number differs.
\newcommand{\digestlabel}[2]{\phantomsection
  \edef\digest@theirs{\getrefnumber{C-#1}}\def\digest@ours{#2}%
  \ifx\digest@theirs\digest@ours\else
    \PackageError{digest}{#1 is #2 here but \digest@theirs\space in the contract}{}\fi
  \def\@currentlabel{#2}\label{#1}}
\makeatother
% A display's local label (amsmath records the \tag, i.e. the contract's number).
\newcommand{\digesteqlabel}[1]{\label{#1}}

% Claims of other contracts cited here (\xref, [[...]]).
\DeclareXref{foundations/def:binary64}{Definition}{Binary64 numbers, declared values and dyadic lattices}{../foundations/contract.pdf}
\DeclareXref{foundations/def:declared-data}{Definition}{Declared data and the parties to a module boundary}{../foundations/contract.pdf}
\DeclareXref{foundations/def:declared-order}{Definition}{Sums in a declared order and sums that round nowhere}{../foundations/contract.pdf}
\DeclareXref{foundations/def:exact-admissibility}{Definition}{Exact admissibility of a declared operator family}{../foundations/contract.pdf}
\DeclareXref{foundations/def:exact-island}{Definition}{The exact island}{../foundations/contract.pdf}
\DeclareXref{foundations/def:factor-enclosure}{Definition}{Factor enclosure of a block}{../foundations/contract.pdf}
\DeclareXref{foundations/def:round-homogeneity}{Definition}{Round homogeneity of a per-round family}{../foundations/contract.pdf}
\DeclareXref{foundations/def:sector-map}{Definition}{Sector maps and block-list operators}{../foundations/contract.pdf}
\DeclareXref{foundations/lem:exact-accumulation}{Lemma}{Exactness of a sum that rounds nowhere}{../foundations/contract.pdf}
\DeclareXref{foundations/rem:computed-completeness}{Remark}{A complete computed sum from an incomplete family}{../foundations/contract.pdf}
\DeclareXref{foundations/rem:summation-order}{Remark}{Dependence of a computed sum on its order}{../foundations/contract.pdf}
\DeclareXref{foundations/req:boundary-admission}{Requirement}{Exact admissibility at every module boundary}{../foundations/contract.pdf}
\DeclareXref{foundations/req:consumer-decides-exactly}{Requirement}{Exact decision by every consumer}{../foundations/contract.pdf}
\DeclareXref{foundations/req:declared-lattice}{Requirement}{The declared lattice condition}{../foundations/contract.pdf}
\DeclareXref{foundations/req:declared-order}{Requirement}{Summation order declared with the data}{../foundations/contract.pdf}
\DeclareXref{foundations/req:exact-emission}{Requirement}{Exact verification before emission}{../foundations/contract.pdf}
\DeclareXref{foundations/req:exact-island}{Requirement}{Exact arithmetic on the island, unbudgeted rounding elsewhere}{../foundations/contract.pdf}
\DeclareXref{foundations/req:tolerance-scope}{Requirement}{Scope of the declared tolerance}{../foundations/contract.pdf}
\DeclareXref{foundations/req:who-snaps}{Requirement}{Snapping by the declarer only}{../foundations/contract.pdf}
\DeclareXref{meat-core/ass:lab}{Assumption}{Each party's laboratory is sealed}{../meat-core/contract.pdf}
\DeclareXref{meat-core/cor:composed-security}{Corollary}{Security of the protocol and its net key}{../meat-core/contract.pdf}
\DeclareXref{meat-core/cor:reduction-constrained}{Corollary}{The reduction for arbitrary constraint lists}{../meat-core/contract.pdf}
\DeclareXref{meat-core/def:admissible-declaration}{Definition}{Admissible declaration}{../meat-core/contract.pdf}
\DeclareXref{meat-core/def:declaration}{Definition}{The protocol declaration as the MEAT core reads it}{../meat-core/contract.pdf}
\DeclareXref{meat-core/def:roundmodel}{Definition}{Round model and licensed attack sets}{../meat-core/contract.pdf}
\DeclareXref{meat-core/prop:a-neutral}{Proposition}{Value neutrality along Alice's sectors}{../meat-core/contract.pdf}
\DeclareXref{meat-core/prop:reference-companion}{Proposition}{The compiled reference companion is an admissible declaration}{../meat-core/contract.pdf}
\DeclareXref{meat-core/prop:trivial-alice-system}{Proposition}{The two source modes coincide at a trivial Alice system}{../meat-core/contract.pdf}
\DeclareXref{meat-core/rem:cgate}{Remark}{Degeneracy of the licensed-class half of the gate}{../meat-core/contract.pdf}
\DeclareXref{meat-core/req:companion-admission}{Requirement}{Admission of the reference companion beside its declaration}{../meat-core/contract.pdf}
\DeclareXref{meat-core/thm:reduction}{Theorem}{Reduction of the rate to the convex program}{../meat-core/contract.pdf}
\DeclareXref{receiver/def:clickalphabet}{Definition}{Click-pattern alphabet}{../receiver/contract.pdf}
\DeclareXref{receiver/def:coarse}{Definition}{Declared coarse-graining and the emitted POVM}{../receiver/contract.pdf}
\DeclareXref{receiver/thm:coarse-graining}{Theorem}{Coarse-graining of the click-pattern POVM}{../receiver/contract.pdf}
\DeclareXref{reductions/ass:meas}{Assumption}{Each measuring party's device as one POVM per round}{../reductions/contract.pdf}
\DeclareXref{reductions/ass:tsm}{Assumption}{Output-state test statistics}{../reductions/contract.pdf}
\DeclareXref{reductions/cor:entries}{Corollary}{The emitted constraint entries}{../reductions/contract.pdf}
\DeclareXref{reductions/def:attset}{Definition}{Admissible attack sets in the consumer's format}{../reductions/contract.pdf}
\DeclareXref{reductions/def:chain}{Definition}{Reduction chain}{../reductions/contract.pdf}
\DeclareXref{reductions/def:joint-family}{Definition}{Joint alphabet, classes and the joint family of a private setting}{../reductions/contract.pdf}
\DeclareXref{reductions/def:licensed}{Definition}{Licensed entries and the party-local form}{../reductions/contract.pdf}
\DeclareXref{reductions/def:link}{Definition}{Typed link}{../reductions/contract.pdf}
\DeclareXref{reductions/def:roundmodel}{Definition}{Round model and attack sequences, typed by source mode}{../reductions/contract.pdf}
\DeclareXref{reductions/def:squashing}{Definition}{Squashing relation}{../reductions/contract.pdf}
\DeclareXref{reductions/domination/def:domination-link}{Definition}{The domination link}{../reductions/contract.pdf}
\DeclareXref{reductions/ex:identity}{Example}{The identity link}{../reductions/contract.pdf}
\DeclareXref{reductions/lem:linear}{Lemma}{Linearity: post-processing, private settings and mixtures}{../reductions/contract.pdf}
\DeclareXref{reductions/lem:party-local}{Lemma}{Party-local transport}{../reductions/contract.pdf}
\DeclareXref{reductions/lem:precomposition}{Lemma}{Licensed admissibility survives precomposition with a source channel}{../reductions/contract.pdf}
\DeclareXref{reductions/lem:row-relaxation}{Lemma}{A rounded constraint entry with its constant relaxed}{../reductions/contract.pdf}
\DeclareXref{reductions/lem:source-transfer}{Lemma}{Exactness of the source map in trusted-source mode}{../reductions/contract.pdf}
\DeclareXref{reductions/rem:intertwining}{Remark}{Classical maps before a squasher, and relabelling after one}{../reductions/contract.pdf}
\DeclareXref{reductions/rem:ordering}{Remark}{Recorded stage order, immaterial to soundness for licensed lists}{../reductions/contract.pdf}
\DeclareXref{reductions/req:factorised-emission}{Requirement}{Emission of several measuring parties as factors}{../reductions/contract.pdf}
\DeclareXref{reductions/req:party-local-tag}{Requirement}{The party-local tag and its scope}{../reductions/contract.pdf}
\DeclareXref{reductions/req:reference-companion}{Requirement}{Emitting the reference companion beside the artefact}{../reductions/contract.pdf}
\DeclareXref{reductions/thm:chain}{Theorem}{Composition of a chain}{../reductions/contract.pdf}
\DeclareXref{reductions/thm:emit}{Theorem}{Exact emission of a reduction instance}{../reductions/contract.pdf}
\DeclareXref{reductions/thm:emitted-product}{Theorem}{The emission of several measuring parties}{../reductions/contract.pdf}
\DeclareXref{reductions/thm:licence}{Theorem}{The attack-set licence}{../reductions/contract.pdf}
\DeclareXref{reductions/thm:transfer}{Theorem}{Security transfer from a constrained squashed model}{../reductions/contract.pdf}
\DeclareXref{reductions/thm:two-sided}{Theorem}{The two-sided reduction in trusted-source mode}{../reductions/contract.pdf}
\DeclareXref{reductions/wpfss/prop:flat-instance}{Proposition}{The flat instance of the post-processed receiver}{../reductions/contract.pdf}
\DeclareXref{transmitter/ass:ens}{Assumption}{Emission of the specified ensemble}{../transmitter/contract.pdf}
\DeclareXref{transmitter/ass:seq}{Assumption}{Sequential round structure}{../transmitter/contract.pdf}
\DeclareXref{transmitter/ass:src}{Assumption}{Retention of the system A in trusted-source mode}{../transmitter/contract.pdf}
\DeclareXref{transmitter/decoy/cor:emitted-decoy}{Corollary}{The emitted decoy declaration}{../transmitter/contract.pdf}
\DeclareXref{transmitter/def:source-segment}{Definition}{Source segment and its interface conditions}{../transmitter/contract.pdf}
\DeclareXref{transmitter/def:source-spec}{Definition}{Source specification}{../transmitter/contract.pdf}
\DeclareXref{transmitter/def:srcert}{Definition}{Source-replacement certificate}{../transmitter/contract.pdf}
\DeclareXref{transmitter/fact:sr}{Fact}{Source replacement for prepare-and-measure protocols}{../transmitter/contract.pdf}
\DeclareXref{transmitter/lem:ensemble-presentation}{Lemma}{Certificate defined by an ensemble presentation}{../transmitter/contract.pdf}
\DeclareXref{transmitter/rem:certificate-forms}{Remark}{Purity and exact declarability of the certificate forms}{../transmitter/contract.pdf}
\DeclareXref{transmitter/rem:da1}{Remark}{The one-dimensional retained system}{../transmitter/contract.pdf}
\DeclareXref{transmitter/rem:segment-scope}{Remark}{Scope of the source-segment conditions}{../transmitter/contract.pdf}
\DeclareXref{transmitter/rem:source-mode}{Remark}{The two source modes}{../transmitter/contract.pdf}
\DeclareXref{transmitter/req:attestation-record}{Requirement}{Pass-through of the retention attestation}{../transmitter/contract.pdf}
\DeclareXref{transmitter/req:failing-certificate}{Requirement}{Refusal of a certificate that fails a decided condition}{../transmitter/contract.pdf}
\DeclareXref{transmitter/req:licence-scope}{Requirement}{Certificate and attestation demanded only at $d_A\ge2$}{../transmitter/contract.pdf}
\DeclareXref{transmitter/req:marginal-fill}{Requirement}{The unique marginal served at $d_A=1$}{../transmitter/contract.pdf}
\DeclareXref{transmitter/thm:certificate-sufficiency}{Theorem}{Checkability and sufficiency of the source-replacement certificate}{../transmitter/contract.pdf}

\begin{document}
\begin{center}
{\LARGE\bfseries Protocol compiler (protocol-compiler)}\\[4pt]
{\large Statements digest, version 0.24.0}
\end{center}
\noindent{\small Every definition, fact, assumption, result and requirement of the \href{contract.pdf}{contract}, with the contract's numbers; remarks, examples, technical steps and all proofs are left out. Each heading opens its statement in the contract.}

\part*{Mathematics}\pdfbookmark[0]{Mathematics}{digest-part-1}

\section*{1\quad The protocol declaration}\pdfbookmark[1]{1 The protocol declaration}{digest-group-protocol-compiler/def:declaration}

% claim protocol-compiler/def:declaration
\digestlabel{def:declaration}{1.1}\pdfbookmark[2]{Definition 1.1 Protocol declaration}{digest-protocol-compiler/def:declaration}
\renewcommand*{\RestatedName}{\hyperref{contract.pdf}{}{claim-protocol-compiler/def:declaration}{Definition~1.1}}%
\begin{digestclaim}[Protocol declaration]
A \emph{protocol declaration}
is a tuple
\begin{equation*}\tag{\getrefnumber{C-eq:declaration}}\digesteqlabel{eq:declaration}
  \Decl=\Big(A,\ \tau_A;\ \Qsys,\ \SecSet;\ \{M_x\}_{x\in\Xset};\
  \{F_v\}_{v\in\Vset};\ \Party;\ a_A,\ a_B;\ \phi;\ r,\ \hold;\
  \CList\Big),
\end{equation*}
with components as follows. $A$ and $\Qsys$ are finite-dimensional systems,
$\Hs{\Qsys}=\bigoplus_{s\in\SecSet}\Hs{\Qsys_s}$ with a finite index set
$\SecSet$ of sectors, sector projectors $\Pi_s$ and dimensions
$d_s=\dim\Hs{\Qsys_s}$, $d_{\Qsys}=\sum_sd_s$, $d_A=\dim\Hs{A}$. The
sector map $\Smap$ is the ordered list
$\big((s,d_s)\big)_{s\in\SecSet}$ of block identifiers with their dimensions,
together with its block kind, a label naming what the identifiers mean
(\xref{foundations/def:sector-map}). It is one object, carried with the
declaration beside the slot $\Qsys,\SecSet$ of \eqref{eq:declaration}, and it
is what fixes $\SecSet$, its total order and $(d_s)_s$.
The
\emph{marginal slot}
holds either a state $\tau_A\in\Seq{A}$, in trusted-source mode, or
the value $\nomarg$, in untrusted-source mode.
$\{M_x\}_{x\in\Xset}$ is a POVM on $A$, given as one matrix per $x$, since $A$
carries no sector decomposition. The words ``state'' and ``POVM'' in the two
sentences above name the intended objects, and they are also exactly the
conditions this contract decides: what a declaration is required to satisfy
is D(i)--D(viii) of Definition~\ref{def:target}, every operator predicate of
which is decided \emph{exactly} on the declared binary64 numbers,
together with the summation condition of Definition~\ref{def:exact}, without
which a computed sum would not be the sum it is read as.

\emph{Operators on $\Qsys$ are block lists.} An operator
on $\Qsys$ is given as one array $F^{(s)}\in\Lin{\Qsys_s}$ per
$s\in\SecSet$, against the sector map $\Smap$, and that list is the
operator: it is a block list against $\Smap$, the operator it denotes is
$\bigoplus_{s}F^{(s)}$, and the
declaration holds no array larger than a single sector block. Where
$|\SecSet|=1$, say $\SecSet=\{s_0\}$, that one block \emph{is} the whole of
$\Qsys$ and is held and emitted as such, $d_{s_0}=d_{\Qsys}$, as on the
Band-A route (Theorem~\ref{thm:band-a}\,(i)). Sector block-diagonality is therefore a
property of the \emph{type}: there is no off-sector entry an operator of this
form could carry, and correspondingly nothing to verify. So
$\{F_v\}_{v\in\Vset}$ is a family of such block lists over a finite alphabet,
one list per $v$; $F_v^{(s)}$ denotes its $s$-block. A state $\omega$ on
$A\Qsys$ is likewise a list, one array $\omega^{(s)}\in\Lin{A\Qsys_s}$ per
$s$, and $\omega=\bigoplus_s\omega^{(s)}$ in the same sense. $\Party$ is a
\emph{party partition}:
two total maps $\rd{\mathsf A}:\Xset\times\Vset\to\Oset{\mathsf A}$ and
$\rd{\mathsf B}:\Xset\times\Vset\to\Oset{\mathsf B}$ onto declared finite
alphabets, each a
\emph{readout map},
together with a party label for $A$ and, for each $s$, a declared tensor
factorization of $\Hs{\Qsys_s}$ with a party label per factor.
$a_A:\Oset{\mathsf A}\to\CAset$ and $a_B:\Oset{\mathsf B}\to\CBset$, each an
\emph{announcement map},
are written jointly
\begin{equation*}\tag{\getrefnumber{C-eq:jointann}}\digesteqlabel{eq:jointann}
  a(x,v)\coloneqq\big(a_A(\rd{\mathsf A}(x,v)),\,a_B(\rd{\mathsf B}(x,v))\big)
  \in\Kann\coloneqq\CAset\times\CBset ,
\end{equation*}
$\phi:\Kann\to\Chatset$ is the acceptance-statistic map,
$r:\Xset\times\Vset\to\Zset$ the raw-key map with
$\Zset=\Zval\sqcup\{\perp\}$ and $\Zval=\{0,\dots,d_Z-1\}$,
$\hold\in\{\mathsf A,\mathsf B\}$ the declared key holder, and
$\CList=\{(\hat O_k,\bowtie_k,\mathrm{rhs}_k)\}_{k=1}^{m}$, $m\ge0$, a finite
constraint list, each entry an operator $\hat O_k$, a relation symbol
$\bowtie_k$ and a right-hand side $\mathrm{rhs}_k$ (which symbols a
well-formed declaration admits is D(vi) of Definition~\ref{def:target}),
with every $\hat O_k$ a block list against the same $\Smap$, Hermitian block
by block.

\emph{The right-hand side of an entry carries a declared tag, and the tag is
part of the datum rather than a property measured of it.} In trusted-source
mode an entry is declared either homogeneous, carrying a real scalar
$c_k$, from which the right-hand side it denotes is \emph{constructed} as
$\Theta_k\coloneqq c_k\tau_A$, or general, carrying a Hermitian matrix
$\Theta_k$ on $A$. In untrusted-source mode only the scalar form exists and
the right-hand side is a real number $\Theta_k=c_k$. So an entry is written
$(\hat O_k,\bowtie_k,c_k)$ exactly when it is declared homogeneous or is
untrusted-mode, and $(\hat O_k,\bowtie_k,\Theta_k)$ when it is declared
general (Definition~\ref{def:homog}). For these two forms the tag
is the type of the right-hand-side datum, a real scalar or a Hermitian
matrix, which is why Definition~\ref{def:homog} calls them the untagged
forms; the third form, party-local, exists only in the general format below
and is declared by its tag (DX6). We write $\kappa=(c_A,c_B)\in\Kann$, a
value of the joint announcement \eqref{eq:jointann}, for an announcement
class and
\begin{equation*}\tag{\getrefnumber{C-eq:chkey}}\digesteqlabel{eq:chkey}
  \Chkey\coloneqq\big\{\hat c\in\Chatset:\ \exists\,\kappa\in\phi^{-1}(\hat c)
  \text{ with }\varsigma(\kappa)=\mathrm{pass}\big\}
\end{equation*}
for the
\emph{key-live}
acceptance values. Here $\varsigma$ is the sifting predicate of the
specification's public part (Definition~\ref{def:spec}\,(d)), read as
$\varsigma(\kappa)=\mathrm{fail}$ on every class
$\kappa\notin a(\Xset\times\Vset)$, which no outcome pair realises; on the
realised classes it is, under the hypotheses of
Theorem~\ref{thm:sifting-and-key}, a witness of D(v) of
Definition~\ref{def:target}
(Theorem~\ref{thm:sifting-and-key}\,(ii)). It is \emph{not} a slot of
\eqref{eq:declaration}, so \eqref{eq:chkey} is the set of acceptance values
to which some realised class passing sifting maps, fixed by the
specification and not by the declaration; the key-live set of the consuming
engine, \xref{meat-core/def:declaration}, takes the same convention on the
classes no outcome pair realises.

\emph{The general format.} Everything above is written for the format in
which $A$ carries one sector, one system is measured, Alice's measurement is
declared as the POVM $\{M_x\}$, the public discussion is two messages neither
of which reads the other, and no entry is declared party-local; we
call it the
\emph{basic format}.
It is the case $|\mathbb{T}|=1$, $|\mathbb{F}|=1$, POVM form, $L=2$ of the
following one, \emph{entry for entry} on every component a derivation reads:
at those values every such object defined in (DX1)--(DX6) is the object of
the same name above, and nothing else is declared. On the party labels of the
tensor factors of each sector block, which no derivation of this module reads
and which the consuming engine's laboratory premise reads as naming the
parties' laboratories, the general format is narrower than the basic one. In general a declaration is the
tuple above with its slots read as follows.
\begin{enumerate}[label=(DX\arabic*),leftmargin=3.2em]
\item \emph{(factorised sector map)} The sector map is an ordered list, each
entry of which is a
\emph{factor}.
Each factor carries a
\emph{system tag},
$A$ or $\Qsys$, a party label, and a sector map of its own, an ordered list of
block identifiers with their dimensions; all factors carry one block kind. At
most one factor is tagged $A$, and where there is one it is the first of the
list; at least one factor is tagged $\Qsys$. The factor tagged $A$ fixes a
finite ordered index set $\mathbb{T}$ and a decomposition
$\Hs{A}=\bigoplus_{t\in\mathbb{T}}\Hs{A_t}$ into subspaces spanned by vectors
of the declared basis of $A$, with projectors $\Pi^{A}_t$ and dimensions
$d_{A_t}$, $d_A=\sum_td_{A_t}$; where no factor is tagged $A$, the system
$A$ is still present and is one Alice sector, $\mathbb{T}=\{t_0\}$ and
$A_{t_0}=A$. A factor tagged $\Qsys$ is a
\emph{measured factor},
and the measured factors form the ordered set $\mathbb{F}$;
factor $f$ fixes a system $\Qsys_f$ with
$\Hs{\Qsys_f}=\bigoplus_{s_f\in\SecSet^{f}}\Hs{\Qsys_{f,s_f}}$, projectors
$\Pi^{f}_{s_f}$ and dimensions $d_{f,s_f}$, and for $s=(s_f)_{f}$
\begin{equation*}\tag{\getrefnumber{C-eq:productmap}}\digesteqlabel{eq:productmap}
  \Qsys=\bigotimes_{f\in\mathbb{F}}\Qsys_f,\qquad
  \SecSet=\prod_{f\in\mathbb{F}}\SecSet^{f},\qquad
  \Hs{\Qsys_s}=\bigotimes_{f}\Hs{\Qsys_{f,s_f}},\qquad
  \Pi_s=\bigotimes_f\Pi^{f}_{s_f},\qquad d_s=\prod_fd_{f,s_f},
\end{equation*}
all products in the order of $\mathbb{F}$, and $\SecSet$ ordered with the
last factor fastest, a product identifier being the factors' identifiers
joined in that order. Each pair $(t,s)\in\mathbb{T}\times\SecSet$ is a
\emph{joint sector};
the joint sectors are ordered with $t$ slowest, and
$\Hs{A\Qsys}=\bigoplus_{(t,s)}\Hs{A_t}\otimes\Hs{\Qsys_s}$, the $A$ index
first inside a block. The party partition labels $A$ with the party of the
factor tagged $A$, where there is one, and otherwise with the party label it
declares for $A$, as in the basic format; and it labels every tensor factor of
$\Hs{\Qsys_{f,s_f}}$ with the party of $f$.
\item \emph{(Alice's data and $\omega$ are block lists)} $\tau_A$, every
operator of Alice's measurement slot (DX4) and the right-hand side $\Theta_k$
of a general entry are lists of one array on $\Hs{A_t}$ per
$t\in\mathbb{T}$, and the list is the operator:
$\tau_A=\bigoplus_t\tau_A^{(t)}$, $M_x=\bigoplus_tM_x^{(t)}$,
$\Theta_k=\bigoplus_t\Theta_k^{(t)}$. A state $\omega$ on $A\Qsys$ is a list
of one array $\omega^{(t,s)}\in\Lin{A_t\Qsys_s}$ per joint sector,
$\omega=\bigoplus_{(t,s)}\omega^{(t,s)}$. Block diagonality along
$\mathbb{T}$ is a property of the type: no off-sector entry can be written.
By itself this typing licenses nothing; whether restricting $\omega$ to such
lists is sound is not settled by this definition.
\item \emph{(factorised family)} The outcome alphabet is the product
$\Vset=\prod_{f\in\mathbb{F}}\Vset^{f}$ of the factors' ordered alphabets, in
the product order, the last factor fastest. What is declared is one family of
block lists $\{F^{f}_{v_f}\}_{v_f\in\Vset^{f}}$ per measured factor, against
that factor's own sector map, together with the rule
\begin{equation*}\tag{\getrefnumber{C-eq:productfamily}}\digesteqlabel{eq:productfamily}
  F_v\coloneqq\bigotimes_{f\in\mathbb{F}}F^{f}_{v_f},\qquad
  F^{(s)}_v=\bigotimes_{f\in\mathbb{F}}F^{f,(s_f)}_{v_f}
  \qquad\big(v=(v_f)_f,\ s=(s_f)_f\big).
\end{equation*}
No entry of a product is declared. That $\{F_v\}_{v\in\Vset}$ is then a
POVM on $\Qsys$, block diagonal for the product sector map with the blocks of
\eqref{eq:productfamily}, follows for the exact numbers the factors declare:
a tensor product of Hermitian positive semidefinite operators is Hermitian
and positive semidefinite,
$\sum_{v}\bigotimes_fF^{f}_{v_f}=\bigotimes_f\sum_{v_f}F^{f}_{v_f}
=\idon{\Qsys}$ by multilinearity, and the tensor product distributes over the
direct sums $F^{f}_{v_f}=\bigoplus_{s_f}F^{f,(s_f)}_{v_f}$.
\item \emph{(the two forms of Alice's measurement slot)} In
\emph{POVM form}
the slot holds $\{M_x\}_{x\in\Xset}$. In
\emph{ensemble form},
which exists in trusted-source mode only, it holds a family
$\{N_x\}_{x\in\Xset}$ of block lists along $\mathbb{T}$, every block
Hermitian and positive semidefinite, with
\begin{equation*}\tag{\getrefnumber{C-eq:ensembleform}}\digesteqlabel{eq:ensembleform}
  \sum_{x\in\Xset}N^{(t)}_x=\tau^{(t)}_A
  \quad\text{and}\quad\tau^{(t)}_A\ne0
  \qquad\text{for every }t\in\mathbb{T},
\end{equation*}
the sum formed entrywise in the declared order of $\Xset$ and
rounding nowhere, and the second condition a comparison of declared entries with the
datum $0$. An ensemble-form declaration \emph{denotes} the POVM-form
declaration with the same remaining slots and
\begin{equation*}\tag{\getrefnumber{C-eq:ensemblepovm}}\digesteqlabel{eq:ensemblepovm}
  M^{(t)}_x\coloneqq\big(\tau^{(t)}_A\big)^{-1/2}N^{(t)}_x
  \big(\tau^{(t)}_A\big)^{-1/2}
  +\frac{\tr N^{(t)}_x}{\tr\tau^{(t)}_A}\,\big(\idon{A_t}-P_t\big),
\end{equation*}
where $P_t$ is the projector onto the support of $\tau^{(t)}_A$ and the
inverse square root is taken on that support and is zero off it. The
operators \eqref{eq:ensemblepovm} are defined and in general irrational; every statement made of an
ensemble-form declaration is made of the declaration it denotes. For a
POVM-form declaration in trusted-source mode put
$N_x\coloneqq\tau_A^{1/2}M_x\tau_A^{1/2}$; for the declaration an ensemble
form denotes, this returns the declared $N_x$.
\item \emph{(transcript protocol)} The announcement slot holds a
\emph{transcript protocol},
an ordered list of $L\ge1$ entries, each a
\emph{message},
message $m$ carrying a sender $\pi_m\in\{\mathsf A,\mathsf B\}$, a finite
alphabet $\mathcal{C}_m$ and a total table
$a_m:\Oset{\pi_m}\times\mathcal{C}_1\times\dots\times\mathcal{C}_{m-1}
\to\mathcal{C}_m$. The joint announcement function is \emph{derived} from the
tables by composition,
\begin{equation*}\tag{\getrefnumber{C-eq:transcript}}\digesteqlabel{eq:transcript}
  \kappa_m(x,v)\coloneqq a_m\big(\rd{\pi_m}(x,v),\,\kappa_1(x,v),\dots,
  \kappa_{m-1}(x,v)\big),\qquad
  a(x,v)\coloneqq\big(\kappa_m(x,v)\big)_{m=1}^{L}\in
  \Kann\coloneqq\prod_{m=1}^{L}\mathcal{C}_m ,
\end{equation*}
and the acceptance-statistic map $\phi$, the sifting predicate $\varsigma$ of
\eqref{eq:chkey} and the holder's key map $r_{\hold}$ of D(viii) of
Definition~\ref{def:target} are maps on this $\Kann$. The two
maps $a_A$, $a_B$ of the basic format are the transcript protocol with $L=2$,
$\pi=(\mathsf A,\mathsf B)$, $\mathcal{C}_1=\CAset$, $\mathcal{C}_2=\CBset$,
$a_1=a_A$ and $a_2(o,\kappa_1)=a_B(o)$. A message is a function of its
sender's readout and of the earlier messages of the same round, and of
nothing else. No measurement of a round depends on a message of that round,
the measurements being the fixed families of (DX3) and (DX4); processing
across rounds is not a round message.
\item \emph{(party-scoped entries)} Each entry of $\CList$ carries the label
of one measured factor, its
\emph{scope}
$f_k$, and its operator $\hat O_k$ is a block list against that factor's own
sector map, Hermitian block by block. Unless the entry is declared
party-local it denotes the
\emph{lifted operator}
$\idon{}\otimes\hat O_k\otimes\idon{}$ on $\Qsys$, with $\hat O_k$ on the
tensor factor $\Qsys_{f_k}$, whose $s$-block is
$\idon{}\otimes\hat O_k^{(s_{f_k})}\otimes\idon{}$, and it is read in the forms stated
above for that lifted operator. With one measured factor the scope is that
factor and the lifted operator is $\hat O_k$ itself. An entry may instead be
declared party-local, in either source mode, carrying a real $c_k$; its
reading is the third form of Definition~\ref{def:homog}.
\end{enumerate}

\emph{What this module decides of the general format} is
Definition~\ref{def:exact} for the operator data and D(i)--D(viii) of
Definition~\ref{def:target} for the rest; that the family denoted by an
ensemble form is a POVM is Theorem~\ref{thm:sifting-and-key}\,(v).
\end{digestclaim}

% claim protocol-compiler/def:exact
\digestlabel{def:exact}{1.3}\pdfbookmark[2]{Definition 1.3 Exact admissibility of declared operator data}{digest-protocol-compiler/def:exact}
\renewcommand*{\RestatedName}{\hyperref{contract.pdf}{}{claim-protocol-compiler/def:exact}{Definition~1.3}}%
\begin{digestclaim}[Exact admissibility of declared operator data]
Declared operator data are binary64 numbers, each the dyadic rational
it denotes and not an approximation of anything
(\xref{foundations/def:binary64}\,(iii)), and every predicate this module
decides on them is decided \emph{exactly}, on those numbers. This definition lists
those predicates, together with the one condition of the list that is a
condition on the arithmetic that forms a sum rather than a predicate on an
operator.

\emph{The declared summation order.} The compilation map emits a
total order on every declared alphabet and on $\SecSet$: the
label-to-index map $\idxm$, which numbers the letters of
each declared alphabet, and the sectors, consecutively from $0$
(Definition~\ref{def:compile}, Theorem~\ref{thm:typing}); the
declared order of an alphabet is its $\idxm$ order. For a finite
family $\{a_j\}_{j\in J}$ of $n\times n$ matrices indexed by a declared
alphabet $J$ with $\idxm$ order $j_1\prec\dots\prec j_{|J|}$, the
sum in the declared order is accumulated as the
computed sequential sum (\xref{foundations/def:declared-order})
\begin{equation*}\tag{\getrefnumber{C-eq:declorder}}\digesteqlabel{eq:declorder}
  \Sigma_0\coloneqq0,\qquad
  \Sigma_i\coloneqq\Sigma_{i-1}+a_{j_i}\ \ (1\le i\le|J|),\qquad
  \Sigma\coloneqq\Sigma_{|J|},
\end{equation*}
formed left to right, entrywise, in binary64, with no re-ordering and no
re-association; the trace of a matrix is likewise accumulated over its row
index in increasing order, which is the order of the emitted array.

\emph{The five conditions.} Let $\{a_j\}_{j\in J}$ be a finite family of
block lists against a sector map $\Smap$, one array $a_j^{(s)}$ on
$\Hs{\Qsys_s}$ per $s\in\SecSet$. The family is exactly admissible
for $\Smap$ when
\begin{itemize}[leftmargin=2.6em]
\item[(T)] \emph{(sector map)} the sector map the supplied family carries is
  $\Smap$, the map this module holds: the same block kind, the same block
  identifiers in the same order, the same dimensions, so that
  $a_j=\bigoplus_sa_j^{(s)}$ holds \emph{by construction}. This is a type
  condition: it is decided, never measured, and an off-block entry is not
  small but inexpressible;
\item[(S)] \emph{(the sum rounds nowhere)} writing $\Sigma^{(s)}$ for
  the total $\Sigma$ that \eqref{eq:declorder} produces on the block family
  $\{a_j^{(s)}\}_j$: for each $s$, every one of the $|J|$ entrywise binary64
  additions forming it is exact, so that the computed $\Sigma^{(s)}$ \emph{is}
  the real sum $\sum_ja_j^{(s)}$;
\item[(H)] \emph{(Hermiticity)} $a_j^{(s)}=a_j^{(s)\dagger}$ for every $(j,s)$;
\item[(P)] \emph{(positivity)} $a_j^{(s)}\succeq0$ for every $(j,s)$; the
  matrix being Hermitian by (H), this says that its spectrum is non-negative;
\item[(C)] \emph{(completeness)} $\Sigma^{(s)}=\idon{\Qsys_s}$ for every $s$.
\end{itemize}
Each of (H), (P), (C) is an \emph{identity or an order relation between
exactly represented numbers} and is decided as such: (H) and (C) entry by
entry, (P) by exact symmetric elimination on the block with the zero-pivot
row test (a negative pivot fails, and a zero pivot passes only when the rest
of its row is zero), cubic in $d_s$. Together the five conditions are exact
admissibility for $\Smap$ with constant $1$, that is with every completeness
sum equal to the identity, in the sense of
\xref{foundations/def:exact-admissibility}, the declared order of $J$ being
its $\idxm$ order: (T), (H), (P) and (S) are the conditions of that
definition bearing the same names, and (C), which reads the computed sum, is
its completeness identity, because under (S) the computed sum $\Sigma^{(s)}$
is the exact sum $\sum_ja_j^{(s)}$
(\xref{foundations/lem:exact-accumulation}).

\emph{The same conditions on Alice's data, where there is one block and no
sector map.} $A$ carries no sector decomposition, so a family on $A$ is one
array per index and (T) is empty for it. $\{M_x\}_{x\in\Xset}$ is
\emph{exactly admissible} when (S), (H), (P) and (C) hold of it on $A$ with
$\idon{A}$ on the right of (C); the declared marginal $\tau_A$ satisfies its
own three conditions when $\tau_A=\tau_A^{\dagger}$, $\tau_A\succeq0$ and
$\tr\tau_A=1$ with the trace accumulated in the row order \emph{without
rounding}, the trace being a sum, read below as an identity of real numbers
exactly as a completeness sum is.

\emph{What is not on this list.} A strict inequality on a declared scalar,
such as the condition of Definition~\ref{def:weighted} that every
public-choice weight be positive, is a domain condition on a declared real
number, decided by its sign and not by any of the five. A relation symbol, a
declared tag and the type of a right-hand side are likewise decided and never
measured (D(vi)). And nothing in this definition is a \emph{norm}: no
quantity here is compared against a tolerance window, and no
defect, the entrywise size of the difference a condition
names, is computed.

\emph{The general format.} For a declaration in the general format of
Definition~\ref{def:declaration}\,(DX1)--(DX6) the same conditions are read
as follows. \emph{Per measured factor:} for every $f\in\mathbb{F}$ the
family $\{F^{f}_{v_f}\}_{v_f\in\Vset^{f}}$ is exactly admissible for the
sector map of $f$ --- (T) against that factor's map, and (S), (H), (P), (C)
on that factor's blocks in that factor's declared order. No condition is
decided on a product block \eqref{eq:productfamily}, none being a datum; that
the product family then satisfies (H), (P) and (C) on $\Qsys$ as identities
of real numbers is shown in (DX3). An entry's operator $\hat O_k$ is decided
(T) and (H) against the map of its scope. \emph{Alice's data, block by
block:} with a factor tagged $A$, condition (T) is not empty for Alice's
data: $\tau_A$, the operators of her measurement slot and a general
$\Theta_k$ are block lists against that factor's map. In POVM form
$\{M^{(t)}_x\}_x$ satisfies (S), (H), (P) and (C) on every $A_t$ with
$\idon{A_t}$ on the right of (C); $\tau_A$ satisfies (H) and (P) block by
block and $\sum_t\tr\tau^{(t)}_A=1$, the traces accumulated in the row order
and then over the blocks in their declared order, rounding nowhere.
\emph{Ensemble form:} the family $\{N^{(t)}_x\}_x$ satisfies (S), (H) and
(P) on every $A_t$, and in place of (C) the identity
$\Sigma^{(t)}=\tau^{(t)}_A$ of \eqref{eq:ensembleform}, entry by entry,
between the declared-order sum and the declared block of $\tau_A$; and
$\tau^{(t)}_A\ne0$ for every $t$, a comparison of declared entries with the
datum $0$. These are all the conditions of the form: nothing is decided of
the denoted $\{M_x\}$ of \eqref{eq:ensemblepovm}, its being a POVM a theorem
(Theorem~\ref{thm:sifting-and-key}\,(v)). At one Alice sector and one measured factor in
POVM form these readings are the paragraphs above.
\end{digestclaim}

% claim protocol-compiler/def:target
\digestlabel{def:target}{1.5}\pdfbookmark[2]{Definition 1.5 Target conditions of a well-formed declaration}{digest-protocol-compiler/def:target}
\renewcommand*{\RestatedName}{\hyperref{contract.pdf}{}{claim-protocol-compiler/def:target}{Definition~1.5}}%
\begin{digestclaim}[Target conditions of a well-formed declaration]
A declaration \eqref{eq:declaration}, \emph{together with the label-to-index
map $\idxm$ emitted alongside it} (Definition~\ref{def:compile}), is a
\emph{well-formed declaration}
when the following hold, every operator predicate below being decided
\emph{exactly} in the sense of Definition~\ref{def:exact} and every
combinatorial condition exactly as well:
\begin{enumerate}[label=D(\roman*),leftmargin=3.2em]
\item all systems are finite-dimensional, all alphabets are finite and
  non-empty, all declared maps are total, and \emph{the sector map agrees}:
  there is one sector map $\Smap$; every declared $\Qsys$-operator is a
  block list against it, with the same block kind, the same
  identifiers in the same order and the same dimensions; and
  $d_{\Qsys}=\sum_{s\in\SecSet}d_s$. That much is condition~(T) of
  Definition~\ref{def:exact}, stated once for the declaration. \emph{And the
  order declared on the sectors is one order and not two}: the restriction of
  $\idxm$ to $\SecSet$ is the total order $\Smap$ fixes. This last clause is a
  condition on the pair $(\Decl,\idxm)$ and not on a slot of
  \eqref{eq:declaration};
\item the marginal slot holds $\tau_A$ with $\tau_A=\tau_A^{\dagger}$,
  $\tau_A\succeq0$ and $\tr\tau_A=1$, the trace summed in the row order and
  rounding nowhere --- so that $\tau_A\in\Seq{A}$
  holds of the declared numbers, and a slot holding the zero matrix is
  inadmissible, its trace being $0$ --- or the value $\nomarg$;
\item \emph{(typing)} $\{F_v\}$ and every $\hat O_k$ are block lists against
  $\Smap$ in the sense of Definition~\ref{def:declaration}, so that
  $\Pi_sF_v\Pi_{s'}=0$ and $\Pi_s\hat O_k\Pi_{s'}=0$ for $s\ne s'$ hold
  \emph{by construction} and are not conditions on data; and
  \emph{(five predicates)} $M_x\succeq0$ for each $x$ and $\{M_x\}$ complete on
  $A$; $F_v^{(s)}\succeq0$ for each $(v,s)$ and $\{F_v^{(s)}\}_v$ complete on
  $\Qsys_s$ for each $s$; and $\hat O_k^{(s)}$ Hermitian for each $(k,s)$,
  each decided exactly, each positivity predicate being the two decisions
  (H) and (P) of Definition~\ref{def:exact} on the same block, and each
  completeness predicate carrying with it the condition~(S) that its
  declared-order sum round nowhere. This is where Hermiticity of Alice's and
  Bob's emitted families enters the target conditions, and it is the only
  place it does;
\item $r$ maps into $\Zval\sqcup\{\perp\}$;
\item \emph{sifting is announcement-determined}: there is
  $\varsigma:\Kann\to\{\mathrm{pass},\mathrm{fail}\}$ with
  $r(x,v)=\perp\iff\varsigma(a(x,v))=\mathrm{fail}$;
\item in trusted-source mode $\bowtie_k\in\{\preceq,=,\succeq\}$ and
  each entry carries one of the two declared forms of
  Definition~\ref{def:homog}, homogeneous with $c_k\in\mathbb{R}$ or
  general with $\Theta_k$ Hermitian on $A$; in
  untrusted-source mode $\bowtie_k\in\{\le,=,\ge\}$ with
  $\Theta_k=c_k\in\mathbb{R}$. The relation symbol and the type of the
  right-hand side, which is what fixes the tag (Definition~\ref{def:homog}),
  are combinatorial and are read exactly; the Hermiticity of a general
  $\Theta_k$ is the exact predicate (H) on one matrix;
\item the feasible set of Definition~\ref{def:feasible} is non-empty,
  $\Feas\ne\varnothing$;
\item \emph{party computability}: each party announces a function of its own
  readout ($a_A$ is defined on $\Oset{\mathsf A}$ and $a_B$ on
  $\Oset{\mathsf B}$, jointly \eqref{eq:jointann}) and the declared holder
  computes the raw key from its own readout and the announcements: there is a
  total $r_{\hold}:\Oset{\hold}\times\Kann\to\Zset$ with
  \begin{equation*}\tag{\getrefnumber{C-eq:partycomp}}\digesteqlabel{eq:partycomp}
    r(x,v)=r_{\hold}\big(\rd{\hold}(x,v),\,a(x,v)\big)
    \qquad\text{for all }(x,v)\in\Xset\times\Vset .
  \end{equation*}
\end{enumerate}
The
\emph{default party partition}
is $\rd{\mathsf A}(x,v)=x$, $\rd{\mathsf B}(x,v)=v$, with $A$ labelled
$\mathsf A$ and every $\Qsys_s$ labelled $\mathsf B$.

\emph{The general format.} For a declaration in the general format of
Definition~\ref{def:declaration}\,(DX1)--(DX6) the conditions read as
follows, and in the basic format --- one Alice sector, one measured
factor, POVM form, two messages neither reading the other and no party-local
entry --- they are the eight above.
D(i): the sector map is one factorised map (at most one factor tagged $A$,
and then the first, at least one tagged $\Qsys$, one block kind), every
declared family on a measured factor and every entry's operator is a block
list against the map of its factor, Alice's data against the factor tagged
$A$, $d_{\Qsys}=\prod_f\sum_{s_f}d_{f,s_f}$, and the restriction of $\idxm$
to each factor's sectors is the order that factor's map fixes, the product
order being derived and not declared. D(ii): $\tau_A$ is a block list along
$\mathbb{T}$ with every block Hermitian and positive semidefinite and total
trace $1$, summed without rounding; in ensemble form additionally every
block is non-zero. D(iii): the typing and the predicates of
Definition~\ref{def:exact} in their general-format reading, per measured
factor and per Alice sector, the ensemble-form conditions in place of those
on $\{M_x\}$ where that form is declared; ensemble form is refused in
untrusted-source mode. D(vi): an entry carries its scope, which names a
measured factor of the map, and one of the forms of
Definition~\ref{def:homog}, the party-local form included, in either
mode, with a real $c_k$; a general $\Theta_k$ is a block list along
$\mathbb{T}$, Hermitian block by block. \emph{And the relation symbol is
read by form and not by mode alone}: an \emph{untagged} entry carries
$\bowtie_k\in\{\preceq,=,\succeq\}$ in trusted-source mode and
$\bowtie_k\in\{\le,=,\ge\}$ in untrusted-source mode, exactly as stated for
D(vi), while an entry declared party-local carries
$\bowtie_k\in\{\preceq,=,\succeq\}$ \emph{in either mode}, its reading being
the operator relation of the last line of \eqref{eq:feas-general}
(Definition~\ref{def:homog}, the relation-symbol clause of the party-local
form); where every row of that entry is $1\times1$ (every Alice block, and
every sector block of every measured factor other than the entry's scope, of
dimension one, $d_A=1$ at one measured factor included) that relation is
between $1\times1$ matrices and the three operator symbols are the three
scalar symbols, which is the degenerate case. Where no entry is party-local,
as in every declaration of the basic format, this clause is D(vi) as
stated. D(viii): the announcement slot is a transcript protocol (DX5),
every table $a_m$ being total on
$\Oset{\pi_m}\times\mathcal{C}_1\times\dots\times\mathcal{C}_{m-1}$, so that
a message reads its sender's readout and earlier messages \emph{by its
signature}; the joint function $a$ is the composition \eqref{eq:transcript},
and \eqref{eq:partycomp} holds with it. D(iv), D(v) and D(vii) keep their
wording, with $a$ the derived function and $\Feas$ the set of
\eqref{eq:feas-general}.
\end{digestclaim}

\section*{2\quad Feasible states and licensed constraints}\pdfbookmark[1]{2 Feasible states and licensed constraints}{digest-group-protocol-compiler/def:feasible}

% claim protocol-compiler/def:feasible
\digestlabel{def:feasible}{2.1}\pdfbookmark[2]{Definition 2.1 Feasible set}{digest-protocol-compiler/def:feasible}
\renewcommand*{\RestatedName}{\hyperref{contract.pdf}{}{claim-protocol-compiler/def:feasible}{Definition~2.1}}%
\begin{digestclaim}[Feasible set]
For a declaration $\Decl$, the
\emph{feasible set}
$\Feas$ is the set of $\omega$ presented as a
sector-block list on $A\Qsys$ --- $\omega=\bigoplus_s\omega^{(s)}$ with $\omega^{(s)}\in\Lin{A\Qsys_s}$,
so that $(\idon{A}\otimes\Pi_s)\omega(\idon{A}\otimes\Pi_{s'})=0$ for
$s\ne s'$ holds by the presentation and is not a condition --- such that
\begin{equation*}\tag{\getrefnumber{C-eq:feas}}\digesteqlabel{eq:feas}
  \omega^{(s)}\succeq0\ \ \forall s,\qquad \tr\omega=1,\qquad
  \omega\models\Decl,\qquad \omega\models\CList ,
\end{equation*}
each read \emph{exactly}, with $\tr\omega=\sum_s\tr\omega^{(s)}$ and with
$\omega^{(s)}\succeq0$ carrying, as in Definition~\ref{def:exact},
Hermiticity of the block. Here $\omega\models\Decl$ is the marginal condition,
$\omega_A=\tau_A$ in trusted-source mode, with
$\omega_A=\sum_s\tr_{\Qsys_s}\omega^{(s)}$, and no condition beyond
$\tr\omega=1$ in untrusted-source mode; and $\omega\models\CList$ reads
\begin{equation*}\tag{\getrefnumber{C-eq:feas-modes}}\digesteqlabel{eq:feas-modes}
\begin{aligned}
  \text{trusted:}\quad&
  \tr_{\Qsys}\!\big[(\idon{A}\otimes\hat O_k)\,\omega\big]
  =\sum_s\tr_{\Qsys_s}\!\big[(\idon{A}\otimes\hat O^{(s)}_k)\,\omega^{(s)}\big]
  \ \bowtie_k\ \Theta_k
  \quad\forall k;\\
  \text{untrusted:}\quad&
  \tr\!\big[(\idon{A}\otimes\hat O_k)\,\omega\big]
  =\sum_s\tr\!\big[(\idon{A}\otimes\hat O^{(s)}_k)\,\omega^{(s)}\big]
  \ \bowtie_k\ c_k
  \quad\forall k ,
\end{aligned}
\end{equation*}
where in the trusted line $\Theta_k$ is the \emph{constructed} right-hand side
$c_k\tau_A$ for an entry declared homogeneous and the \emph{declared}
Hermitian matrix for one declared general
(Definition~\ref{def:homog}). The set carries no parameter, admits no
tolerance window and relaxes no conjunct; in particular positivity is never
relaxed.

\emph{The general format.} For a declaration in the general format of
Definition~\ref{def:declaration}\,(DX1)--(DX6), $\Feas$ is the set of
$\omega$ presented as a list over the
joint sectors,
$\omega=\bigoplus_{(t,s)}\omega^{(t,s)}$ with
$\omega^{(t,s)}\in\Lin{A_t\Qsys_s}$ and every block positive semidefinite, so
that block diagonality along $\mathbb{T}$ and along $\SecSet$ both hold by
the presentation, subject to the same conditions $\tr\omega=1$,
$\omega\models\Decl$ and $\omega\models\CList$. Write
$\hat O^{\uparrow(s)}_k\coloneqq\idon{}\otimes\hat O^{(s_{f_k})}_k\otimes\idon{}$
for the $s$-block of the lifted operator of an entry scoped to $f_k$
(DX6), and for a party-local entry scoped to $f$ split
$s=(s_{f^c},s_f)$, with $\Qsys_{f^c,s_{f^c}}$ the product of the other
factors' blocks. The conditions then read, block by block,
\begin{equation*}\tag{\getrefnumber{C-eq:feas-general}}\digesteqlabel{eq:feas-general}
\begin{aligned}
  &\textstyle\sum_{(t,s)}\tr\omega^{(t,s)}=1
  &&\text{(normalization)};\\
  &\textstyle\sum_{s}\tr_{\Qsys_s}\omega^{(t,s)}=\tau^{(t)}_A
  \quad\forall t
  &&\text{(marginal, trusted-source mode)};\\
  &\textstyle\sum_{s}\tr_{\Qsys_s}\!\big[(\idon{A_t}\otimes
  \hat O^{\uparrow(s)}_k)\,\omega^{(t,s)}\big]\bowtie_k\Theta^{(t)}_k
  \quad\forall t
  &&\text{(trusted-source mode)};\\
  &\textstyle\sum_{(t,s)}\tr\!\big[(\idon{A_t}\otimes
  \hat O^{\uparrow(s)}_k)\,\omega^{(t,s)}\big]\bowtie_kc_k
  &&\text{(untrusted-source mode)};\\
  &\textstyle\sum_{s_f\in\SecSet^{f}}\tr_{\Qsys_{f,s_f}}\!\Big[\big(\idon{}
  \otimes(\hat O^{(s_f)}_k-c_k\idon{})\big)\,
  \omega^{(t,(s_{f^c},s_f))}\Big]\bowtie_k0
  \quad\forall(t,s_{f^c})
  &&\text{(party-local, either mode)},
\end{aligned}
\end{equation*}
the third line with $\Theta^{(t)}_k=c_k\tau^{(t)}_A$ for an entry declared
homogeneous, the fourth for a scalar entry, and the last as an operator on
$\Hs{A_t}\otimes\Hs{\Qsys_{f^c,s_{f^c}}}$. Each
line is the corresponding operator condition on $A\Qsys$: the marginal
$\omega_A$, the entry value $\tr_{\Qsys}[(\idon{A}\otimes\hat O_k)\omega]$,
$\tau_A$ and every $\Theta_k$ are block diagonal along $\mathbb{T}$, the
party-local value
$\tr_{\Qsys_f}[(\idon{}\otimes(\hat O_k-c_k\idon{\Qsys_f}))\,\omega]$ is block
diagonal along the pairs $(t,s_{f^c})$, and an equality or a semidefinite
inequality between operators block diagonal for one decomposition holds
exactly when it holds block by block. With $|\mathbb{T}|=1$, one
measured factor and no party-local entry, \eqref{eq:feas-general} is
\eqref{eq:feas} with \eqref{eq:feas-modes}, line for line. With one measured
factor the party-local line reads
$\sum_s\tr_{\Qsys_s}[(\idon{A_t}\otimes(\hat O^{(s)}_k-c_k\idon{}))
\omega^{(t,s)}]\bowtie_k0$ on $\Hs{A_t}$, that is
$\tr_{\Qsys}[(\idon{A}\otimes\hat O_k)\omega]\bowtie_kc_k\,\omega_A$. The
set is convex and compact: every line is an affine equality or an affine
semidefinite or scalar inequality, and $\omega\succeq0$ with $\tr\omega=1$.
\end{digestclaim}

% claim protocol-compiler/def:homog
\digestlabel{def:homog}{2.5}\pdfbookmark[2]{Definition 2.5 Licensed constraint list}{digest-protocol-compiler/def:homog}
\renewcommand*{\RestatedName}{\hyperref{contract.pdf}{}{claim-protocol-compiler/def:homog}{Definition~2.5}}%
\begin{digestclaim}[Licensed constraint list]
Each entry of $\CList$ carries its right-hand side in a declared form, and
the form is part of the datum rather than a property measured of it
(Definition~\ref{def:declaration}). There are two \emph{untagged} forms,
called so because their tag is the type of their datum and needs no field of
its own: a
\emph{homogeneous}
entry declares a real scalar $c_k$, a
\emph{general}
entry declares a Hermitian $\Theta_k$ on $A$; in untrusted-source mode only
the scalar form exists. For these two forms
the tag \emph{is} the type of the right-hand side an entry carries, and it is
read off that type: an entry whose right-hand-side slot holds a real scalar
is homogeneous, and one whose slot holds a Hermitian matrix on $A$ is
general. \emph{That
reading is a type decision}: it inspects which of two objects a slot holds,
compares no numbers, reads no matrix entry and measures nothing. The third
form, party-local, is declared by its tag and is not read off the type of the
datum. $\CList$ is
\emph{licensed for the declared source mode}
when every entry is in the form that mode admits:
\begin{enumerate}[label=(\alph*),leftmargin=2.2em]
\item \emph{trusted-source mode:} every entry is declared
  homogeneous, carrying $c_k\in\mathbb{R}$; the right-hand side it denotes is
  then \emph{constructed} as $\Theta_k\coloneqq c_k\tau_A$ from the declared
  $c_k$ and the declared $\tau_A$, and the entry of \eqref{eq:feas-modes}
  reads $\tr_{\Qsys}[(\idon{A}\otimes\hat O_k)\,\omega]\bowtie_kc_k\tau_A$;
\item \emph{untrusted-source mode:} every right-hand side is a real number
  $\Theta_k=c_k$ and the entry reads
  $\tr[(\idon{A}\otimes\hat O_k)\,\omega]\bowtie_kc_k$.
\end{enumerate}
Licensing is therefore a \emph{type} check on the tag, and it compares no
numbers: what is inspected is which of the two forms an entry declares; no
operator identity is tested against any number, and nothing here is decided
at a window. \emph{Only in cases (a) and
(b) are the entries written $(\hat O_k,\bowtie_k,c_k)$}; a
general entry is written $(\hat O_k,\bowtie_k,\Theta_k)$ and is the form the
notation table records under $\Theta_k$. A general entry is admissible
input, every structural claim of this contract being proved for it, but it is
not licensed, so a list containing one carries no key-length claim
(Theorem~\ref{thm:gates}).

\emph{The party-local form, a third licensed form.} In the general format of
Definition~\ref{def:declaration}\,(DX6) an entry may be declared
\emph{party-local}:
it carries its scope $f$, a block list $\hat O_k$ against the
sector map of the measured factor $f$ and a real $c_k$, it is
written
$(\hat O_k,\bowtie_k,c_k;f)$, and its state-level reading, in either source
mode, is
\begin{equation*}\tag{\getrefnumber{C-eq:partylocal-entry}}\digesteqlabel{eq:partylocal-entry}
  \tr_{\Qsys_f}\!\Big[\big(\idon{}\otimes
  (\hat O_k-c_k\,\idon{\Qsys_f})\big)\,\omega\Big]\;\bowtie_k\;0
  \qquad\text{as an operator on the systems of }\omega
  \text{ other than }\Qsys_f ,
\end{equation*}
those systems being $A$ and the measured factors $\Qsys_{f'}$, $f'\ne f$;
equivalently $\tr_{\Qsys_f}[(\idon{}\otimes\hat O_k)\,\omega]\bowtie_k
c_k\tr_{\Qsys_f}\omega$, the entry being homogeneous with respect to the
marginal that is actually there. Block by block it is the last line of
\eqref{eq:feas-general}. The list is then \emph{licensed} when every entry
is either in the form the source mode admits, (a) or (b), read for its
lifted operator, or is declared party-local. This is still a type
check and compares no numbers: what is inspected is the declared tag. A
party-local entry's right-hand-side slot holds a real scalar, as a
homogeneous entry's does, so the type of the datum does not separate the
two: an entry is party-local exactly when it carries that tag.

\emph{The soundness condition of a party-local entry.} A party-local entry
$(\hat O_k,\bowtie_k,c_k;f)$ meets the
\emph{soundness condition of a party-local entry}
when every attack reaches $\Qsys_f$ through a channel $\Lambda_f$ acting on
that factor alone which carries the spectral certificate
$\Lambda_f^{\dagger}(\hat O_k)\bowtie_kc_k\idon{}$, that is, the adjoint of
the channel applied to the entry's operator stands in the entry's relation to
$c_k$ times the identity.

\emph{What the tag changes, case by case.} With one measured factor and
trusted-source mode, \eqref{eq:partylocal-entry} reads
$\tr_{\Qsys}[(\idon{A}\otimes\hat O_k)\,\omega]\bowtie_kc_k\,\omega_A$, which
under the marginal condition $\omega_A=\tau_A$ is the form (a) of the same
$(\hat O_k,\bowtie_k,c_k)$. With one measured factor and $d_A=1$, in either
mode, it is the comparison of real numbers
$\tr[\hat O_k\omega]\bowtie_kc_k\tr\omega=c_k$, the form (b). \emph{With one
measured factor, untrusted-source mode and $d_A>1$ it is neither}: it is an
operator inequality on $A$ against the marginal $\omega_A$ that is actually
there, and the scalar entry (b) of the same $(\hat O_k,\bowtie_k,c_k)$ is its
trace, implied by it and not equivalent to it. With more than one measured
factor it implies, in either mode, the untagged reading of the same entry,
which is obtained from it by the partial trace over the other measured
factors (and over $A$ in untrusted-source mode), and it is in general not
implied by it. So the tag never weakens an entry, and of the four one-factor
cases, by source mode and by whether $d_A=1$, it changes something only in
untrusted-source mode at $d_A>1$.

\emph{Which relation symbols a tagged entry carries, in either mode.} An entry
declared party-local carries $\bowtie_k\in\{\preceq,=,\succeq\}$ \emph{in
either source mode}, because \eqref{eq:partylocal-entry} is a relation
between \emph{operators} in either mode: its left-hand side is an operator on
the systems of $\omega$ other than $\Qsys_{f_k}$, which are $A$ and the other
measured factors. The scalar case is the degenerate one: where
$d_{A_t}\prod_{f\ne f_k}d_{f,s_f}=1$ for every row --- every Alice block, and
every sector block of every measured factor other than $f_k$, of dimension
one, which with one measured factor means every Alice block of dimension one,
whether $d_A=1$ or $d_A=|\mathbb{T}|>1$ --- each row is a comparison of
$1\times1$ matrices, that is of real numbers, and the three operator symbols
are then the three scalar symbols $\le$, $=$, $\ge$. With one measured factor
and $d_A=1$ there is one row, and the reading is the form (b) of the same
$(\hat O_k,\bowtie_k,c_k)$, as the case-by-case paragraph states; otherwise
there is one such comparison per row, and the reading relates to the untagged
forms as that paragraph states for its case. \emph{Nothing about the
untagged forms changes}: an untagged entry carries the symbols its mode
admits, and the operator symbols beside a tag are a property of the third
form and of no other. No entry of a declaration in the basic format is
tagged.
\end{digestclaim}

\section*{3\quad Specification, compilation map and typing}\pdfbookmark[1]{3 Specification, compilation map and typing}{digest-group-protocol-compiler/def:spec}

% claim protocol-compiler/def:spec
\digestlabel{def:spec}{3.1}\pdfbookmark[2]{Definition 3.1 Protocol specification}{digest-protocol-compiler/def:spec}
\renewcommand*{\RestatedName}{\hyperref{contract.pdf}{}{claim-protocol-compiler/def:spec}{Definition~3.1}}%
\begin{digestclaim}[Protocol specification]
A \emph{protocol specification} is a tuple
\begin{equation*}
  \Spec=\big(\Spec_{\mathrm{src}},\ \Spec_{\mathrm{rec}},\ \pcChain,\
  \Spec_{\mathrm{pub}},\ \Spec_{\mathrm{con}}\big)
\end{equation*}
whose parts are:

\emph{(a) Source part $\Spec_{\mathrm{src}}$.} A source specification
in the sense of \xref{transmitter/def:source-spec}, presented to this module as
the \emph{transmitter artefact} in its declared mode: a finite non-empty
alphabet $\Xset$ with its declared order, a finite-dimensional system
$A$ with declared computational basis, a POVM $\{M_x\}_{x\in\Xset}$ on $A$
--- itself possibly presented as a choice-weighted family
(Definition~\ref{def:weighted}) --- and the specified signal ensemble
$\{p_x,\rho_x\}_{x\in\Xset}$ on the transmitted system $\Asig$
(\xref{transmitter/ass:ens}). The mode is read off the artefact's
marginal slot and is never a second datum. In
trusted-source mode the slot holds $\tau_A\in\Seq{A}$ --- at $d_A=1$
the slot the transmitter filled with $[1]$ included --- and at $d_A\ge2$ the
artefact additionally carries a source-replacement certificate
$(\tau_A,\Psi,\{p_x,\rvirt_x\}_{x\in\Xset})$ on $A\Avirt$ in the sense of
\xref{transmitter/def:srcert}, written against the virtual ensemble
the source segment \emph{starts from}, which lives on the
virtual transmitted system $\Avirt$, the source descriptor of
$\Xi_1$ (\xref{transmitter/def:source-segment}\,(S1), (S3)). When
$L_{\mathrm s}=0$ that system and that ensemble are the specified ones:
$\Avirt=\Asig$ and $\rvirt_x=\rho_x$. In untrusted-source mode the slot
holds $\nomarg$ and there is no certificate.

\emph{(b) Receiver part $\Spec_{\mathrm{rec}}$.} A finite-dimensional physical
receiver system with a POVM $\{\Fph_w\}_{w\in\Vph}$, possibly presented as a
choice-weighted family; and, optionally, a public choice whose weight
depends on the private setting of the measuring apparatus --- a
stochastic matrix $P(f\,|\,t)$ with the setting map
$t:\Vset\to\mathcal{T}$ it is read through --- to be applied by
this module \emph{after} the reduction chain, in the
setting-dependent form of Definition~\ref{def:weighted}.

\emph{(c) Reduction chain $\pcChain=(\Chsrc,\Chrec)$.} A reduction chain in
the sense of \xref{reductions/def:chain}, presented to this module as a
record: an ordered \emph{pair} of ordered lists of typed links, the
\emph{source segment} $\Chsrc=(\Xi_1,\dots,\Xi_{L_{\mathrm s}})$ with
$L_{\mathrm s}\ge0$, which is a source segment in the sense of
\xref{transmitter/def:source-segment} --- its links act on the transmitted
system $\Asig$ and replace Alice's signal ensemble --- and the \emph{receiver
segment} $\Chrec=(\Lambda_1,\dots,\Lambda_L)$ with $L\ge1$, whose links act on
the receiver system and replace Bob's measurement. What this module reads of a
link, stated in its own words, is this: each link of either segment carries a
type tag, source and target system descriptors, a status tag, its channel as
data where the link type has one, the constraint entries it transports, and a
licence slot. A receiver-stage link in the sense of
\xref{reductions/def:link} carries more --- its reachable-set statement and
its carried operator pair --- and this module reads none of it. The last link
of $\Chrec$ has target descriptor $(\Qsys,\SecSet,(d_s)_s)$ with the
sector projectors $\Pi_s$, and target POVM $\{F_v\}_{v\in\Vset}$. The two
segments are separate lists because they have different types; the pair is
composed source segment first. The interface conditions on a chain that
this module decides on the record are those of
Definition~\ref{def:chaincond}.

\emph{(d) Public part $\Spec_{\mathrm{pub}}$.} A party partition
$\Party$ (readouts $\rd{\mathsf A},\rd{\mathsf B}$ onto $\Oset{\mathsf A}$,
$\Oset{\mathsf B}$, and the party labels), announcement maps
$a_A:\Oset{\mathsf A}\to\CAset$ and $a_B:\Oset{\mathsf B}\to\CBset$, a sifting
predicate $\varsigma:\Kann\to\{\mathrm{pass},\mathrm{fail}\}$, a key rule
$r_{\hold}^{0}:\Oset{\hold}\times\Kann\to\Zval$, a holder
$\hold\in\{\mathsf A,\mathsf B\}$, a finite acceptance alphabet $\Chatset$ and
a map $\phi:\Kann\to\Chatset$. All maps are total.

\emph{(e) Constraint part $\Spec_{\mathrm{con}}$.} A finite, possibly empty,
list of entries. Each $\hat O_k$ is presented as a block list on
$\Qsys$ together with the sector map it is written against, which is
carried with it as for every other supplied $\Qsys$-operator family, and is
Hermitian block by block. Each entry carries a right-hand-side datum of one of
two types, and the type fixes the entry's tag: homogeneous with a real
$c_k$, or general with a Hermitian $\Theta_k$ on $A$, in trusted-source
mode; a real $c_k$ in untrusted-source mode (Definition~\ref{def:homog}). The
constraint part also carries the optional attack-set licence record and an
optional feasibility witness $\omega^{\star}$, likewise a block list with its
own map. The map each $\hat O_k$ carries is compared with the emitted $\Smap$
by check V1, an entry declared here being on the same footing as one
transported by a link. The witness's map is not compared there: $\omega^{\star}$
is a block list on $A\Qsys$ and not a $\Qsys$-operator, it is no slot of
\eqref{eq:declaration}, and branch~(c) of check V7 is where it is required to
be a block list against $\Smap$ (Definition~\ref{def:validator}).

\emph{The general format.} A specification to be compiled into the general
format of Definition~\ref{def:declaration}\,(DX1)--(DX6) has the same five
parts, read as follows. At one Alice sector, one measured factor,
POVM form, two messages neither of which reads the other and no
party-local entry these readings are (a)--(e) above, part for part on every
component a derivation reads: at those values every such object named below
is the object of the same name there and nothing else is presented. On the
party labels of the tensor factors of each $\Hs{\Qsys_s}$, which no
derivation of this module reads and which the consuming engine's laboratory
premise reads as naming the parties' laboratories, (DX1) restricts the party
partition of (d). The party set is unchanged: it is
$\{\mathsf A,\mathsf B\}$, with the two readout maps of $\Party$. What
is new is that several measured factors may carry a party label and that the
factor tagged $A$ carries one too (DX1).

\emph{(a$'$) Source part.} The artefact's marginal slot may hold a block list
$(\tau^{(t)}_A)_{t\in\mathbb{T}}$ along the Alice sectors the artefact
declares, and Alice's measurement arrives in one of the two forms of (DX4):
the POVM slot, or, in trusted-source mode only, the ensemble slot holding the
block lists $\{N_x\}_{x\in\Xset}$ of \eqref{eq:ensembleform}. The
source-replacement certificate may accordingly be declared in
pure form, in block form, or as an
ensemble presentation, all three in the sense of
\xref{transmitter/def:srcert}. Which form a certificate is in is read off the
artefact and is never a second datum, exactly as the mode is. That a form may
be \emph{presented} is the transmitter's rule and is stated here. Which forms
this module \emph{admits}, and what gate (G2) decides of each form, is
Definition~\ref{def:eligible}'s and is stated once there.

\emph{(b$'$) Receiver part.} $\Spec_{\mathrm{rec}}$ is an ordered list, one
entry per measured factor $f\in\mathbb{F}$ and in the order of $\mathbb{F}$: a
party label; a finite-dimensional physical receiver system with a POVM
$\{F^{\mathrm{ph},f}_{w_f}\}_{w_f\in\mathcal{V}^{\mathrm{ph},f}}$, possibly
presented as a choice-weighted family; and, optionally, a
setting-dependent public choice at that factor --- a stochastic matrix
$P_f(\,\cdot\mid\cdot\,)$ with the setting map on that factor's alphabet it is
read through --- to be applied after that factor's segment, in the
setting-dependent form of Definition~\ref{def:weighted}. No entry of the list
carries an operator on another factor's system, and no product of physical
families is presented: the product is the rule \eqref{eq:productfamily},
formed by no one.

\emph{(c$'$) Reduction chain.} $\pcChain=(\Chsrc,\Chrec)$ with $\Chsrc$ as
in (c) and $\Chrec=(\pcChain^{f}_{\mathrm{rec}})_{f\in\mathbb{F}}$ one
receiver segment per measured factor, each
$\pcChain^{f}_{\mathrm{rec}}=(\Lambda^{f}_1,\dots,\Lambda^{f}_{L_f})$,
$L_f\ge1$, a list of receiver-stage links whose source descriptor is factor
$f$'s physical receiver system and whose last target descriptor is
$(\Qsys_f,\SecSet^{f},(d_{f,s_f})_{s_f})$ with target POVM
$\{F^{f}_{v_f}\}_{v_f\in\Vset^{f}}$. What this module reads of a link is what
(c) lists and nothing more, of each segment separately: the segments are
separate lists because their links act on different systems, and no link of
one is composable with a link of another. Where the factors' data are the
ordered list of the factors' artefacts that the reductions module emits,
factor $f$'s receiver segment and target family are that factor's
own, and the product rule of \xref{reductions/thm:emitted-product} is
emitted in place of a joint family; that list
is what the factor list of (DX1) is built from, in its order. The interface
conditions on a chain that this module decides on the record are those of
Definition~\ref{def:chaincond}, per factor for the receiver segments.

\emph{(d$'$) Public part.} In place of $a_A$ and $a_B$ the public part
declares the transcript of (DX5): the number $L\ge1$ of messages and,
for each $m$, its sender $\pi_m$, its finite alphabet $\mathcal{C}_m$ and
its total table
$a_m:\Oset{\pi_m}\times\mathcal{C}_1\times\dots\times\mathcal{C}_{m-1}
\to\mathcal{C}_m$. The sifting predicate $\varsigma$, the key rule
$r^{0}_{\hold}$, the holder $\hold$, the acceptance alphabet $\Chatset$ and
$\phi$ are declared on the derived $\Kann=\prod_{m}\mathcal{C}_m$ of
\eqref{eq:transcript}, which is a set determined by the tables and declared by
no one. All tables are total. The party partition $\Party$ is as in (d), with
a party label per measured factor and, where there is one, for the factor
tagged $A$.

\emph{(e$'$) Constraint part.} Each entry presents, in addition to what (e)
lists, its scope: the label of one measured factor of the chain; and
$\hat O_k$ is a block list on $\Qsys_{f_k}$ together with the sector map of
that factor, Hermitian block by block. Its right-hand-side datum is of one of
the three forms of Definition~\ref{def:homog}: a real $c_k$ tagged
homogeneous, a Hermitian $\Theta_k$ tagged general and then itself a block
list along $\mathbb{T}$, or a real $c_k$ tagged party-local. The
arriving tag is compared with the type of that datum at check V6, which also
decides that the scope names a measured factor of the emitted map. An optional
witness $\omega^{\star}$ is a block list over the joint sectors with
its own map, compared at branch~(c) of V7 and nowhere else, for the reason
(e) gives.
\end{digestclaim}

% claim protocol-compiler/def:well-formed-spec
\digestlabel{def:well-formed-spec}{3.2}\pdfbookmark[2]{Definition 3.2 Well-formed specification}{digest-protocol-compiler/def:well-formed-spec}
\renewcommand*{\RestatedName}{\hyperref{contract.pdf}{}{claim-protocol-compiler/def:well-formed-spec}{Definition~3.2}}%
\begin{digestclaim}[Well-formed specification]
A protocol specification $\Spec$ (Definition~\ref{def:spec}) is a
\emph{well-formed specification}
when its declared data are as declared:
(i) all alphabets are finite and non-empty;
(ii) all systems are finite-dimensional with the declared dimensions;
(iii) all tabulated maps are total on their stated domains;
(iv) all declared POVM families are exactly admissible: every family
other than the physical receiver family $\{\Fph_w\}_{w\in\Vph}$ as presented,
block by block in the sense of Definition~\ref{def:exact}; and the physical
receiver family as presented (in the general format, each
measured factor's, on that factor's physical receiver system), on its
physical receiver space read as one block, condition~(T) being empty for it
as Definition~\ref{def:exact} reads Alice's data: where it is presented
directly, (S), (H), (P) and (C) of that definition hold of it in the
declared order of its alphabet; where it is presented as a
choice-weighted family (Definition~\ref{def:weighted}), its
conditional families are exactly admissible, block by block against the map
they carry, and condition~(S) holds of its own declared-order sum. A
choice-weighted family whose conditional families carry a sector map is a
block list against that map, and reading it as that block list or as one
block decides the same (S), (H), (P) and (C), its off-block entries being
exact zeros;
(v) all public-choice weights are strictly positive --- a domain condition on
a declared scalar (Definition~\ref{def:weighted}) --- and sum to one;
(vi) every setting-dependent weighting, at whichever measured system or
measured factor it is declared, is a stochastic matrix, its
entries non-negative and its rows summing to one exactly (the
setting-dependent form of Definition~\ref{def:weighted});
(vii) every constraint operator is a block list against the upstream
module's sector map and Hermitian block by block, and every entry
carries its declared right-hand-side tag.
\end{digestclaim}

% claim protocol-compiler/def:compile
\digestlabel{def:compile}{3.3}\pdfbookmark[2]{Definition 3.3 The compilation map}{digest-protocol-compiler/def:compile}
\renewcommand*{\RestatedName}{\hyperref{contract.pdf}{}{claim-protocol-compiler/def:compile}{Definition~3.3}}%
\begin{digestclaim}[The compilation map]
The \emph{compilation map} is the map $\Spec\mapsto\Decl$ that fills
\eqref{eq:declaration} as follows.
\begin{itemize}[leftmargin=1.6em]
\item \emph{Alice's system and measurement.} $A$, $d_A$ and $\{M_x\}$ from
  $\Spec_{\mathrm{src}}$, weighted by Definition~\ref{def:weighted} where a
  choice family is presented.
\item \emph{Alice's marginal.} $\tau_A$ is what the transmitter artefact's
  marginal slot holds: in trusted-source mode the matrix in the slot --- at
  $d_A=1$ the value $[1]$ the transmitter filled the slot with, the unique
  element of $\Seq{A}$ --- and
  $\tau_A=\nomarg$ where the slot is empty. The mode is read off that slot and
  off nothing else. A certificate the artefact carries is consumed by gate
  (G2) of Definition~\ref{def:eligible}.
\item \emph{The measured system and its family.} $(\Qsys,\SecSet,(d_s)_s)$
  and $\{F_v\}$ from the last link of $\Chrec$. Where $\Spec_{\mathrm{rec}}$
  declares a setting-dependent public choice, the family is weighted
  after the chain by the setting-dependent form of
  Definition~\ref{def:weighted}: the emitted family is then
  $N_{v,f}=P(f\,|\,t(v))\,F_v$ over the support of $P$, the emitted outcome
  alphabet is that support, and the setting map $t$ through which $P$ is read
  is the class map of the joint alphabet of the chain's producer where
  the chain was emitted by the reductions module
  (\xref{reductions/def:joint-family}\,(a)), and is declared with the family
  otherwise.
\item \emph{The public maps.} $\Party$, $a_A$, $a_B$, $\phi$, $\hold$
  verbatim from $\Spec_{\mathrm{pub}}$, declared on the emitted alphabets.
\item \emph{The constraint list.} $\CList$ is the concatenation of
  $\Spec_{\mathrm{con}}$'s entries with the entries transported by the links
  of $\Chsrc$ and then those of $\Chrec$, each segment in its own order, each
  entry with the datum its right-hand-side slot holds and with the tag that
  datum's type determines (Definition~\ref{def:homog}).
\item \emph{The raw-key map.}
\begin{equation*}\tag{\getrefnumber{C-eq:rdef}}\digesteqlabel{eq:rdef}
  r(x,v)\coloneqq
  \begin{cases}
    r^{0}_{\hold}\big(\rd{\hold}(x,v),\,a(x,v)\big),
      & \varsigma(a(x,v))=\mathrm{pass},\\[0.3ex]
    \perp, & \varsigma(a(x,v))=\mathrm{fail},
  \end{cases}
  \qquad\text{with }d_Z\coloneqq|\Zval| .
\end{equation*}
\end{itemize}

\emph{The index map.} Alongside $\Decl$ the compiler emits the
\emph{index map}
$\idxm$, the label-to-index map: a choice, for each of $\Xset$, $\Vset$,
$\SecSet$, $\Oset{\mathsf A}$, $\Oset{\mathsf B}$, $\CAset$, $\CBset$,
$\Chatset$ and $\Zset$, of a bijection onto an initial segment of
$\mathbb{Z}_{\ge0}$; every emitted list, matrix and function table is indexed
through $\idxm$, and every semantic label is carried as metadata only. Its
component on $\SecSet$ is the total order the sector map $\Smap$ already
fixes, and not a second choice: one order on the sectors, required by the last
clause of D(i) (Definition~\ref{def:target}) and decided by V1. And $\idxm$ is
what fixes the declared summation order of Definition~\ref{def:exact} on every
alphabet it indexes; it indexes no basis of a matrix, a trace being
accumulated in the row order of the emitted array.

\emph{The binding rule.} On the remaining alphabets the choice is the
compiler's, under one binding rule. Where a supplied family arrives with a
declared order on the alphabet that indexes it --- as the target POVM of the
last link of $\Chrec$ does, its declarer having established its own
completeness sum in that order --- the compiler takes \emph{that} order as
$\idxm$'s component there. So the order in which condition~(S) of
Definition~\ref{def:exact} is decided on a supplied family here is the order
that family was declared in, and no family is refused here for an ordering
its declarer never chose. Where no order arrives with an alphabet the choice is free, and (S) is
decided in the order the compiler chose and emits. This rule prescribes
nothing to the declarer of a family: it constrains only which bijection this
module picks.

Where a choice-weighted presentation is resolved, this same rule settles the
order the decision on each conditional family is read in, and adds nothing to
it. The family the weighting produces is indexed by one alphabet --- $\Xset$
where the presentation is Alice's, the physical receiver alphabet where it is
a measuring party's, which in Band~A is the emitted $\Vset$ --- and where a declared order
arrives with that alphabet the rule above takes \emph{it}, here as everywhere:
$\Xset$ arrives ordered with the source part
(Definition~\ref{def:spec}\,(a)), and the physical alphabet with the order the
fibre sum of Definition~\ref{def:chaincond}\,(C4) is accumulated in.
Condition~(S) on the conditional family of a choice value is then read in the
\emph{restriction of that order} to the elements that choice value
contributes, which is a total order on $\Uset_t$ however the record listed
it; the two conditions are independent, only the order restricting. Where no order arrives with that alphabet the
free choice is made as Definition~\ref{def:weighted} says, by the record's own
listing, and it is emitted. No declarer's order is overridden and no
presentation is refused for an order. Where the setting-dependent weighting is applied,
the emitted alphabet is the support of $P$ and its order is the product order
of Definition~\ref{def:weighted}: $v$ in its declarer's declared order and,
within each $v$, $f$ in the declared order of the choice alphabet, the
pairs off the support omitted; that is the order (S) is decided in on the
emitted family and the order $\idxm$ emits on it.

\emph{The general format.} For a specification presented as
Definition~\ref{def:spec}'s (a$'$)--(e$'$) the map fills the general format
of Definition~\ref{def:declaration}\,(DX1)--(DX6) as follows. In the
basic format --- one Alice sector, one measured factor, POVM form, two
messages neither reading the other and no party-local entry --- every clause
below is the clause of the same subject above, slot for slot, and the map
emits what those clauses emit.

\emph{The factor list} (DX1) is assembled and not declared as a whole: the
factor tagged $A$ is the Alice sector map the transmitter artefact declares,
where it declares one, and it is the first of the list; where the artefact
declares none, $\mathbb{T}=\{t_0\}$ and no $A$ factor is emitted. The factors
tagged $\Qsys$ are the entries of $\Spec_{\mathrm{rec}}$ in their declared
order, each carrying that entry's party label and the sector map of the last
target descriptor of that entry's receiver segment, which
Definition~\ref{def:chaincond}\,(C1) decides per factor. The one block kind
is the kind those maps carry, and the factors' agreement in it is decided by
the same four-component comparison V1 makes, taken on the factor list
(Definition~\ref{def:validator}, V1).
The product sector map, its order, the product alphabet and the
joint sectors are derived from that list by (DX1) and (DX3) and are
declared, formed and emitted by no one: what is emitted is the factor list and
the rules.

\emph{Alice's slots.} $\tau_A$ is what the artefact's marginal slot holds,
the block list along $\mathbb{T}$ included; Alice's measurement slot is the
artefact's, in the form the artefact declares --- the POVM slot or the
ensemble slot of (DX4) --- weighted by Definition~\ref{def:weighted} in that
form where a choice family is presented. The mode is read off the
marginal slot as above.

\emph{Per measured factor.} $(\Qsys_f,\SecSet^{f},(d_{f,s_f})_{s_f})$ and
$\{F^{f}_{v_f}\}_{v_f\in\Vset^{f}}$ from the last link of that factor's
receiver segment and, where that factor's entry of $\Spec_{\mathrm{rec}}$
declares a setting-dependent public choice, weighted after that segment by the
setting-dependent form of Definition~\ref{def:weighted}, so that factor $f$'s
emitted family is $P_f(\,\cdot\mid\cdot\,)$ times that factor's delivered
family over the support of $P_f$, and factor $f$'s emitted alphabet
$\Vset^{f}$ is that support in its product order. Each factor's weighting is
applied to its own family and to nothing else. No product of two factors'
families is formed by this map, whether emitted or handed on; so the subject of the
representability decision of Definition~\ref{def:weighted} is, as in the basic
format, the products \emph{that definition's} two constructions form, and no
product block \eqref{eq:productfamily} is a product this map forms.

\emph{The public part and the raw key.} $\Party$, $\phi$, $\hold$ verbatim
from $\Spec_{\mathrm{pub}}$ and the transcript verbatim: the sender,
the alphabet and the table of each message, declared on the emitted
alphabets. The joint announcement function $a$ is not a slot: it is the
composition \eqref{eq:transcript} of those tables, and $r$ is \eqref{eq:rdef}
with \emph{that} $a$, which is what Theorem~\ref{thm:sifting-and-key}\,(iv), (vi) is
stated of. $\CList$ is the concatenation above, each entry carrying in
addition its scope, and each with the datum its right-hand-side slot
holds and the tag that datum's type determines; for a party-local entry that
is the tag field it arrives with beside a real scalar
(Definition~\ref{def:homog}).

\emph{The index map in the general format.} $\idxm$ has one component per
factor alphabet $\Vset^{f}$ and one per factor sector set $\SecSet^{f}$, and,
where an $A$ factor is emitted, one on $\mathbb{T}$, which is the total order
that factor's own map fixes. Its components on $\Vset$ and on $\SecSet$ are
the product orders (DX1) and (DX3) derive from the factors' orders and are not
further choices, as its component on $\SecSet$ above is not; the joint sectors
are ordered with $t$ slowest by (DX1). The binding rule above is applied per
factor: where a factor's target family arrives with a declared order on its
alphabet, the compiler takes that order there. So condition~(S) of
Definition~\ref{def:exact} is decided per factor, in that factor's arriving
order, and per Alice sector on Alice's data; no sum over a product alphabet is
accumulated by this map or by its consumer.
\end{digestclaim}

% claim protocol-compiler/thm:typing
\digestlabel{thm:typing}{3.4}\pdfbookmark[2]{Theorem 3.4 Typing of an emitted declaration}{digest-protocol-compiler/thm:typing}
\renewcommand*{\RestatedName}{\hyperref{contract.pdf}{}{claim-protocol-compiler/thm:typing}{Theorem~3.4}}%
\begin{digestclaim}[Typing of an emitted declaration]
Let $\Spec$ be a well-formed specification
(Definition~\ref{def:well-formed-spec}) satisfying
Assumption~\ref{ass:spec}, and
let $\Decl$ be a declaration that the compilation map of Definition~\ref{def:compile}
\emph{emits}, that is, one for which no check of
Definition~\ref{def:validator} refused. Then $\Decl$, together with the index
map $\idxm$ emitted alongside it, satisfies the target conditions D(i), D(ii),
D(iv) and D(vi) of Definition~\ref{def:target}, and $\idxm$ is a bijection on
each alphabet it indexes, so that every emitted operator list is totally
ordered and its index is recoverable from its semantic label.
\end{digestclaim}

% claim protocol-compiler/def:phi
\digestlabel{def:phi}{3.6}\pdfbookmark[2]{Definition 3.6 Declared statistics functional}{digest-protocol-compiler/def:phi}
\renewcommand*{\RestatedName}{\hyperref{contract.pdf}{}{claim-protocol-compiler/def:phi}{Definition~3.6}}%
\begin{digestclaim}[Declared statistics functional]
For a declaration $\Decl$ and $\omega\in\Lin{A\Qsys}$ put
\begin{equation*}\tag{\getrefnumber{C-eq:jointdist}}\digesteqlabel{eq:jointdist}
  P_{\Decl}[\omega](x,v)\coloneqq\tr\!\big[(M_x\otimes F_v)\,\omega\big],
  \qquad
  \Phimap[\omega](\hat c)\coloneqq
  \sum_{(x,v):\,\phi(a(x,v))=\hat c}P_{\Decl}[\omega](x,v).
\end{equation*}
Both are defined by trace formulas and are therefore defined, and affine, on
all of $\Lin{A\Qsys}$. The feasible set $\Feas$ of
Definition~\ref{def:feasible} is contained in $\Pos{A\Qsys}$, and on it the
values are a probability distribution
(Theorem~\ref{thm:measurements-and-statistics}\,(iv), and (iv$'$) in the
general format).
\end{digestclaim}

% claim protocol-compiler/thm:measurements-and-statistics
\digestlabel{thm:measurements-and-statistics}{3.7}\pdfbookmark[2]{Theorem 3.7 Exact measurement families and the statistics functional}{digest-protocol-compiler/thm:measurements-and-statistics}
\renewcommand*{\RestatedName}{\hyperref{contract.pdf}{}{claim-protocol-compiler/thm:measurements-and-statistics}{Theorem~3.7}}%
\begin{digestclaim}[Exact measurement families and the statistics functional]
Let $\Spec$ be a well-formed specification
(Definition~\ref{def:well-formed-spec}) satisfying
Assumption~\ref{ass:spec}, and
let $\Decl$ be emitted by Definition~\ref{def:compile} in the
basic format. Then:
\begin{enumerate}[label=(\roman*),leftmargin=2.2em]
\item $\{M_x\}_{x\in\Xset}$ is a POVM on $A$ and, for each $s\in\SecSet$, the
block family $\{F_v^{(s)}\}_{v\in\Vset}$ is a POVM on $\Qsys_s$ --- each
predicate holding \emph{of the declared numbers}, decided exactly --- and each
of the completeness sums $\sum_xM_x$ and $\sum_vF_v^{(s)}$, formed in the
declared order, rounds nowhere, so that each computed sum \emph{is} the sum it
is written as (Definition~\ref{def:exact});
\item every $F_v$ and every $\hat O_k$ is sector-block-diagonal
\emph{by construction}: it is emitted as the list
$\big(F_v^{(s)}\big)_{s}$, respectively $\big(\hat O_k^{(s)}\big)_{s}$, and
the operator it denotes is $\bigoplus_sF_v^{(s)}$. The two formulations of
block-diagonality remain equivalent --- for an operator $F$ on $\Qsys$,
$\Pi_sF\Pi_{s'}=0$ for all $s\ne s'$ if and only if
$F=\bigoplus_s\Pi_sF\Pi_s$ --- and that equivalence is what licenses the
\emph{type}: a list of blocks denotes exactly the operators the vanishing
condition picks out, and no other. It is not checked, because the condition
has no data to be checked on;
\item hence D(iii) holds: its typing half by (ii) and check V1; its first
four predicates --- those on $\{M_x\}$ and on each block family
$\{F_v^{(s)}\}_v$ --- by (i); and its fifth, Hermiticity of
$\hat O_k^{(s)}$ for each $(k,s)$, by the last clause of check V3
(Definition~\ref{def:validator}), which decides it exactly on the emitted
constraint entries and which (i) does not state; and
\item $P_{\Decl}[\omega]$ of \eqref{eq:jointdist} and $\Phimap[\omega]$ are
affine in $\omega$, $\Phimap[\omega]$ depends on $\omega$ only through
$P_{\Decl}[\omega]$, and the fibres of $\phi\circ a$ partition
$\Xset\times\Vset$, so $\Phimap[\omega]$ is the pushforward of
$P_{\Decl}[\omega]$ and inherits its total mass. For every block list $\omega$
with $\omega^{(s)}\succeq0$ for each $s$ and $\tr\omega=1$ --- in particular
for every $\omega\in\Feas$ --- $P_{\Decl}[\omega]$ is a probability
distribution on $\Xset\times\Vset$ and $\Phimap[\omega]$ one on $\Chatset$.
\end{enumerate}

\emph{The general format.} Let $\Decl$ instead be emitted by the
general-format clauses of Definition~\ref{def:compile}, in the general format
of Definition~\ref{def:declaration}\,(DX1)--(DX6). Then (i)--(iv) hold with
the following readings, of which the four above are the case of the basic
format.
\begin{enumerate}[label=(\roman*$'$),leftmargin=2.6em]
\item for every measured factor $f$ and every $s_f\in\SecSet^{f}$ the
block family $\{F^{f,(s_f)}_{v_f}\}_{v_f\in\Vset^{f}}$ is a POVM on
$\Qsys_{f,s_f}$; and for every Alice sector $t$, in POVM form
$\{M^{(t)}_x\}_{x\in\Xset}$ is a POVM on $A_t$, and in ensemble form
every $N^{(t)}_x$ is Hermitian and positive semidefinite, with
$\sum_xN^{(t)}_x=\tau^{(t)}_A$ and $\tau^{(t)}_A\ne0$ --- each predicate
holding of the declared numbers, decided exactly, and each of these sums,
formed in its declared order, rounding nowhere. The product family
\eqref{eq:productfamily} is a POVM on $\Qsys$ and, in ensemble form, the
family \eqref{eq:ensemblepovm} is a POVM on $A$, both as identities of real
numbers, and neither is formed, decided or summed;
\item every $F^{f}_{v_f}$ is a block list against factor $f$'s sector map and
every $\hat O_k$ one against the sector map of its scope, and $\tau_A$, every
operator of Alice's measurement slot and every general $\Theta_k$ is a block
list along $\mathbb{T}$, each block diagonal by construction as in (ii); and
the product family is block diagonal for the product sector map,
$\Pi_sF_v\Pi_{s'}=0$ for $s\ne s'$, with the $s$-block of
\eqref{eq:productfamily};
\item hence D(iii) holds in its general-format reading
(Definition~\ref{def:target}); and
\item (iv) holds for every list $\omega$ over the
joint sectors with every block positive semidefinite and
$\tr\omega=1$ --- in particular for every $\omega\in\Feas$ --- with $F_v$ the
product \eqref{eq:productfamily} and, in ensemble form, $M_x$ the operator
\eqref{eq:ensemblepovm} of the declaration $\Decl$ denotes: $P_{\Decl}[\omega]$
is a probability distribution on $\Xset\times\Vset$ and $\Phimap[\omega]$ one
on $\Chatset$.
\end{enumerate}
\end{digestclaim}

\section*{4\quad Public choices as measurement weights}\pdfbookmark[1]{4 Public choices as measurement weights}{digest-group-protocol-compiler/def:weighted}

% claim protocol-compiler/def:weighted
\digestlabel{def:weighted}{4.1}\pdfbookmark[2]{Definition 4.1 Choice-weighted families}{digest-protocol-compiler/def:weighted}
\renewcommand*{\RestatedName}{\hyperref{contract.pdf}{}{claim-protocol-compiler/def:weighted}{Definition~4.1}}%
\begin{digestclaim}[Choice-weighted families]
Let $S$ be a finite-dimensional system, $\Tpub$ a finite non-empty
\emph{public-choice alphabet}
with declared weights $q_t>0$, $\sum_{t\in\Tpub}q_t=1$, and for each
$t\in\Tpub$ a POVM $\{N_{u|t}\}_{u\in\Uset_t}$ on $S$. The
\emph{choice-weighted family}
is indexed by the disjoint union
$\Vset\coloneqq\bigsqcup_{t\in\Tpub}\Uset_t\times\{t\}$ and defined by
\begin{equation*}\tag{\getrefnumber{C-eq:weighted}}\digesteqlabel{eq:weighted}
  N_{u,t}\coloneqq q_t\,N_{u|t},\qquad (u,t)\in\Vset .
\end{equation*}
Several independent choices compose by iterating \eqref{eq:weighted}; the
resulting weight is the product of the individual weights. \emph{What is
declared, and therefore what is decided, is the innermost family and each
weight list}: an iterated presentation declares $\{N_{u|t}\}_u$ for each value
of each choice together with one weight list per choice, the decision on the
conditional families is taken on those, and the intermediate
families are this module's own products, which inherit admissibility by
Lemma~\ref{lem:weighted}\,(i) at each step and carry no second declared
condition; the representability decision is taken on every product the map
forms, the intermediate ones included. This is the
\emph{setting-independent form}:
the weight of a choice value is one declared number, the same in every round
and whatever else either party's apparatus did.

\emph{A weighting that depends on a station's setting.} Let
$\{F_v\}_{v\in\Vset}$ be a POVM on $S$ over a finite alphabet $\Vset$ with a
declared order --- the target family of the last link of $\Chrec$, as
Definition~\ref{def:compile} takes it --- and let
$t:\Vset\to\mathcal{T}$ be a declared total map onto a finite non-empty
\emph{setting alphabet},
the
\emph{setting component}
of the outcome: the value of the private per-round setting of the measuring
apparatus that the outcome $v$ records. Where the chain was emitted by the reductions module, $t$ is the
class map of that module's joint alphabet
(\xref{reductions/def:joint-family}\,(a): $\Vset=\mathcal{Y}\times\mathcal{T}$
with classes $\Vset_t$); where the family is declared directly, $t$ is
declared with it, the singleton $\mathcal{T}=\{\ast\}$ being the case of no
setting. Let $\mathcal{F}$ be a finite non-empty
\emph{choice alphabet}
with a declared order and let $P(f\,|\,t)$,
$(t,f)\in\mathcal{T}\times\mathcal{F}$, be a \emph{stochastic matrix}:
$P(f\,|\,t)\ge0$ and $\sum_{f\in\mathcal{F}}P(f\,|\,t)=1$ for every
$t\in\mathcal{T}$. The
\emph{support of the weighting}
is
\begin{equation*}\tag{\getrefnumber{C-eq:support}}\digesteqlabel{eq:support}
  \Vset_P\coloneqq\big\{(v,f)\in\Vset\times\mathcal{F}:\ P(f\,|\,t(v))>0\big\},
\end{equation*}
ordered by the product order --- $v$ in the declared order of $\Vset$ and,
within each $v$, $f$ in the declared order of $\mathcal{F}$, the pairs off
the support omitted --- and the choice-weighted family in its
\emph{setting-dependent form}
is
\begin{equation*}\tag{\getrefnumber{C-eq:weighted-t}}\digesteqlabel{eq:weighted-t}
  N_{v,f}\coloneqq P(f\,|\,t(v))\,F_v,\qquad (v,f)\in\Vset_P .
\end{equation*}
The setting-independent form is the case $P(f\,|\,t)=q_f$ for every $t$, in
which \eqref{eq:weighted-t} is \eqref{eq:weighted} with
$\Tpub=\mathcal{F}$ and $N_{v|f}=F_v$ for every $f$ --- one conditional
family, weighted by a choice that does not see the setting; conversely,
\eqref{eq:weighted} with conditional families that differ with the choice is
not of the form \eqref{eq:weighted-t}, and the two forms are two constructions
overlapping in that case. \emph{In the overlap the two constructions do not
agree on an order, and the order is a property of the presentation and not of
the family}: the alphabet the first form produces is
$\bigsqcup_f\Uset_f\times\{f\}$, listed choice-major, while the support
\eqref{eq:support} of the second is listed outcome-major, so the two are
transposes of one another and one family declared through the two records has
its condition~(S) decided in two different orders. (S) is order-dependent by
construction (Definition~\ref{def:exact}), and the two orders are not
reconciled: what fixes the order is the record that was declared, and a
declarer that declares the same family both ways must expect the
decision its own record's order gives. \emph{$P$ is read through the setting
and through nothing else.} It is declared on $\mathcal{T}\times\mathcal{F}$
and reaches an outcome only through $t(v)$. Zero entries of $P$ are admitted, and the support
is what they are for: a pair $(v,f)$ with $P(f\,|\,t(v))=0$ is not an outcome,
is not emitted and indexes nothing, so every emitted element of
\eqref{eq:weighted-t} carries a strictly positive weight, exactly as $q_t>0$
requires of every element of \eqref{eq:weighted}. Which choice values a given
setting admits is therefore declared by the zero pattern of $P$ and not by a
zero-weight outcome. The
form is applied by the compiler \emph{after} the reduction chain, to the
family the last link delivered (Definition~\ref{def:compile}).

\emph{The reading of each declared condition.} The weights are
binary64 numbers and the conditional families binary64 matrices. So
$\sum_tq_t=1$ is decided
\emph{exactly}, with the sum accumulated in the declared order of $\Tpub$ ---
the order in which the declared public-choice record lists $\Tpub$, fixed
there because $\idxm$ does not index $\Tpub$
(Definition~\ref{def:compile}) --- and
required to round nowhere; and ``$\{N_{u|t}\}_u$ is a
POVM on $S$'' is exact admissibility in the sense of
Definition~\ref{def:exact}, per block where $S$ carries sectors: (H), (P) and
(C) per block; (S) on the completeness sum $\sum_{u\in\Uset_t}N_{u|t}$,
accumulated in the order stated below and required to round nowhere; and (T)
as the agreement --- in block kind, in block identifiers, in their order and in
their dimensions --- of the maps \emph{the conditional families of one
presentation themselves carry}: there is one such map for the presentation,
and the family the weighting produces is a block list against
\emph{it}. On a system carrying no sector decomposition, Alice's $A$, (T) is
empty. This condition is decided of each $t\in\Tpub$ separately, on the
declared numbers.

\emph{In which order (S) is read on a conditional family.} The family the
weighting produces is indexed by one alphabet, and that alphabet carries one
order: $\Xset$ where the presentation is Alice's, which arrives with its
declared order (Definition~\ref{def:spec}\,(a)); the physical receiver
alphabet where it is a measuring party's, which is the order the fibre sum of
Definition~\ref{def:chaincond}\,(C4) is accumulated in and, in Band A, is
the emitted alphabet's own. Condition~(S) on the conditional family
$\{N_{u|t}\}_u$ is decided in the \emph{restriction of that order} to the
elements the choice value $t$ contributes --- a total order on $\Uset_t$
whatever the presentation listed, so no compatibility of the record's listing
with it is required and none is decided. Where a declared order arrives with
that alphabet the compiler takes it, by the binding rule of
Definition~\ref{def:compile}, to which this definition adds nothing; where
none arrives, the choice the compiler makes and emits is the record's own
listing, $t$ in the order it lists $\Tpub$ in and, within each $t$, $u$ in the
order it lists $\Uset_t$ in. \emph{No presentation is refused for an order}:
what is fixed here is which order a decision is read in, not what a declarer
may declare.

\emph{The remaining declared conditions, at the same reading.} The strict
inequality $q_t>0$ is a domain condition on a declared real scalar, decided by
its sign, and Definition~\ref{def:exact} carries no
strict-inequality row because none is wanted --- $q_t\ne0$ is what
Lemma~\ref{lem:weighted}\,(ii) needs, and a negative $q_t$ would reverse the
sign of an element of the emitted family, which no public choice of known bias
can mean. For the setting-dependent form the same reading applies to $P$: each
entry is a binary64 number, $P(f\,|\,t)\ge0$ is decided by sign, each row sum
$\sum_fP(f\,|\,t)=1$ is decided exactly in the declared order of
$\mathcal{F}$ and required to round nowhere, the support \eqref{eq:support}
is decided by the sign of each entry, and ``$\{F_v\}$ is a POVM on $S$'' is
the exact admissibility the chain's (C3) already decided of the family the
last link delivered. \emph{This form has no second conditional family to
decide}: the choice value is drawn after the measurement, $\{F_v\}$ is the
whole of what the decision on the conditional families has as its subject
here, and (C3) is where it was taken.

\emph{What this module constructs, and the two emission decisions.} The
products $q_tN_{u|t}$ of \eqref{eq:weighted} and $P(f\,|\,t(v))F_v$ of
\eqref{eq:weighted-t} are formed \emph{exactly} on the declared dyadic
numbers. A product of two binary64 numbers need not be one, so the compiler
decides, before emitting anything, that every entry of every $N_{u,t}$,
respectively of every $N_{v,f}$, is a binary64 number --- on every product
this map forms, whether the product is emitted or handed on --- and that the
emitted family satisfies Definition~\ref{def:exact} --- (H), (P), (C) and, on
its own declared-order sum over $\Vset$, respectively over $\Vset_P$ in its
product order, (S). These are the two \emph{emission decisions}.

\emph{The lattice of the products.} Let the arriving entries be integer
multiples of $2^{-k}$ of modulus at most $1$, and let each $q_t$ be an
integer multiple of $2^{-k_q}$; then every entry of $q_tN_{u|t}$ is an integer
multiple of $2^{-(k+k_q)}$ of modulus at most $1$, so the constructed family
lies on a lattice $k_q$ steps finer than the one it was handed and is exactly
binary64 whenever $k+k_q\le53$. Where a $q_t$ is not a power of two the finer
lattice can leave binary64 altogether: the decimal literal $0.9$ is an integer
multiple of $2^{-53}$ and of no coarser one, so $k+k_q>53$ for every $k\ge1$
and the exact product of a declared entry by it need not be a binary64
number. The same arithmetic governs $P(f\,|\,t(v))F_v$ with the exponent
$k_P$ of the entries of $P$ in place of $k_q$: a $P$ whose non-zero entries
are dyadic with small exponent --- $\tfrac12$, $\tfrac14$, $1$ --- composes
with any family the chain delivers at $k\le53-k_P$, and a decimal literal in
$P$ can leave binary64 exactly as a decimal $q_t$ can.

\emph{The general format.} In the general format of
Definition~\ref{def:declaration}\,(DX1)--(DX6) both forms above are read
\emph{on the system of the party that declares them}, and nothing else of
this definition changes. At one measured factor, in POVM form
and with no factor tagged $A$, the readings below are the text above,
sentence for sentence.

\emph{Which system $S$ is.} For a presentation of Alice's it is $A$, and (T)
is \emph{not empty where a factor is tagged $A$}: the conditional families
are then block lists along $\mathbb{T}$ and (T) is the four-component
agreement of the maps they carry with that factor's map, as on any other
sectored system. For a presentation of a measuring party's it is that party's
\emph{own} measured factor: one presentation per factor, decided against that
factor's own map; \emph{a presentation indexed by a product alphabet has no
declared system to be a family on} --- the product is the rule
\eqref{eq:productfamily} and is no datum (Definition~\ref{def:compile}, the
general format). So the decision ``$\{N_{u|t}\}_u$ is
a POVM on $S$'' and the products this module forms are per factor, and the
representability decision keeps exactly the subject it has: the products
\emph{these two constructions} form, and no product block of
\eqref{eq:productfamily}.

\emph{The ensemble form of Alice's slot.} Where Alice's measurement slot
is declared in ensemble form (DX4), a choice-weighted presentation of it
declares, for each $t\in\Tpub$, a family $\{N_{x|t}\}_{x\in\Xset}$ of block
lists along $\mathbb{T}$ with every block Hermitian and positive
semidefinite, and \emph{in place of the POVM condition} the ensemble
condition of \eqref{eq:ensembleform}: for every Alice sector $t'$,
$\sum_{x}N^{(t')}_{x|t}=\tau^{(t')}_A$, the sum accumulated in the order
stated above and rounding nowhere. That condition is decided \emph{of each
$t$ separately and on the declared numbers}, with (H), (P) per block, that
identity in place of (C), and (T) as above. \emph{Why the
construction is the same construction.} With $N_{x,t}\coloneqq q_tN_{x|t}$ of
\eqref{eq:weighted}, the exact real numbers the data denote satisfy
\begin{equation*}\tag{\getrefnumber{C-eq:weighted-ensemble}}\digesteqlabel{eq:weighted-ensemble}
  \sum_{(x,t)\in\Vset}N^{(t')}_{x,t}
  =\sum_{t\in\Tpub}q_t\sum_{x\in\Xset}N^{(t')}_{x|t}
  =\Big(\sum_{t\in\Tpub}q_t\Big)\tau^{(t')}_A=\tau^{(t')}_A
\end{equation*}
for every $t'$, by the decision just taken and by $\sum_tq_t=1$; so a
choice-weighted presentation in ensemble form whose conditional ensembles are
admitted denotes a weighted family with the ensemble identity of
\eqref{eq:ensembleform}. \emph{That is a consistency statement and not a
substitute for the decision on the emitted family}: the emitted family's own
(S), (H), (P) and its identity with the declared $\tau_A$ are decided at check
V3, on the emitted blocks and in the emitted order, and neither condition
implies the other: \eqref{eq:weighted-ensemble} is an identity of real
numbers and the emitted decision is about the additions that compute it. In
untrusted-source mode there is no ensemble form and this paragraph has
no subject.

\emph{The setting-dependent form, per measured factor.} It is applied by the
compiler after \emph{that factor's} receiver segment, to the family that
factor's last link delivered, with that factor's own setting map, support
\eqref{eq:support} and product order; the setting component read is a
component of \emph{that factor's} outcome, the private per-round setting of
the station that measures that factor, whichever party's label the factor
carries. This form has no second
conditional family here either, per factor: (C3) of
Definition~\ref{def:chaincond}, read for that factor, is where the delivered
family was decided.
\end{digestclaim}

% claim protocol-compiler/lem:weighted
\digestlabel{lem:weighted}{4.2}\pdfbookmark[2]{Lemma 4.2 Properties of the choice-weighted families}{digest-protocol-compiler/lem:weighted}
\renewcommand*{\RestatedName}{\hyperref{contract.pdf}{}{claim-protocol-compiler/lem:weighted}{Lemma~4.2}}%
\begin{digestclaim}[Properties of the choice-weighted families]
With the notation of Definition~\ref{def:weighted}, and reading its conditions
exactly:
\begin{enumerate}[label=(\roman*),leftmargin=2.2em]
\item $\{N_{u,t}\}_{(u,t)\in\Vset}$ is a POVM on $S$; and so is
$\{N_{v,f}\}_{(v,f)\in\Vset_P}$ of the setting-dependent form, over the
support of $P$;
\item if $S=\Qsys$ carries the sector decomposition and every $N_{u|t}$ is
sector-block-diagonal, then so is every $N_{u,t}$, and conversely (for
$q_t>0$); equivalently, on block lists the weighting acts blockwise,
$N_{u,t}^{(s)}=q_t\,N_{u|t}^{(s)}$, so a weighted family is a block list
against the same sector map as the conditional families it was built from.
The same holds of the setting-dependent form: every $N_{v,f}$ is a scalar
multiple of the $F_v$ it was built from, $N_{v,f}^{(s)}=P(f\,|\,t(v))\,F_v^{(s)}$
block by block, and $N_{v,f}$ is sector-block-diagonal iff $F_v$ is, the
weight being strictly positive on the support;
\item \emph{(the choice marginal, for the setting-independent form)} for every
$\omega\in\Lin{S}$ with $\tr\omega=1$ and every $t\in\Tpub$,
\begin{equation*}\tag{\getrefnumber{C-eq:choicemass}}\digesteqlabel{eq:choicemass}
  \sum_{u\in\Uset_t}\tr[N_{u,t}\,\omega]=q_t ,
\end{equation*}
independently of $\omega$; and for any complementary system $R$ and any
$\omega\in\Lin{RS}$ of unit trace likewise
$\sum_{u}\tr[(\idon{R}\otimes N_{u,t})\,\omega]=q_t$, the operator being
written with its factors in the declared order of the tensor factors. For a
choice declared on Bob's side, $S=\Qsys$ and $R=A$, this reads
$\sum_u\tr[(\idon{A}\otimes N_{u,t})\omega]=q_t$; for a choice declared on
Alice's side, $S=A$ and $R=\Qsys$, it reads
$\sum_u\tr[(N_{u,t}\otimes\idon{\Qsys})\omega]=q_t$. \emph{This part is
asserted of the setting-independent form and of it alone};
\item \emph{(what is asserted of the setting-dependent form instead)}
forgetting the choice value returns the family the chain delivered,
\begin{equation*}\tag{\getrefnumber{C-eq:forgetf}}\digesteqlabel{eq:forgetf}
  \sum_{f:\,(v,f)\in\Vset_P}N_{v,f}=F_v\qquad\text{for every }v\in\Vset ,
\end{equation*}
an exact operator identity on $S$; and the marginal of the choice value is
the declared matrix applied to the \emph{setting marginal} of the outcome
statistics,
\begin{equation*}\tag{\getrefnumber{C-eq:fmarginal}}\digesteqlabel{eq:fmarginal}
  \sum_{v:\,(v,f)\in\Vset_P}\tr[N_{v,f}\,\omega]
  =\sum_{t\in\mathcal{T}}P(f\,|\,t)\sum_{v:\,t(v)=t}\tr[F_v\,\omega]
  \qquad\text{for every }\omega\in\Lin{S},
\end{equation*}
which depends on $\omega$ only through the \emph{setting marginal}
$m_t(\omega)\coloneqq\tr[\Pi_t\omega]$, $\Pi_t\coloneqq\sum_{v:\,t(v)=t}F_v$,
with $\Pi_t\succeq0$ and $\sum_t\Pi_t=\idon{S}$. If $P(f\,|\,t)=q_f$ for
every $t$ it equals $q_f$ for every unit-trace $\omega$ --- the
setting-independent case, and \eqref{eq:choicemass} in the form the first
construction gives it. If every class sum $\Pi_t$ is a multiple
$\lambda_t\idon{S}$ of the identity --- as for the identity link, where the
emitted family is the physical joint family --- it is state-independent as
well, with value $\sum_tP(f\,|\,t)\lambda_t$. Otherwise it is, in general, a
function of the state, and \emph{no identity of the form
\eqref{eq:choicemass} is asserted for this form}.
\end{enumerate}
\end{digestclaim}

% claim protocol-compiler/thm:choice-encoding
\digestlabel{thm:choice-encoding}{4.4}\pdfbookmark[2]{Theorem 4.4 Encoding of public choices in the measurement operators}{digest-protocol-compiler/thm:choice-encoding}
\begin{mainclaim}
\renewcommand*{\RestatedName}{\hyperref{contract.pdf}{}{claim-protocol-compiler/thm:choice-encoding}{Theorem~4.4}}%
\begin{digestclaim}[Encoding of public choices in the measurement operators]
Let a per-round public choice of known bias --- a basis choice, a choice
between test and generation rounds, a probe setting --- be declared on either
party's side with alphabet $\Tpub$, weights $(q_t)$ and conditional POVMs
$\{N_{u|t}\}$, under Assumption~\ref{ass:coin}. Then the compiler forms the
choice-weighted family \eqref{eq:weighted}, and (iv) below holds. Where that
family is the emitted one --- the presentation is Alice's, or it is a
measuring party's and the last link of that party's receiver segment
delivers the weighted family itself, on its own alphabet, as its
target POVM, as on the Band A route
(Definition~\ref{def:banda}\,(b)) --- the compiler emits it, and (i)--(iii)
hold as well:
\begin{enumerate}[label=(\roman*),leftmargin=2.2em]
\item the emitted family is a POVM and preserves sector block-diagonality
(Lemma~\ref{lem:weighted}\,(i),(ii)), so the encoding cannot break D(iii) ---
and it is emitted only after the two decisions of
Definition~\ref{def:weighted}, so that what is emitted is binary64 data whose
own predicates D(iii) is then decided on
(Theorem~\ref{thm:measurements-and-statistics}\,(i), and (i$'$) in the
general format);
\item the encoding is \emph{exact}: for every state, the marginal probability
of the choice value $t$ computed from the emitted declaration equals the
declared weight $q_t$ (Lemma~\ref{lem:weighted}\,(iii));
\item the choice value is a component of the joint outcome, hence enters every
downstream quantity only through $P_{\Decl}[\omega]$ and
$\Phimap[\omega]$ of Definition~\ref{def:phi}; no factor $q_t$ is emitted
anywhere outside a POVM element; and
\item the alternative encoding --- declaring the conditional families
$\{N_{u|t}\}$ together with a separate state-independent fraction $q_t$ to be
applied outside the operators --- is never emitted, and a specification
requesting it is refused.
\end{enumerate}
Where the presentation is a measuring party's and the last link of that
party's receiver segment delivers any other family, the weighted family is
the physical receiver family handed to the chain
(Definition~\ref{def:spec}\,(b)) and is not emitted: it is a POVM on the
physical receiver space, and for every unit-trace state there the marginal
probability of $t$ computed from it is $q_t$
(Lemma~\ref{lem:weighted}\,(i), (iii)); no factor $q_t$ is emitted anywhere
outside a POVM element, and (iv) holds. The emitted family is that link's
target POVM, whose own conditions are (C3)'s and check V3's
(Definition~\ref{def:chaincond}), and
nothing here asserts (ii) or (iii) of it; the declared test-statistics model
is the output-state model, as (vii) says of the setting-dependent form.
Let, further, a public choice $f\in\mathcal{F}$ be declared whose weight
depends on the private setting of the apparatus that measures the
measured system, as a
stochastic matrix $P(f\,|\,t)$ read through the setting map
$t:\Vset\to\mathcal{T}$ of the family $\{F_v\}_{v\in\Vset}$ the last link of
$\Chrec$ delivered --- the setting-dependent form of
Definition~\ref{def:weighted}; in the general format, at a measured factor,
read through that factor's setting map on the family the last link of that
factor's receiver segment delivered, with $S$ that factor's system --- under
Assumption~\ref{ass:coin}. Then the compiler applies it \emph{after} the
reduction chain, emits the family \eqref{eq:weighted-t},
$N_{v,f}=P(f\,|\,t(v))F_v$ over the support $\Vset_P$, and:
\begin{enumerate}[label=(\roman*),leftmargin=2.2em,resume]
\item the emitted family is a POVM and preserves sector block-diagonality
(Lemma~\ref{lem:weighted}\,(i),(ii)), every element being a scalar multiple
of an element the chain delivered, so the encoding cannot break D(iii); and it
is emitted only after the two decisions of Definition~\ref{def:weighted} on
the products $P(f\,|\,t(v))F_v$, so that what is emitted is binary64 data
whose own predicates D(iii) is then decided on
(Theorem~\ref{thm:measurements-and-statistics}\,(i), and (i$'$) in the
general format);
\item \emph{(soundness: the relabelling commutes with the chain)} for every
channel $\Lambda$ into $S$ and every family $\{G_v\}_{v\in\Vset}$ on its
input space with $\Lambda^{\dagger}(F_v)=G_v$ for every $v$,
\begin{equation*}\tag{\getrefnumber{C-eq:relabel-sound}}\digesteqlabel{eq:relabel-sound}
  \Lambda^{\dagger}\big(N_{v,f}\big)=P(f\,|\,t(v))\,G_v
  \qquad\text{for every }(v,f)\in\Vset_P :
\end{equation*}
the weighted target and the identically weighted input family stand in the
same relation under the \emph{same} channel, with the same target space, the
same sector map, the same constraint entries transported along the chain and
every pull-back unchanged. In this module's setting $G$ is the
$\gamma$-coarse-grained physical family of
Definition~\ref{def:chaincond}\,(C4), $G_v=\sum_{w\in\gamma^{-1}(v)}\Fph_w$,
so the emitted family and the physical family weighted by
$P(f\,|\,t(\gamma(w)))$ satisfy the pullback identity \eqref{eq:pullback}
with the outcome correspondence $(w,f)\mapsto(\gamma(w),f)$ restricted to the
support, wherever (C4) is decidable, and are related as (C4)'s families are
wherever it is not. This is \xref{reductions/lem:linear}\,(i) for the
stochastic matrix from $\Vset$ to $\Vset\times\mathcal{F}$ with entries
$\delta_{vv'}P(f\,|\,t(v'))$;
\item \emph{(what is not asserted)} the exactness clause (ii) is asserted for
the setting-independent form and for it alone. For this form the marginal of
$f$ is the quantity \eqref{eq:fmarginal} of
Lemma~\ref{lem:weighted}\,(iv), a function of the state through the setting
marginal, equal to the declared $q_f$ for every state when $P(f\,|\,t)=q_f$
and state-dependent in general otherwise; nothing here asserts otherwise, no
state-independent fraction is emitted, and the declared test-statistics model
is the output-state model.
What (iv) says of the fixed-fraction encoding holds of this form unchanged;
\item \emph{(refusals)} a weighting presented through anything but the
setting map --- indexed by $\Vset$, by the device outcome, or by any
component of $v$ other than $t(v)$ --- is refused; a $P$ with a negative
entry or a row whose declared-order sum is not exactly $1$ is refused; and
the two emission decisions of Definition~\ref{def:weighted} refuse, for a
product that is not representable or with the predicate that failed. Each
refusal emits no declaration.
\end{enumerate}
\end{digestclaim}
\end{mainclaim}

\section*{5\quad Sifting, key computability and feasibility}\pdfbookmark[1]{5 Sifting, key computability and feasibility}{digest-group-protocol-compiler/thm:sifting-and-key}

% claim protocol-compiler/thm:sifting-and-key
\digestlabel{thm:sifting-and-key}{5.1}\pdfbookmark[2]{Theorem 5.1 Announcement locality, sifting and party computability}{digest-protocol-compiler/thm:sifting-and-key}
\renewcommand*{\RestatedName}{\hyperref{contract.pdf}{}{claim-protocol-compiler/thm:sifting-and-key}{Theorem~5.1}}%
\begin{digestclaim}[Announcement locality, sifting and party computability]
Let $\Spec$ be a well-formed specification
(Definition~\ref{def:well-formed-spec}) satisfying
Assumption~\ref{ass:spec}, so
that in particular $r^{0}_{\hold}$ maps into $\Zval$
(Assumption~\ref{ass:spec}\,(ii)), and let $\Decl$ be emitted by
Definition~\ref{def:compile}. Then:
\begin{enumerate}[label=(\roman*),leftmargin=2.2em]
\item \emph{(announcement locality)} $a_A$ is a function of
$\rd{\mathsf A}(x,v)$ alone and $a_B$ of $\rd{\mathsf B}(x,v)$ alone, and the
joint announcement is \eqref{eq:jointann};
\item \emph{(D(v))} the specification's $\varsigma$ --- consumed, and not
itself a slot of \eqref{eq:declaration} --- witnesses announcement-determined
sifting: $r(x,v)=\perp\iff\varsigma(a(x,v))=\mathrm{fail}$;
\item \emph{(D(viii))} $r$ is party-computable: the total map
\begin{equation*}\tag{\getrefnumber{C-eq:rhold}}\digesteqlabel{eq:rhold}
  r_{\hold}(o,\kappa)\coloneqq
  \begin{cases}
    r^{0}_{\hold}(o,\kappa), & \varsigma(\kappa)=\mathrm{pass},\\
    \perp, & \varsigma(\kappa)=\mathrm{fail},
  \end{cases}
  \qquad (o,\kappa)\in\Oset{\hold}\times\Kann,
\end{equation*}
satisfies \eqref{eq:partycomp}.
\item \emph{(transcript protocols)} if the
announcement slot holds a transcript protocol
(Definition~\ref{def:declaration}\,(DX5)), then for
every $m$ the $m$-th component of $a(x,v)$ is a function of
$\rd{\pi_m}(x,v)$ and of the components before it, and of nothing else; $a$ is
total; and (ii) and (iii) hold with this $a$. For $L=2$,
$\pi=(\mathsf A,\mathsf B)$ and $a_2$ independent of $\kappa_1$ this is (i);
\item \emph{(ensemble form)} if Alice's measurement slot is declared
in ensemble form and the conditions of Definition~\ref{def:exact} for that
form hold, then the operators \eqref{eq:ensemblepovm} satisfy
$M^{(t)}_x\succeq0$ and $\sum_xM^{(t)}_x=\idon{A_t}$ for every $t$, and
$(\tau^{(t)}_A)^{1/2}M^{(t)}_x(\tau^{(t)}_A)^{1/2}=N^{(t)}_x$: the denoted
declaration satisfies the POVM clause of D(iii) on $A$ as an identity of real
numbers, and $\tr[M_x\tau_A]=\sum_t\tr N^{(t)}_x$.
\item \emph{(the compiled general format)} if $\Spec$ is presented as
Definition~\ref{def:spec}\,(a$'$)--(e$'$) and $\Decl$ is emitted by the
general-format clauses of Definition~\ref{def:compile}, then: the
announcement slot of $\Decl$ holds the transcript $\Spec_{\mathrm{pub}}$
declares, message for message; the joint announcement function of $\Decl$ is
\eqref{eq:transcript} evaluated at \emph{those} tables and is a slot of
nothing; and Alice's measurement slot holds the form the transmitter artefact
declares. Consequently (iv) applies to $\Decl$, hence (ii) and (iii) hold of
it, and (v) applies to it exactly where that form is the ensemble form.
\end{enumerate}
Moreover, under the default party partition these read:
$a_A:\Xset\to\CAset$, $a_B:\Vset\to\CBset$, and $r(x,v)=r_{\mathsf A}(x,a(x,v))$
for $\hold=\mathsf A$, respectively $r(x,v)=r_{\mathsf B}(v,a(x,v))$ for
$\hold=\mathsf B$.
\end{digestclaim}

% claim protocol-compiler/thm:feasibility
\digestlabel{thm:feasibility}{5.5}\pdfbookmark[2]{Theorem 5.5 Feasibility by an explicit or a verified candidate}{digest-protocol-compiler/thm:feasibility}
\renewcommand*{\RestatedName}{\hyperref{contract.pdf}{}{claim-protocol-compiler/thm:feasibility}{Theorem~5.5}}%
\begin{digestclaim}[Feasibility by an explicit or a verified candidate]
Let $\Decl$ be a candidate declaration satisfying D(i)--D(iii), and let
$\Feas$ be as in Definition~\ref{def:feasible}. (The hypothesis
is on the conditions alone: nothing below uses that $\Decl$ was emitted by
Definition~\ref{def:compile}, and Theorem~\ref{thm:validator-soundness} applies this theorem to
candidates.)
\begin{enumerate}[label=(\roman*),leftmargin=2.2em]
\item If $\CList=\varnothing$ then $\Feas\ne\varnothing$. Explicitly, the
block list
\begin{equation*}\tag{\getrefnumber{C-eq:witness-empty}}\digesteqlabel{eq:witness-empty}
  \omega^{\star(s)}=
  \begin{cases}
    \tau_A\otimes\dfrac{\idon{\Qsys_s}}{d_{\Qsys}},
      & \text{trusted-source mode},\\[1.2ex]
    \dfrac{\idon{A}}{d_A}\otimes\dfrac{\idon{\Qsys_s}}{d_{\Qsys}},
      & \text{untrusted-source mode},
  \end{cases}
  \qquad s\in\SecSet ,
\end{equation*}
which denotes $\tau_A\otimes\idon{\Qsys}/d_{\Qsys}$, respectively
$\idon{A}/d_A\otimes\idon{\Qsys}/d_{\Qsys}$, lies in $\Feas$.
\item If $\CList\ne\varnothing$ and $\Decl$ satisfies D(vi) as well, which
makes each constraint leg of \eqref{eq:feas} a well-posed comparison, then
$\Feas\ne\varnothing$ holds if and only
if some block list $\omega$ satisfies the conditions of \eqref{eq:feas}. The
compiler decides those conditions \emph{exactly}, by finitely many arithmetic
operations and per block, on at most two candidates and in a fixed order:
first the \emph{default candidate} \eqref{eq:witness-empty}, which satisfies
every conjunct of \eqref{eq:feas} except possibly the constraint legs
(by the proof of (i)) and which the compiler forms in order to test and never
emits; and, if that candidate fails a constraint leg, a \emph{supplied}
witness $\omega^{\star}$, consumed and verified and never emitted. It refuses
exactly when the default candidate fails and either no witness is supplied or
the supplied one fails.
\end{enumerate}
\emph{The general format.} For a candidate declaration in the general format
of Definition~\ref{def:declaration}\,(DX1)--(DX6) satisfying D(i)--D(iii),
and D(vi) for part (ii), with $\Feas$ the set of \eqref{eq:feas-general},
both parts hold with the
joint sectors in place of the sectors. In (i) the default candidate is
the block list over the joint sectors
\begin{equation*}\tag{\getrefnumber{C-eq:witness-general}}\digesteqlabel{eq:witness-general}
  \omega^{\star(t,s)}=
  \begin{cases}
    \tau^{(t)}_A\otimes\dfrac{\idon{\Qsys_s}}{d_{\Qsys}},
      & \text{trusted-source mode},\\[1.2ex]
    \dfrac{\idon{A_t}}{d_A}\otimes\dfrac{\idon{\Qsys_s}}{d_{\Qsys}},
      & \text{untrusted-source mode},
  \end{cases}
  \qquad (t,s)\in\mathbb{T}\times\SecSet ,
\end{equation*}
and in (ii) the verification is the per-joint-block one listed in the proof,
which forms no array larger than one joint block. At $|\mathbb{T}|=1$ and one
measured factor \eqref{eq:witness-general} is \eqref{eq:witness-empty}
and the two lists of decisions coincide item for item.
\end{digestclaim}

\section*{6\quad The emitted validator}\pdfbookmark[1]{6 The emitted validator}{digest-group-protocol-compiler/def:validator}

% claim protocol-compiler/def:validator
\digestlabel{def:validator}{6.1}\pdfbookmark[2]{Definition 6.1 The emitted validator}{digest-protocol-compiler/def:validator}
\renewcommand*{\RestatedName}{\hyperref{contract.pdf}{}{claim-protocol-compiler/def:validator}{Definition~6.1}}%
\begin{digestclaim}[The emitted validator]
Given a candidate declaration \eqref{eq:declaration} with its emitted
index map, the checks are as follows. \emph{Every operator condition
below is decided exactly}, in the sense of Definition~\ref{def:exact}, and
every combinatorial condition exactly as well.
\begin{description}[leftmargin=3.0em,style=nextline]
\item[V1 (finiteness, totality, and the sector map)] every declared alphabet is
  finite and non-empty and every declared map is tabulated on all of its stated
  domain; every declared system carries a finite dimension;
  $d_{\Qsys}=\sum_{s\in\SecSet}d_s$; the restriction of $\idxm$ to $\SecSet$ is
  the total order the sector map fixes; \emph{and the sector maps
  agree}: there is one emitted map $\Smap$, and the map carried by each
  supplied $\Qsys$-operator family is compared against it in \emph{block
  kind}, in the \emph{list of block identifiers}, in their \emph{order} and in
  their \emph{dimensions}. A disagreement on any of the four is refused,
  naming the first position at which the lists differ. The families compared
  are exactly the $\Qsys$-operator families of \eqref{eq:declaration}, which
  is the list D(i) quantifies over and no wider one: the emitted family
  $\{F_v\}$, which the last link of $\Chrec$ delivers as its target POVM
  against the map the chain's producer wrote it against, and each constraint
  entry of $\CList$, declared in $\Spec_{\mathrm{con}}$ or transported by a
  link of either segment. Those two slots are the whole of
  \eqref{eq:declaration} that carries a $\Qsys$-operator
  (Definition~\ref{def:declaration}), so the list is exhaustive by
  construction rather than by enumeration; the objects outside it are compared
  where they belong. The entries of
  $\Spec_{\mathrm{con}}$ are compared as well as the transported ones:
  Definition~\ref{def:compile} fills $\CList$ from both sources,
  Definition~\ref{def:spec}\,(e) has each declared $\hat O_k$ carry its own
  sector map, and an entry written against a different map is the error this
  check exists for, whichever source it arrived from. The comparison is
  combinatorial --- it compares labels and integers, never matrix entries ---
  and it is the whole of the structural verification the emitted declaration
  receives: block-diagonality is carried by the type
  (Theorem~\ref{thm:measurements-and-statistics}\,(ii)), so the one structural error that can still
  arise is two parties disagreeing about what the sectors are, and V1 is the
  check for it.
\item[V2 (marginal slot)] the slot holds $\nomarg$, or a matrix $\tau_A$ on $A$
  --- one block, $A$ carrying no sector decomposition --- satisfying
  $\tau_A=\tau_A^{\dagger}$, $\tau_A\succeq0$ and $\tr\tau_A=1$, each decided
  exactly, and with the trace accumulated over the row index verified to
  round nowhere; a failure of the last is refused.
\item[V3 (operators, block by block)]
  $M_x=M_x^{\dagger}$ and $M_x\succeq0$ for each $x\in\Xset$, and
  $\sum_xM_x=\idon{A}$ with the sum, in the emitted index order, verified to
  round nowhere, a failure of which is refused;
  $F_v^{(s)}=F_v^{(s)\dagger}$ and $F_v^{(s)}\succeq0$ for each
  $(v,s)\in\Vset\times\SecSet$, and, for each $s$,
  $\sum_vF_v^{(s)}=\idon{\Qsys_s}$ with that sum verified to round nowhere, a
  failure being refused with the sector and the outcome index at which the
  first inexact addition occurred; and
  $\hat O_k^{(s)}=\hat O_k^{(s)\dagger}$ for each $(k,s)$. Every one of these
  is decided exactly. There is no off-block condition: an emitted
  $\Qsys$-operator is a block list, so there is no $\Pi_sF_v\Pi_{s'}$ or
  $\Pi_s\hat O_k\Pi_{s'}$ to evaluate and nothing that could be non-zero.
\item[V4 (key alphabet)] every tabulated value of $r$ lies in
  $\Zval\sqcup\{\perp\}$.
\item[V5 (sifting consistency)] for all $(x,v)$ and $(x',v')$ with
  $a(x,v)=a(x',v')$: $r(x,v)=\perp\iff r(x',v')=\perp$.
\item[V6 (constraint typing, including the tag)] when the marginal
  slot is filled: $\bowtie_k\in\{\preceq,=,\succeq\}$, and the entry's
  right-hand side is of one of the two admitted types, which is what fixes its
  tag (Definition~\ref{def:homog}): homogeneous, whose datum is of
  real scalar type, or general, whose datum is a matrix on $A$ with
  $\Theta_k=\Theta_k^{\dagger}$ decided exactly. When the slot holds
  $\nomarg$: $\bowtie_k\in\{\le,=,\ge\}$ and the right-hand side is of real
  scalar type. A real scalar right-hand side is accepted in
  trusted-source mode at every $d_A$, and is the licensed form there
  (Definition~\ref{def:homog}). The relation symbol, the type of the
  right-hand side and the tag that type fixes are compared exactly; an entry
  whose right-hand-side slot holds neither admitted object, or a complex-typed
  right-hand side in untrusted-source mode, is refused by type and not
  measured. The check also reads an object \eqref{eq:declaration} does not
  carry, the tag field of an \emph{arriving} entry, in $\Spec_{\mathrm{con}}$
  or transported by a link: where an arriving entry carries an explicit tag
  field, V6 compares that field with the type of the datum it names, and a
  disagreement is refused. An entry of \eqref{eq:declaration} in
  either of the two forms above is the triple
  $(\hat O_k,\bowtie_k,\mathrm{rhs}_k)$ and holds no tag field independent of
  the type of its datum, the emitted tag being read off that type
  (Definition~\ref{def:homog}); so on a candidate declaration this clause has
  no subject, refuses nothing, and is not among the conjuncts D(vi) is paired
  with (Theorem~\ref{thm:validator-soundness}, proof). An entry carrying no tag field is not
  refused: its tag is the type of its datum and is read off it, which is what
  lets an entry declared as
  a real constant, carrying no tag field, be read here.
\item[V7 (feasibility; one-sided, three branches)] V7 passes in exactly three
  cases:
  \begin{enumerate}[label=(\alph*),leftmargin=2.2em]
  \item $\CList=\varnothing$; or
  \item $\CList\ne\varnothing$ and the default candidate
    \eqref{eq:witness-empty}, which the compiler forms and does not emit,
    passes the per-block verification of Theorem~\ref{thm:feasibility}\,(ii); or
  \item $\CList\ne\varnothing$, the default candidate fails, and a witness
    $\omega^{\star}$ is supplied \emph{as a block list against $\Smap$} --- the
    map it carries compared with $\Smap$ in block kind, in identifiers, in
    their order and in their dimensions, exactly as V1 compares the map of a
    declared $\Qsys$-operator family --- and passes that same verification.
    This is the one place a witness's sector map is compared, and it is
    here rather than at V1 because $\omega^{\star}$ is a block list on
    $A\Qsys$, is no slot of \eqref{eq:declaration}, and is quantified over by
    no target condition (Definition~\ref{def:target}).
  \end{enumerate}
  Otherwise V7 fails and the candidate is refused, the refusal naming, where a
  witness was supplied and rejected on its map, the first position at which
  the two lists differ. The verification in (b) and (c) is the $|\SecSet|$
  block positivity decisions --- Hermiticity and semidefiniteness per block
  --- the trace, the marginal equality in trusted-source mode, and one leg per
  constraint entry, all decided exactly, on quantities formed in exact
  arithmetic (Theorem~\ref{thm:feasibility}\,(ii)). It is taken on a
  candidate satisfying D(i)--D(iii) and D(vi), which V1--V3 and V6 decide; a
  candidate failing one of them is refused there. V7 is the one check of the
  list that
  is \emph{not} equivalent to its target condition: it implies D(vii) and is
  implied by nothing weaker than a candidate that verifies
  (Theorem~\ref{thm:validator-soundness}\,(iii)).
\item[V8 (party computability)] $a_A$ and $a_B$ are tabulated on
  $\Oset{\mathsf A}$, $\Oset{\mathsf B}$; and for all $(x,v)$, $(x',v')$ with
  $\rd{\hold}(x,v)=\rd{\hold}(x',v')$ and $a(x,v)=a(x',v')$:
  $r(x,v)=r(x',v')$.
\end{description}

\emph{The general format.} On a candidate declaration in the general format
of Definition~\ref{def:declaration}\,(DX1)--(DX6) the eight checks read as
follows, every operator condition again being decided exactly. In the
basic format --- one Alice sector, one measured factor,
POVM form, two messages neither reading the other and no party-local
entry --- each reading is the check of the
same name above, conjunct for conjunct, so on such a candidate the two
readings admit and refuse alike.

\emph{V1 in the general format.} Finiteness and totality as above, every
table of the transcript included: each $a_m$ is tabulated on all of
$\Oset{\pi_m}\times\mathcal{C}_1\times\dots\times\mathcal{C}_{m-1}$, which is
also where the signature clause of D(viii) lives: a table whose stated
domain names a later message's alphabet is not of that type and is refused
here. \emph{The sector map is one factorised map}: at most one factor tagged
$A$ and, where there is one, first in the list; at least one factor tagged
$\Qsys$; \emph{one block kind across the factors}, a disagreement being
refused, naming the first factor at which the kind differs;
$d_{\Qsys_f}=\sum_{s_f}d_{f,s_f}$ for every factor and
$d_{\Qsys}=\prod_fd_{\Qsys_f}$; and, where an $A$ factor is declared,
$d_A=\sum_td_{A_t}$. The restriction of $\idxm$ to each factor's sector set
is the total order that factor's own map fixes, and its restriction to
$\mathbb{T}$ is the order the $A$ factor's map fixes; the product orders on
$\SecSet$ and $\Vset$ and the order on the joint sectors are derived
by (DX1) and (DX3) and are not compared, there being no second datum to
compare them with. \emph{And the maps agree, factor by factor}: the map
carried by each declared family $\{F^{f}_{v_f}\}_{v_f}$ is compared against
the map of factor $f$, and the map carried by each entry's operator
$\hat O_k$ against the map of that entry's scope, in the same four
components, with the same refusal naming the first differing position.
\emph{Alice's declared data are compared too, against the map of the
factor tagged $A$}: the map carried by $\tau_A$, the map carried by each
operator of Alice's measurement slot --- $M_x$ in POVM form, $N_x$ in
ensemble form --- and the map carried by each $\Theta_k$ of an entry declared
general are each compared against that factor's map in the same four
components (block kind, block identifiers, their order, dimensions), with the
same refusal naming the first differing position. That comparison is what
makes condition~(T) of Definition~\ref{def:exact} decided for Alice's data,
which its general-format paragraph calls non-empty for them; it is decided
here and nowhere else, and without it a candidate whose Alice block list
carried a different block kind, a different identifier list or a different
order would pass every check and violate D(i). Where no factor is tagged $A$
there is one Alice block and no map to compare, as in V1 above. \emph{The list
of compared objects is again exhaustive by construction}: the objects D(i)
quantifies over are the per-factor families, the entries' operators and
Alice's declared data. The product family is no datum, and the
lifted operator of a scoped entry is no datum; the descriptors that the
last links of a chain produce, decided per factor at
Definition~\ref{def:chaincond}\,(C1), and a supplied witness, compared at
branch~(c) of V7, stay outside the list; and the physical families stay outside
it too, one per factor.

\emph{V2 in the general format.} The slot holds $\nomarg$, or a block list
$(\tau^{(t)}_A)_t$ against the $A$ factor's map with
$\tau^{(t)}_A=\tau^{(t)\dagger}_A$ and $\tau^{(t)}_A\succeq0$ for every $t$
and $\sum_t\tr\tau^{(t)}_A=1$, each trace accumulated over the row index and
the block traces then added in the declared order of $\mathbb{T}$, with
\emph{every} one of those additions verified to round nowhere; a failure of
the last is refused. \emph{Where Alice's measurement slot is in ensemble form,
additionally $\tau^{(t)}_A\ne0$ for every $t$}, a comparison of declared
entries with the datum $0$, a failure being refused naming the sector. At
$|\mathbb{T}|=1$ this is V2 above, the one block being the whole of $A$.

\emph{V3 in the general format.} \emph{Per measured factor}: for every
$f\in\mathbb{F}$, every $v_f\in\Vset^{f}$ and every $s_f\in\SecSet^{f}$,
$F^{f,(s_f)}_{v_f}$ Hermitian and positive semidefinite, and for each $s_f$
the sum $\sum_{v_f}F^{f,(s_f)}_{v_f}=\idon{\Qsys_{f,s_f}}$ with its
declared-order sum verified to round nowhere, a failure being refused with
the \emph{factor}, the sector and the outcome index at which the first
inexact addition occurred. \emph{On Alice's data, block by block along
$\mathbb{T}$}: in POVM form $M^{(t)}_x$ Hermitian and positive
semidefinite for every $(x,t)$ and $\sum_xM^{(t)}_x=\idon{A_t}$ for every
$t$, the sum in the emitted index order verified to round nowhere, a failure
being refused with the Alice sector; in ensemble form $N^{(t)}_x$
Hermitian and positive semidefinite for every $(x,t)$ and, \emph{in place of
the completeness identity}, the identity $\sum_xN^{(t)}_x=\tau^{(t)}_A$ of
\eqref{eq:ensembleform} decided entry by entry against the declared block of
$\tau_A$, a failure being refused with the sector and the first differing
entry, with that declared-order sum verified to round nowhere, a failure
being refused with the sector and the index of the first inexact addition.
\emph{An ensemble-form slot beside a marginal slot holding $\nomarg$ is
refused}; this is D(iii)'s general-format clause, and it is a decision on
which of two slots is filled, not on a number. And $\hat O_k^{(s_{f_k})}$ is
Hermitian for every entry and every block of its scope. \emph{Nothing is
decided on a product block \eqref{eq:productfamily}, on a lifted operator or
on the denoted family \eqref{eq:ensemblepovm}}: none of the three is a datum,
and that the first is a POVM on
$\Qsys$ and the third a POVM on $A$ are theorems
(Definition~\ref{def:declaration}\,(DX3), Theorem~\ref{thm:sifting-and-key}\,(v)).

\emph{V4, V5, V7 in the general format.} V4 and V5 keep their wording, with
$a$ the function \eqref{eq:transcript} derives from the transcript's tables;
V5's pairs are compared on the values of that function and on nothing about
how it was built. V7 keeps its three branches and its refusal, and what its
branches (b) and (c) verify is the per-\emph{joint}-sector verification of
Theorem~\ref{thm:feasibility}\,(ii) in the general format, on a candidate
satisfying D(i)--D(iii) and D(vi), in their general-format reading, which
V1--V3 and V6 decide (a candidate failing one of them is refused there): the
default candidate is
\eqref{eq:witness-empty} read per joint sector, and a supplied
$\omega^{\star}$ is required to be a block list over the joint sectors, its
map compared in the same four components at branch~(c) and there alone.

\emph{V6 in the general format.} When the marginal slot is filled:
$\bowtie_k\in\{\preceq,=,\succeq\}$, and the entry's right-hand side is of
one of the \emph{three} admitted forms, which with the arriving tag field is
what fixes its tag (Definition~\ref{def:homog}): homogeneous with a real
scalar, general with a block list $(\Theta^{(t)}_k)_t$ along $\mathbb{T}$
decided Hermitian block by block, or party-local with a real scalar. When the
slot holds $\nomarg$: the right-hand side is of real scalar type, either
untagged or tagged party-local. \emph{The relation symbol is decided by form
and not by mode alone}: for an \emph{untagged} entry
$\bowtie_k\in\{\preceq,=,\succeq\}$ when the slot is filled and
$\bowtie_k\in\{\le,=,\ge\}$ when it holds $\nomarg$, exactly as in V6 above;
for an entry declared party-local, $\bowtie_k\in\{\preceq,=,\succeq\}$
\emph{in either case}, its reading being an operator relation
(Definition~\ref{def:homog}, the relation-symbol clause of the party-local
form; Definition~\ref{def:target}, D(vi) in the general format), with the
scalar symbols admitted as well where every row of that entry is $1\times1$
--- every Alice block, and every sector block of every measured factor other
than the entry's scope, of dimension one, $d_A=1$ at one measured factor
included --- the three operator symbols being the three scalar symbols there.
This is a decision on a declared symbol beside a declared tag and measures
nothing, and at an entry not declared party-local it is the clause of V6
above word for word, so nothing that clause admits or refuses moves.
\emph{Every entry's scope names a measured factor of the emitted sector map},
a comparison of a label against the factor list, a failure being refused
naming the entry; this is the general-format half of D(vi) and is decided and
never measured, as the relation symbol is. \emph{A party-local entry is the
one that arrives carrying that tag value}, and it carries a tag field the
type of its datum does not determine, that type not separating it from a
homogeneous one. On a candidate carrying one, the arriving-tag clause of V6
above therefore has a subject, and it is this comparison: the datum is of the
type \emph{that} tag admits, a real scalar, and the scope names a measured
factor; a party-local tag beside a matrix right-hand side is refused. The
sentence of V6 above that an entry holds no tag field independent of the type
of its datum is true of the two untagged forms, homogeneous and general, and
of them alone. \emph{Nothing is inferred}: an entry arriving without the
party-local tag is never read party-locally. So the arriving-tag clause
refuses no party-local entry beside a real scalar, which a literal reading of it
would otherwise do, and wherever no entry is tagged the reading is that of V6
above (Definition~\ref{def:homog}).

\emph{V8 in the general format.} The announcement clause reads the
transcript: every $a_m$ is tabulated on all of its stated domain, which is
V1's clause restated where D(viii) pairs with it; the key clause is unchanged,
with $a$ the derived function.
\end{digestclaim}

% claim protocol-compiler/thm:validator-soundness
\digestlabel{thm:validator-soundness}{6.2}\pdfbookmark[2]{Theorem 6.2 Soundness, completeness and decidability of the validator}{digest-protocol-compiler/thm:validator-soundness}
\begin{mainclaim}
\renewcommand*{\RestatedName}{\hyperref{contract.pdf}{}{claim-protocol-compiler/thm:validator-soundness}{Theorem~6.2}}%
\begin{digestclaim}[Soundness, completeness and decidability of the validator]
For a candidate declaration with finite tabulated data:
\begin{enumerate}[label=(\roman*),leftmargin=2.2em]
\item \emph{(soundness, on all eight conditions)} if V1--V8 all
pass then D(i)--D(viii) all hold, of the declared numbers;
\item \emph{(completeness, on seven of the eight)} if D(i)--D(vi)
and D(viii) hold then V1--V6 and V8 all pass;
\item \emph{(V7 is one-sided, and the gap is exactly its third branch)} the
converse of (i) fails, and it fails only at V7: D(vii) does \emph{not} imply
V7. For a witness, take untrusted-source mode with $\Qsys$ a qubit
carrying the single sector $\SecSet=\{s_0\}$, every other slot well-formed,
and the one-entry list
$\CList=\{(\dyad{0}_{\Qsys},\ \ge,\ 1)\}$, whose entry reads
$\tr[(\idon{A}\otimes\dyad{0})\,\omega]\ge1$. Then $\hat O_1$ is a
one-block list, Hermitian, the typing is that of D(vi), and
$\omega=\eta\otimes\dyad{0}$ is feasible for any state $\eta\in\Seq{A}$ --- it
is Hermitian and positive semidefinite outright, has unit trace, and its
constraint value is exactly $1$ --- so $\Feas\ne\varnothing$ and D(i)--D(viii)
all hold. But $\CList\ne\varnothing$, and the default candidate
\eqref{eq:witness-empty} that branch (b) of V7 tries first gives
$\tr[(\idon{A}\otimes\dyad{0})(\idon{A}/d_A\otimes\idon{\Qsys}/2)]=\tfrac12$,
which fails $\ge1$; so with no witness supplied branch (c) is unavailable, V7
fails and the well-formed declaration is refused. Exactly the same
happens for every declaration with $\CList\ne\varnothing$ presented without a
witness whose feasible set is non-empty but does not contain
\eqref{eq:witness-empty}. The refusal errs on the safe side --- a refused
declaration is never consumed --- and the incompleteness is deliberate: deciding D(vii) in general is a
semidefinite feasibility problem this module does not solve;
\item \emph{(decidability, and what is computed)} each of V1--V6 and V8 is
decided by finite enumeration over the declared alphabets together with
finitely many arithmetic operations on the declared matrices; V7 is decided in
the same sense on the candidate of its branch (b) or (c), and is otherwise a
refusal. The arithmetic part consists of \emph{exact decisions}: an
identity, a semidefiniteness, or the exactness of an addition. The following
counts give the complete applicable predicate inventory for well-typed data,
not the number an early-terminated execution necessarily performs.
Write $I_{\mathrm{tr}}=1$ in trusted-source mode and $I_{\mathrm{tr}}=0$ when
the marginal slot holds $\nomarg$. The complete arithmetic inventory of V2,
V3 and V6 is
\begin{equation*}\tag{\getrefnumber{C-eq:predcount}}\digesteqlabel{eq:predcount}
  \underbrace{3I_{\mathrm{tr}}}_{\text{V2}}
  \;+\;\underbrace{2|\Xset|+1}_{\text{V3 on }A}
  \;+\;\underbrace{|\SecSet|\,\big(2|\Vset|+1\big)}_{\text{V3 on }\Qsys}
  \;+\;\underbrace{m\,|\SecSet|}_{\text{V3 on }\CList}
  \;+\;\underbrace{m_{\text{gen}}}_{\text{V6}}
\end{equation*}
predicate decisions, where $m_{\text{gen}}\le m$ counts the entries declared
general --- the only ones carrying a matrix whose Hermiticity is
decided; V6's remaining content is $m$ comparisons of a relation symbol, a tag
and a type, which are decided and not computed. \emph{Beside these stand the
summation decisions of condition~(S)}: one on $\tr\tau_A$ in trusted
mode and none on the absent marginal in untrusted mode; one on
$\sum_xM_x$; and one per sector on $\sum_vF_v^{(s)}$, so
$|\SecSet|+1+I_{\mathrm{tr}}$ of them,
each a pass over the additions of one declared-order sum. And V7, on the
candidate it verifies, takes $2|\SecSet|+1$ decisions --- Hermiticity and
semidefiniteness per block, and the trace --- one further in trusted mode for
the marginal equality, and one decision per constraint entry: semidefiniteness
for an inequality, exact entrywise equality for the relation $=$
(Theorem~\ref{thm:feasibility}\,(ii)); in untrusted mode a constraint leg is one
comparison of real numbers instead. Every one of them is taken on a matrix of
size $d_A$, $d_s$ or $d_Ad_s$ for a single $s$; \emph{the validator
assembles no array larger than one such matrix}, which where $|\SecSet|=1$ is
$d_{\Qsys}$ or $d_Ad_{\Qsys}$ square, the one block being the whole space
there (Theorem~\ref{thm:band-a}\,(i)).
In trusted mode the V2 term is $3$ and not $4$: D(ii)'s named Hermiticity
condition and the Hermiticity half of its positivity condition are the same
decision, taken once. In untrusted mode V2 decides only the marginal slot's
$\nomarg$ branch, with no marginal arithmetic or marginal summation predicate.
\end{enumerate}

\emph{The general format.} All four parts hold for a candidate declaration in
the general format of Definition~\ref{def:declaration}\,(DX1)--(DX6), with the
checks read as in the general-format paragraph of
Definition~\ref{def:validator} and the conditions as in that of
Definition~\ref{def:target}: (i) soundness on all eight and (ii) completeness
on D(i)--D(vi) and D(viii); (iii) with the \emph{same} witness, which is a
general-format declaration at $|\mathbb{T}|=1$, one measured factor and
no party-local entry; and (iv) with the predicate decisions of V2, V3 and V6
numbering
\begin{equation*}\tag{\getrefnumber{C-eq:predcount-general}}\digesteqlabel{eq:predcount-general}
  \underbrace{I_{\mathrm{tr}}\big(2|\mathbb{T}|+1
    +|\mathbb{T}|\,\mathbb 1[\text{ensemble form}]\big)}_{\text{V2}}
  \;+\;\underbrace{|\mathbb{T}|\,\big(2|\Xset|+1\big)}_{\text{V3 on }A}
  \;+\;\underbrace{\sum_{f\in\mathbb{F}}|\SecSet^{f}|
    \big(2|\Vset^{f}|+1\big)}_{\text{V3 per measured factor}}
  \;+\;\underbrace{\sum_{k=1}^{m}|\SecSet^{f_k}|}_{\text{V3 on }\CList}
  \;+\;\underbrace{m_{\text{gen}}\,|\mathbb{T}|}_{\text{V6}} ,
\end{equation*}
beside them the summation decisions of condition~(S) --- one on
$\tr\tau_A$ in trusted mode only, one per Alice sector on Alice's family,
and one per factor and sector on that factor's family, so
$|\mathbb{T}|+\sum_f|\SecSet^{f}|+I_{\mathrm{tr}}$ of them --- and V7, on the
candidate it verifies, $2|\mathbb{T}|\,|\SecSet|+1$ decisions, $|\mathbb{T}|$
further in trusted mode for the marginal equalities, and per entry $k$ either
$|\mathbb{T}|$ operator comparisons, one comparison of real numbers in
untrusted mode, or, for a party-local entry,
$|\mathbb{T}|\,|\SecSet|/|\SecSet^{f_k}|$ operator comparisons.
Each operator comparison is one semidefiniteness decision for an inequality
or one exact entrywise equality decision for $=$
(Theorem~\ref{thm:feasibility}\,(ii), general format). \emph{Every one of them is
taken on a matrix of size $d_{A_t}$, $d_s$ or $d_{A_t}d_s$ for a single joint
sector, or of size $d_{A_t}d_s/d_{f,s_f}$ for a single party-local row;
the validator assembles no array larger than one joint block.} In the
basic format --- $|\mathbb{T}|=1$, one measured factor, POVM form, two
messages neither reading the other and no party-local entry ---
\eqref{eq:predcount-general} is
\eqref{eq:predcount} term for term, and every other count of this paragraph is
the corresponding count of (iv).
\end{digestclaim}
\end{mainclaim}

\section*{7\quad Consuming the reduction chain}\pdfbookmark[1]{7 Consuming the reduction chain}{digest-group-protocol-compiler/def:chaincond}

% claim protocol-compiler/def:chaincond
\digestlabel{def:chaincond}{7.1}\pdfbookmark[2]{Definition 7.1 Interface conditions on a chain}{digest-protocol-compiler/def:chaincond}
\renewcommand*{\RestatedName}{\hyperref{contract.pdf}{}{claim-protocol-compiler/def:chaincond}{Definition~7.1}}%
\begin{digestclaim}[Interface conditions on a chain]
A reduction chain $\pcChain=(\Chsrc,\Chrec)$ presented to the compiler, with
$\Chsrc=(\Xi_1,\dots,\Xi_{L_{\mathrm s}})$ and
$\Chrec=(\Lambda_1,\dots,\Lambda_L)$, satisfies the
\emph{interface conditions on a chain}
when the following hold.

\emph{Receiver segment.}
\begin{enumerate}[label=(C\arabic*),leftmargin=3.0em]
\item \emph{(descriptor match)} the target descriptor of $\Lambda_i$ equals the
source descriptor of $\Lambda_{i+1}$ for $1\le i<L$; the source descriptor of
$\Lambda_1$ is the physical receiver space of $\Spec_{\mathrm{rec}}$; and the
target descriptor of $\Lambda_L$ is $(\Qsys,\Smap)$ --- that is, the last
link's target descriptor \emph{is} the sector map. That equality is
decided by the same four-component comparison V1 makes on the map of a
declared $\Qsys$-operator family (Definition~\ref{def:validator});
\item \emph{(outcome correspondence)} the declared outcome correspondence is a
tabulated total surjection $\gamma:\Vph\to\Vset$ from the \emph{physical}
alphabet onto the \emph{emitted} one --- the coarse-graining direction, of
which a bijection is the special case;
\item \emph{(target POVM)} $\{F_v\}_{v\in\Vset}$ is a family of
block lists against $\Smap$ --- \emph{typing}, held by the
emitted form and not checked --- and for each $s$ the block family
$\{F_v^{(s)}\}_v$ is a POVM on $\Qsys_s$, decided \emph{exactly} on the
emitted blocks themselves: $F_v^{(s)}$ Hermitian and positive semidefinite
for each $v$, $\sum_vF_v^{(s)}=\idon{\Qsys_s}$, and the declared-order sum
forming it rounding nowhere (Definition~\ref{def:exact}). This condition is
check V3 restricted to $\{F_v\}$ and states the same list of conjuncts
(Theorem~\ref{thm:chain-consumption}\,(ii));
\item \emph{(pullback, where checkable)} if every link of $\Chrec$ supplies its
channel as data then, with
$\Lambda\coloneqq\Lambda_L\circ\dots\circ\Lambda_1$,
\begin{equation*}\tag{\getrefnumber{C-eq:pullback}}\digesteqlabel{eq:pullback}
  \Lambda^{\dagger}(F_v)=\sum_{w\in\gamma^{-1}(v)}\Fph_w
  \qquad\text{for every }v\in\Vset ,
\end{equation*}
an \emph{exact} matrix identity, decided entry by entry on the emitted matrices
and on the supplied channel data, with the fibre sum accumulated in the
declared order of $\Vph$. Where some link supplies no channel as data,
only a theorem that constructs one, the identity is not decidable from the
emitted data and is carried instead as an obligation of the chain's
producer, Assumption~\ref{ass:link};
\item \emph{(slots)} each link carries its constraint entries and its licence
slot, both possibly empty, and its status tag.
\end{enumerate}

\emph{Source segment.} The source segment satisfies the
interface conditions (S1)--(S4) of
\xref{transmitter/def:source-segment} --- descriptor match with the
virtual transmitted system $\Avirt$ as the source descriptor of
$\Xi_1$, on which the certificate of \xref{transmitter/def:srcert} is
written, and every descriptor Alice-side; index preservation; the source
identity $\Xi(\rvirt_x)=\rho_x$ for every $x$, that definition's equation
(S3), where every link supplies its channel as data; and the slots ---
\emph{decided by this module on the presented data}, exactly as that
definition states them. When $L_{\mathrm s}=0$ the segment is empty, $\Xi$ is
the identity, $\Avirt=\Asig$ and $\rvirt_x=\rho_x$, so (S3) holds by
inspection.

\emph{Between the segments.}
\begin{enumerate}[label=(X\arabic*),leftmargin=3.0em]
\item \emph{(order)} the chain is composed source segment first and receiver
segment second. This is not a tag and is not checked: it is the shape of the
record. A source-map link has no slot in $\Chrec$ to be written into, and a
link presented in the wrong segment fails that segment's descriptor rule ---
(C1) for a source map, since $\Asig$ is not the physical receiver space, and
(S1) for a squashing map, since $\Qsys$ is not an Alice-side descriptor.
\end{enumerate}

\emph{At which reading (C4) and (S3) are decided.} Both are \emph{matrix
equalities on declared data} --- a composite channel supplied as
binary64 data, applied to binary64 POVM elements on one side and to binary64
signal states on the other --- and both are decided \emph{exactly}, entry by
entry, on the numbers the compiler holds. \emph{The composite, the adjoint
and the fibre sum are formed in exact arithmetic as well}, so the identity
decided is the identity written; that
the operands are binary64 is not what makes the decision exact, a product of
two binary64 numbers need not be one. Exactness here is an obligation of
the chain's producer (Assumption~\ref{ass:link}): a producer whose composite
lands off the identity by a rounding has not delivered the reduction the
interface asks for, and repairs its own emission; no window is opened here. Everything else
in this definition --- descriptors, alphabets, fibres, slots, segments --- is
combinatorial and is read exactly, as D(vi)'s relation symbols are.

\emph{The general format: one receiver segment per measured factor.}
For a chain presented as Definition~\ref{def:spec}\,(c$'$), with
$\Chrec=(\pcChain^{f}_{\mathrm{rec}})_{f\in\mathbb{F}}$, the conditions
(C1)--(C5) are decided \emph{of each segment separately, and of nothing
across two segments}. (C1): within segment $f$ the target descriptor of a link
is the source descriptor of the next, the source descriptor of
$\Lambda^{f}_1$ is factor $f$'s physical receiver system, and the target
descriptor of $\Lambda^{f}_{L_f}$ is
$(\Qsys_f,\SecSet^{f},(d_{f,s_f})_{s_f})$ --- \emph{that factor's own sector
map, and not the product map}, which is derived by
Definition~\ref{def:declaration}\,(DX1) and is the target descriptor of
nothing; the same four-component comparison decides it. (C2): one tabulated total surjection
$\gamma_f$ from factor $f$'s physical alphabet onto the alphabet factor $f$'s
last link delivers. (C3): the target family $\{F^{f}_{v_f}\}$ is a family of
block lists against factor $f$'s map and, for each $s_f$, a POVM on
$\Qsys_{f,s_f}$ decided exactly on that factor's blocks --- which is check V3
restricted to that factor, conjunct for conjunct, in its general-format
reading. (C4): where every link of segment $f$ supplies its channel as data,
with $\Lambda_f\coloneqq\Lambda^{f}_{L_f}\circ\dots\circ\Lambda^{f}_1$,
\begin{equation*}\tag{\getrefnumber{C-eq:pullback-factor}}\digesteqlabel{eq:pullback-factor}
  \Lambda_f^{\dagger}\big(F^{f}_{v_f}\big)
  =\sum_{w_f\in\gamma_f^{-1}(v_f)}F^{\mathrm{ph},f}_{w_f}
  \qquad\text{for every }v_f ,
\end{equation*}
an exact matrix identity on that factor's matrices, the fibre sum accumulated
in the declared order of that factor's physical alphabet. \emph{The product
composite $\bigotimes_f\Lambda_f$ and the product identity on $\{F_v\}$ are
formed and decided by no one}: the identity
$\big(\bigotimes_f\Lambda_f\big)^{\dagger}(F_v)
=\bigotimes_f\Lambda_f^{\dagger}(F^{f}_{v_f})$ is a consequence of
\eqref{eq:pullback-factor} and of the product rule
\eqref{eq:productfamily} --- it is \xref{reductions/thm:emitted-product}\,(ii),
named here and decided here by nothing --- and (C4) is the per-factor
identity and nothing wider. (C5) is per link, as above. Where some link of
segment $f$ supplies no channel as data, only a theorem that constructs one,
\eqref{eq:pullback-factor}
for \emph{that} segment is not decidable from the emitted data and is carried
as Assumption~\ref{ass:link}, exactly as above and per factor.

\emph{A non-empty source segment, and the conditions it brings.} \emph{A
presented chain may carry one.} The rule by which a non-empty source segment
composes with a receiver segment is a theorem of the reductions module
(\xref{reductions/lem:precomposition}, \xref{reductions/thm:two-sided}, with
\xref{reductions/def:chain}), in
trusted-source mode; and the module that admits source-link kinds may
emit an artefact carrying a link of an admitted kind. What follows is what
this module decides of such a segment, on the presented data and exactly;
whether a presented chain carries one is data of the record and is asserted
here of no artefact.
\begin{itemize}[leftmargin=2.6em]
\item \emph{(S1)--(S4) as \xref{transmitter/def:source-segment}
states them}, which is what the source-segment paragraph above already says;
and within (S3), \emph{the split by kind of witness}: for a link supplying its
channel as data the per-link identity is decided here, entry by entry, with
the composite formed in exact arithmetic; for a link supplying a
theorem constructor it is \emph{not decided here at all} but is the conclusion of the
claim that link's obligation record names, and the composite identity
is the composition of the per-link ones. So what this module decides of (S3)
is the data links' share of it, and the constructor links' share travels as
the obligation Assumption~\ref{ass:link} names ---
Theorem~\ref{thm:chain-consumption}\,(iv), (v) --- and is discharged nowhere in this
contract.
\item \emph{(A1)--(A3) of that definition, which are decided on data and are
not part of (S1)--(S4)}: that the triple a link's obligation record carries
--- its kind, the claim id it names, and the assumption short ids it names ---
appears \emph{entry by entry} in the soundness registry the artefact
carrying the segment emits, both components of that registry entry's
assumption list compared, so that a record naming a claim's closure and
omitting the source model's attestation fails as one naming neither
does; that every assumption of that entry is in the artefact's assumption
ledger; and that the segment is presented inside an artefact of the module
that admits its links. These are comparisons of finitely many declared
strings and of a finite list against a finite list, and \emph{they are decided
here because they are decidable on the presented record}, which is why they
are named beside (S1)--(S4) rather than left to Assumption~\ref{ass:link}.
\emph{What they do not decide} is that the named claim's conclusion applies to
these descriptors --- family membership is the emitting module's assertion and
no consumer re-proves a theorem --- and a link supplying its channel as data
carries no obligation record and is decided by (S3) instead.
\item \emph{the licensed-only rule.} Beside a non-empty source segment
\emph{every} entry of the emitted $\CList$ is required to carry a real scalar
as its right-hand side --- homogeneous, party-local, or an untrusted-mode
scalar (Definition~\ref{def:homog}).
\item \emph{the declared source mode.} \emph{A declaration carrying a
non-empty source segment is in trusted-source mode}: the rule fails
where the marginal slot of \eqref{eq:declaration} holds $\nomarg$
(Definition~\ref{def:declaration}), the untrusted-source mode, while
$L_{\mathrm s}\ge1$.
\end{itemize}
\emph{(X1) in the general format.} The order rule is unchanged and is still
the type of the record: a source-map link has no slot in any receiver segment
to be written into, and a link presented in the wrong list fails that list's
descriptor rule --- (C1) of the factor whose segment it was written into, or
(S1). In the reductions module the composite of a chain is a pair, and the
reason a source link may not follow a receiver link is a theorem of that
module (\xref{reductions/lem:precomposition},
\xref{reductions/thm:two-sided}) beside the published warrant
\cite[Rem.~9.2]{tupkary2026decoy}; this module types the record, decides the
descriptors, and states neither.
\end{digestclaim}

% claim protocol-compiler/thm:chain-consumption
\digestlabel{thm:chain-consumption}{7.5}\pdfbookmark[2]{Theorem 7.5 Decided identities and named obligations of a consumed chain}{digest-protocol-compiler/thm:chain-consumption}
\renewcommand*{\RestatedName}{\hyperref{contract.pdf}{}{claim-protocol-compiler/thm:chain-consumption}{Theorem~7.5}}%
\begin{digestclaim}[Decided identities and named obligations of a consumed chain]
Let a chain be presented to the compiler, with the
interface conditions of
Definition~\ref{def:chaincond}. Then:
\begin{enumerate}[label=(\roman*),leftmargin=2.2em]
\item (C1)--(C5) and (S1)--(S4) are decided by finite enumeration and finitely
many arithmetic operations on the emitted data, and a chain failing any of
them is refused. (X1) is not decided,
because it is not a condition on data: it is the type of the record;
\item (C3) is exactly the part of D(iii) concerning $\{F_v\}$, in both halves
--- typing against $\Smap$, and the per-block predicates and summation
condition of Definition~\ref{def:exact} --- hence a chain passing the interface
conditions yields an emitted $\{F_v\}$ for which Theorem~\ref{thm:measurements-and-statistics}
applies. Where the Bob-side weighting of Definition~\ref{def:weighted} is
applied after the chain (Definition~\ref{def:compile}), (C3) is decided on the family the chain delivered,
the emitted family is its weighting $N_{v,f}=P(f\,|\,t(v))F_v$, check V3
decides the emitted family itself (Theorem~\ref{thm:measurements-and-statistics}\,(i)), and
Lemma~\ref{lem:weighted}\,(i), (ii) is what says the weighting cannot break
what (C3) established; the identity of conjunct lists asserted here is between
(C3) and D(iii) on the chain's family;
\item \emph{(what a verified (C4) buys)} when the composite receiver channel
$\Lambda$ is given as data, (C4) is \emph{equivalent} to the statement that
measuring the emitted family $\{F_v\}$ after $\Lambda$ reproduces the
\emph{$\gamma$-coarse-grained} physical measurement statistics:
\begin{equation*}\tag{\getrefnumber{C-eq:pullbackstats}}\digesteqlabel{eq:pullbackstats}
  \tr[F_v\,\Lambda(\rho)]=\sum_{w\in\gamma^{-1}(v)}\tr[\Fph_w\,\rho]
\end{equation*}
for every $v\in\Vset$ and every operator $\rho$ on the physical receiver space.
Consistency with completeness holds \emph{under one hypothesis this contract
names and does not verify}: if the supplied channel data are trace-preserving,
so that $\Lambda^{\dagger}$ is unital, then summing \eqref{eq:pullback} over
$v$ returns completeness of the physical family and nothing more, whereas the
same identity read with a coarse-graining in the opposite direction --- a
total surjection $\gamma':\Vset\to\Vph$, with
$\Lambda^{\dagger}(F_v)=\Fph_{\gamma'(v)}$ --- does not in general, forcing
every physical element covered more than once by $\gamma'$ to vanish. That
trace preservation is
Assumption~\ref{ass:link}, not a check: nothing here evaluates it;
\item \emph{(what a verified (S3) buys)} when the composite source channel
$\Xi$ is given as data, (S3) of
\xref{transmitter/def:source-segment} --- the source identity,
that definition's equation, decided by this module on the presented data
(Definition~\ref{def:chaincond}) --- says that the virtual source, pushed
through $\Xi$, is the specified source: $\Xi(\rvirt_x)=\rho_x$ for every $x$,
hence
\begin{equation*}\tag{\getrefnumber{C-eq:srcresid}}\digesteqlabel{eq:srcresid}
  (\idmap_X\otimes\Xi)\Big(\sum_{x\in\Xset}p_x\dyad{x}_X\otimes\rvirt_x\Big)
  =\sum_{x\in\Xset}p_x\dyad{x}_X\otimes\rho_x ,
\end{equation*}
and, averaging, $\sum_xp_x\Xi(\rvirt_x)=\sum_xp_x\rho_x$: the
prepare-and-measure states of the emitted virtual description map onto the
specified ones, index by index, with the weights unchanged by (S2). The
compiler decides the source identity (S3) and nothing beyond it; and
\item the following are \emph{named and never discharged here}: that each link
of either segment is a sound reduction --- that its channel exists as claimed
for a link supplying only a theorem constructor, that a link supplying
its channel as data supplies a completely positive and trace-preserving map,
and that every reachable state after the reduction satisfies every transported
constraint entry --- and that any restricted attack set claimed by a
link is licensed. These are obligations of the chain's producer, the
soundness of each link being Assumption~\ref{ass:link}. No claim of this contract discharges them for a
presented chain in general --- for the Band A identity link, which
Definition~\ref{def:banda}\,(b) fixes with its channel supplied as data, its
outcome correspondence and its target POVM,
Theorem~\ref{thm:band-a}\,(v) discharges every clause of
Assumption~\ref{ass:link} by exhibition --- and
no proof below invokes them.
\end{enumerate}
\emph{The general format.} For a chain presented as
Definition~\ref{def:spec}\,(c$'$), with the interface conditions read as in the
general-format paragraph of Definition~\ref{def:chaincond}, (i)--(v) hold with
the following readings, of which the five above are the case of one
measured factor and an empty source segment.
\begin{enumerate}[label=(\roman*$'$),leftmargin=2.6em]
\item (C1)--(C5) are decided per measured factor and (S1)--(S4) as before, all
by finite enumeration and finitely many arithmetic operations, and a chain
failing any of them is refused.
\emph{And (A1)--(A3), the licensed-only rule and the declared-source-mode
rule are decided as well}: the first three by comparing finitely many declared
strings and one finite list against another, the fourth by reading $m$
declared tags, and the fifth by reading one declared slot. (X1) is again not
decided, being the type of the record.
\item (C3) is exactly the part of D(iii) concerning the declared measured
families, \emph{factor by factor}: for each $f$, (C3) read for $f$ and
D(iii)'s clauses on $\{F^{f}_{v_f}\}$ are the same list of conjuncts, so a
chain passing the interface conditions yields per-factor families for which
Theorem~\ref{thm:measurements-and-statistics}\,(i$'$) applies. \emph{That the product family is a POVM on
$\Qsys$ is on neither side of that pairing}, being a consequence of the
factors' conditions and no condition on data
(Definition~\ref{def:declaration}\,(DX3), Theorem~\ref{thm:validator-soundness}, general
format). The Bob-side weighting sentence applies per factor, to that factor's
delivered family.
\item \emph{(what a verified (C4) buys, per factor and jointly)} when every
link of segment $f$ supplies its channel as data,
\eqref{eq:pullback-factor} is equivalent to
$\tr[F^{f}_{v_f}\Lambda_f(\rho_f)]
=\sum_{w_f\in\gamma_f^{-1}(v_f)}\tr[F^{\mathrm{ph},f}_{w_f}\rho_f]$ for every
$v_f$ and every operator $\rho_f$ on factor $f$'s physical receiver space.
\emph{And the joint statement follows for every operator, product or not}: if
every segment supplies its channels as data then, with
$\Lambda\coloneqq\bigotimes_f\Lambda_f$, $\gamma\coloneqq\prod_f\gamma_f$ and
$F^{\mathrm{ph}}_w\coloneqq\bigotimes_fF^{\mathrm{ph},f}_{w_f}$,
\begin{equation*}\tag{\getrefnumber{C-eq:pullback-product}}\digesteqlabel{eq:pullback-product}
  \Lambda^{\dagger}(F_v)=\sum_{w\in\gamma^{-1}(v)}F^{\mathrm{ph}}_w ,
  \qquad\text{hence}\qquad
  \tr[F_v\,\Lambda(\rho)]=\sum_{w\in\gamma^{-1}(v)}\tr[F^{\mathrm{ph}}_w\,\rho]
\end{equation*}
for every $v\in\Vset$ and every operator $\rho$ on the product physical space.
\emph{Neither side of \eqref{eq:pullback-product} is formed or decided by
anything}: it is a consequence of the per-factor identities, and the objects in
it --- the product channel, the product physical family, the product target
family --- are data of no one. The consistency reading holds under the same
hypothesis read per factor: if every factor's supplied channel data are trace
preserving, so that each $\Lambda_f^{\dagger}$ is unital, then so is
$\Lambda^{\dagger}$ and summing \eqref{eq:pullback-product} over $v$ returns
completeness of the product physical family and nothing more. That trace
preservation is Assumption~\ref{ass:link}, per factor, and is evaluated
nowhere.
\item (S3) is as in (iv), with the split by kind of witness: what a
\emph{decided} (S3) buys is \eqref{eq:srcresid} and its averaged form, for the
data links' share of the composite; for a link supplying a theorem constructor
this module decides nothing and what its share buys is what the claim its
obligation record names concludes, which this contract neither states nor
asserts.
\item the list of obligations named and never discharged here gains three
items, each per the general format: that each link of \emph{each} receiver
segment is a sound reduction, factor by factor; that a non-empty source
segment composes with the receiver segments as a reduction, which is
\xref{reductions/thm:two-sided}, used by no step here; and that the claim an
obligation record names applies to that link's descriptors, which (A1)--(A3)
do \emph{not} decide. These are again the obligations of (v),
Assumption~\ref{ass:link} and the attack-set licence, and the emitting
module's own assertion; no claim of this contract discharges them, save as (v) says of
the Band A identity link, and no proof below invokes them.
\end{enumerate}
\end{digestclaim}

\section*{8\quad The Band-A route}\pdfbookmark[1]{8 The Band-A route}{digest-group-protocol-compiler/def:banda}

% claim protocol-compiler/def:banda
\digestlabel{def:banda}{8.1}\pdfbookmark[2]{Definition 8.1 Band-A specification}{digest-protocol-compiler/def:banda}
\renewcommand*{\RestatedName}{\hyperref{contract.pdf}{}{claim-protocol-compiler/def:banda}{Definition~8.1}}%
\begin{digestclaim}[Band-A specification]
A specification is
\emph{Band A}
when (a) its receiver part declares a POVM $\{\Fph_w\}_{w\in\Vph}$ directly on a
finite-dimensional system, with no characterization step and no photon-number
structure, and the emitted sector set is the singleton $\SecSet=\{s_0\}$ with
$\Pi_{s_0}=\idon{\Qsys}$; and (b) its reduction chain has an \emph{empty
source segment}, $L_{\mathrm s}=0$ --- the source segment of
\xref{transmitter/def:source-segment} at length zero, required here of a
Band-A chain and of no other --- and a receiver segment of one link, $L=1$,
$\Lambda_1=\idmap_{\Qsys}$ supplied as data, with status \emph{exact},
empty constraint slot and empty licence slot, whose outcome correspondence is
the identity map $\gamma=\idmap$ on $\Vset=\Vph$ and whose
target POVM is the declared family itself:
$F_v=\Fph_v$, entry for entry, for every $v\in\Vset$.
\end{digestclaim}

% claim protocol-compiler/thm:band-a
\digestlabel{thm:band-a}{8.2}\pdfbookmark[2]{Theorem 8.2 The Band-A route}{digest-protocol-compiler/thm:band-a}
\renewcommand*{\RestatedName}{\hyperref{contract.pdf}{}{claim-protocol-compiler/thm:band-a}{Theorem~8.2}}%
\begin{digestclaim}[The Band-A route]
For a Band-A specification:
\begin{enumerate}[label=(\roman*),leftmargin=2.2em]
\item the block structure of a Band-A declaration is the one-block structure:
$\Smap$ is the singleton map $\big((s_0,d_{s_0})\big)$, $d_{\Qsys}=d_{s_0}$,
and every $\Qsys$-operator is a one-element block list, so that its
single block \emph{is} the whole operator. Two consequences follow, and they
differ. The typing half of D(iii) is not vacuous but trivial to satisfy: an
operator is a one-element list, and V1's sector-map comparison is one
identifier and one dimension. The five predicates of D(iii) are \emph{not}
vacuous at all --- with $|\SecSet|=1$ the per-block predicates are the
full-space ones, decided once each on $d_{\Qsys}\times d_{\Qsys}$ matrices, and
a Band-A declaration is decided exactly like any other, condition~(S)
included;
\item the interface conditions of
Definition~\ref{def:chaincond} hold by inspection, save the predicates and
summation condition of (C3), which are decided: descriptors match because
source and target of $\idmap$ coincide; the outcome correspondence is the
identity bijection $\gamma=\idmap$ on $\Vset=\Vph$, as
Definition~\ref{def:banda}\,(b) requires; (C3)'s typing half holds
because the declared receiver family is already a one-element block list
against the singleton map, and its predicates and summation condition are the
declared receiver POVM's own, so (C3) holds exactly when that family is
exactly admissible, which is decided on it and is not given by
inspection; the pullback identity (C4) reads
$\idmap^{\dagger}(F_v)=F_v=\Fph_v$ and holds \emph{identically} --- entry for
entry, since $\idmap^{\dagger}=\idmap$, every fibre $\gamma^{-1}(v)$ is a
singleton, and the delivered $F_v$ \emph{is} the declared $\Fph_v$, the same
matrix, by Definition~\ref{def:banda}\,(b); the source conditions (S1)--(S4) of
\xref{transmitter/def:source-segment} are vacuous, quantifying over an empty
list of links --- the empty source segment being that definition's at
$L_{\mathrm s}=0$, which Definition~\ref{def:banda}\,(b) requires of a Band-A
chain and which is therefore a hypothesis of this theorem --- and the
virtual ensemble of (S3) is the specified one, so the source identity
reads $\rho_x=\rho_x$; the ordering (X1) holds because the source segment is
empty; and the slots of (C5) are present and empty. The one link of the
receiver segment is the identity link.
\emph{So both conditions that a general chain has to decide on computed data
are met here by exhibition}, and what a verified pullback identity and a
verified source identity buy on a general chain holds on this path without
appeal to any obligation of the chain's producer: measuring the
family the link delivers after $\idmap$ reproduces the physical statistics,
$\tr[F_v\,\idmap(\rho)]=\tr[\Fph_v\,\rho]$ for every $v$ and every operator
$\rho$, and $\idmap$ is exactly trace-preserving \emph{by exhibition}, the
supplied channel data being exactly that map rather than binary64 data
trusted to be trace-preserving, so its adjoint is unital and summing the
pullback identity over $v$ returns $\sum_v\Fph_v=\idon{}$, which on this path
is the same statement as (C3)'s own completeness, $\{\Fph_v\}$ being
$\{F_v\}$: the identity holds as a triviality and carries no content this
path did not already have;
\item the link transports no constraint entry, so $\CList$ is the
specification's own list $\Spec_{\mathrm{con}}$ and the constraint-soundness
obligation of the link (the transported-constraint clause of
Assumption~\ref{ass:link}) is vacuous; where $\Spec_{\mathrm{con}}$ is empty,
$\CList=\varnothing$ and the two-part gate of Theorem~\ref{thm:gates} is not
engaged;
\item consequently no emitted slot, and no claim of this contract used in the
Band-A path, depends on an artefact of the receiver module or of the
reductions module. The residual physical content of the receiver is
Assumption~\ref{ass:ideal}, and it is an idealization of the receiver, not a
reduction of it; an entry the specification declares itself engages the
gate, and with it the attack-set licence of Assumption~\ref{ass:atk}, as on
every route; and
\item \emph{(the identity link is sound)} the Band-A chain satisfies every
condition of Assumption~\ref{ass:link}: (i)(a) the channel of its one link,
$\idmap_{\Qsys}$, exists, the link supplying it as data; (i)(b) that link
satisfies the pullback identity \eqref{eq:pullback}; (i)(c) the protocol
using the measurement the link delivers is the protocol using the declared
receiver family $\{\Fph_v\}$, the delivered $\{F_v\}$ being that family, so
security of the one is security of the other, and where a setting-dependent
weighting $P(f\,|\,t)$ is applied after the link
(Definition~\ref{def:weighted}), the emitted family
$\{P(f\,|\,t(v))\,F_v\}$ is the identically weighted declared family
$\{P(f\,|\,t(v))\,\Fph_v\}$, so the same holds of the weighted protocol; (ii) its source segment is empty; (iii) its link transports no
constraint entry; and (iv) $\idmap_{\Qsys}$ is completely positive and
trace-preserving. No clause of Assumption~\ref{ass:link} is therefore a
premise on this path; that the declared family is the measurement the
receiver performs is Assumption~\ref{ass:ideal}, and no condition on the link.
\end{enumerate}
\end{digestclaim}

\section*{9\quad The emission gates}\pdfbookmark[1]{9 The emission gates}{digest-group-protocol-compiler/def:eligible}

% claim protocol-compiler/def:eligible
\digestlabel{def:eligible}{9.1}\pdfbookmark[2]{Definition 9.1 Key-length eligibility}{digest-protocol-compiler/def:eligible}
\renewcommand*{\RestatedName}{\hyperref{contract.pdf}{}{claim-protocol-compiler/def:eligible}{Definition~9.1}}%
\begin{digestclaim}[Key-length eligibility]
An emitted declaration is \emph{key-length-eligible}
when all of the following hold:
\begin{enumerate}[label=(G\arabic*),leftmargin=3.0em]
\item V1--V8 all pass --- whereupon D(i)--D(viii) all hold of the declared
numbers, by the soundness half of Theorem~\ref{thm:validator-soundness}\,(i);
\item \emph{(a consumption gate on the transmitter artefact)} if the
marginal slot is filled and $d_A\ge2$: the transmitter artefact
carries the source-replacement licence record of
\xref{transmitter/def:srcert} --- the certificate
$(\tau_A,\Psi,\{p_x,\rvirt_x\})$ with its attestation (SR3) --- whose
\emph{five decidable conditions} the transmitter decided exactly on the same
data (\xref{transmitter/thm:certificate-sufficiency}\,(i)) and which \emph{this
module re-decides on the emitted data}: $\Psi=\Psi^{\dagger}$; $\Psi^2=\Psi$;
$\tr\Psi=1$, the trace accumulated over the row index and
rounding nowhere; (SR1); and (SR2) --- each
\emph{exactly}, as matrix and scalar identities. The transmitter's decision is
not taken on its word. \emph{No condition of (G2) is stated on
a vector.} The certificate's declared datum is the operator $\Psi$; the unit
vector $\ket{\psi}$ with $\Psi=\dyad{\psi}$ exists by $\Psi^2=\Psi$ and
$\tr\Psi=1$ and is never itself declared (\xref{transmitter/def:srcert}).
Clause (SR3) carries no number and is not decidable
(Theorem~\ref{thm:gates}\,(i)): (G2) asks for its presence only. At $d_A=1$
no record is demanded, the licence being vacuous there;
\item if $m\ge1$: (a) $\CList$ lies in the licensed class of the emitted source
mode (Definition~\ref{def:homog}) --- \emph{a type check on each entry's
right-hand-side tag, comparing no numbers} --- \emph{and}
(b) the attack-set licence record, the attestation of
Assumption~\ref{ass:atk}, is present.
\end{enumerate}

\emph{(G2) by form of the certificate.} The five conditions above are those
of a certificate declared in pure form: a pure-form certificate is
re-decided exactly as stated, and refused when one of the five fails.

A certificate declared in block form is refused at this gate by its
declared form, whatever its blocks, by row~3 of the gate outcome table of
Definition~\ref{def:gate-outcome-table}. A block-form
certificate (\xref{transmitter/def:srcert}) is a positive semidefinite
unit-trace block list on which that definition imposes no purity, so
one with more than one non-zero block, or whose one non-zero block is not a
rank-one projector, also fails $\Psi^2=\Psi$; but the refusal does not rest on
that, since one whose one non-zero block is a rank-one projector can meet all
five conditions above.

An ensemble presentation declares no $\Psi$ and no $\{M_x\}$ to
re-decide. For it the gate decides, in place of the five conditions, what the
presentation declares: (E1), (E2) and (E3) of \xref{transmitter/def:srcert}
on the declared weighted ensemble, and that the emitted declaration is in
ensemble form with its Alice slots \emph{the} ones the presentation
defines --- the sector structure of the factor tagged $A$ is the
matched structure, and for every $t$ and $x$ the block $N^{(t)}_x$ is entry
for entry the transpose of the declared block $\sigma^{\mathrm v,(t)}_x$ and
$\tau^{(t)}_A$ the transpose of their declared-order sum. \emph{A presentation
accompanying a POVM-form declaration is refused}, nothing tying the declared
$\{M_x\}$ to it, which is the transmitter's own rule for a POVM-form consumer.
In every form, (G2) asks for the presence of (SR3).

\emph{(G3)(a) with the party-local form.} The licensed class of
Definition~\ref{def:homog} contains, in either source mode, the entries
declared party-local; the decision remains a type check on each
entry's tag and compares no numbers. A list carrying such an entry is
key-length-eligible only with the attack-set licence record of (b), as every
non-empty list is, and what that record attests for a tagged entry is a
certificate of the module that produces the chain, on the scoped factor's own
channel (Definition~\ref{def:homog}).
\end{digestclaim}

% claim protocol-compiler/def:gate-outcome-table
\digestlabel{def:gate-outcome-table}{9.3}\pdfbookmark[2]{Definition 9.3 The gate outcome table}{digest-protocol-compiler/def:gate-outcome-table}
\renewcommand*{\RestatedName}{\hyperref{contract.pdf}{}{claim-protocol-compiler/def:gate-outcome-table}{Definition~9.3}}%
\begin{digestclaim}[The gate outcome table]
The
\emph{gate outcome table}
is the following table of the failures of the gates of
Definition~\ref{def:eligible} and their outcomes; where the failures of one declaration meet the
premises of a refusal row and of an emission row both, the outcome is the
refusal.

\medskip
\noindent\emph{Gate failure to emitted outcome.} No row substitutes a default
value, drops an entry, changes the source mode or synthesizes an attestation.
\phantomsection

\noindent{\small
\begin{tabular}{@{}>{\raggedright\arraybackslash}p{0.58\textwidth}
                  >{\raggedright\arraybackslash}p{0.36\textwidth}@{}}
\hline
\emph{What failed} & \emph{Outcome} \\
\hline
(G1): any of V1--V8 & refusal, naming the check \\
(G2): trusted mode at $d_A\ge2$, certificate \emph{absent}, and
$m=0$ &
none (no outcome of the compiler) \\
(G2): trusted mode \emph{at any $d_A$}, certificate \emph{present}
but one
of the five conditions of Definition~\ref{def:eligible}\,(G2) fails ---
$\Psi$ Hermitian, $\Psi^2=\Psi$, $\tr\Psi=1$ rounding nowhere, (SR1), (SR2) &
refusal, naming the condition \\
(G2): trusted mode at $d_A\ge2$, certificate absent, and $m\ge1$ &
none (no outcome of the compiler) \\
(G3)(a): $\CList$ not in the licensed class --- some entry declared
general in trusted-source mode &
emission, not key-length-eligible \\
(G3)(b): the attack-set licence record absent &
emission, not key-length-eligible \\
\hline
\end{tabular}}

\noindent The cell of rows 2 and 4 assigns the compiler no outcome.
\end{digestclaim}

% claim protocol-compiler/thm:gates
\digestlabel{thm:gates}{9.4}\pdfbookmark[2]{Theorem 9.4 The emission gates and their outcomes}{digest-protocol-compiler/thm:gates}
\begin{mainclaim}
\renewcommand*{\RestatedName}{\hyperref{contract.pdf}{}{claim-protocol-compiler/thm:gates}{Theorem~9.4}}%
\begin{digestclaim}[The emission gates and their outcomes]
\begin{enumerate}[label=(\roman*),leftmargin=2.2em]
\item \emph{(decidability)} (G1) is decidable by Theorem~\ref{thm:validator-soundness}; (G2)'s
conditions --- the five of Definition~\ref{def:eligible}\,(G2), namely $\Psi$
Hermitian, $\Psi^2=\Psi$, $\tr\Psi=1$ with the trace rounding nowhere, (SR1)
and (SR2) --- are finitely many exact identities on the emitted data
(\xref{transmitter/thm:certificate-sufficiency}\,(i)); (G3)(a) is decided by
\emph{reading the type of the datum each entry's right-hand-side slot holds,
which is its tag}, and by nothing else --- it compares no matrix against any
other and constructs no right-hand side; (G3)(b) is a presence test. (G2)'s
clause (SR3) and (G3)(b)'s record are attestations and are \emph{not}
decidable; they are tested for presence only;
\item \emph{(outcome of a failure)} If any gate fails, the compiler emits no
key-length-eligible declaration, and which of the \emph{two} permitted
outcomes it produces --- refusal, or emission of a declaration that is not
key-length-eligible --- is fixed by the gate outcome table
(Definition~\ref{def:gate-outcome-table}) and is not an
implementer's choice; where the failures of one declaration meet the premises
of a refusal row and of an emission row both, the outcome is the refusal. No
protocol specification meets the premise of rows 2 and 4, the two rows of the
gate outcome table that describe a certificate-absent case: by
Definition~\ref{def:spec}\,(a) a transmitter artefact in trusted-source mode at
$d_A\ge2$ carries a source-replacement certificate. The compiler
never drops a constraint entry silently, never
substitutes a default, never changes the source mode it was handed, and never
synthesizes a licence.
\item \emph{(no assertion)} Passing (G3)(b) records that an attestation was
supplied. It asserts nothing about attack-set inclusion, and no claim of this
contract does. Likewise passing (G2) records that the transmitter's licence
record is present and that its five decided conditions hold of the emitted
data: it asserts nothing about the apparatus, and what it buys for the run is
bought under Assumption~\ref{ass:rhom} --- the certificate is decided once, for
the one declaration that describes every round, which is where this module
spends round homogeneity (\xref{transmitter/fact:sr} drops the per-round
superscript of the source state under exactly that assumption of the
consumer's).
\end{enumerate}

\emph{The general format.} For a declaration in the general format of
Definition~\ref{def:declaration}\,(DX1)--(DX6) the three parts hold with the
following readings, of which the three above are the case of one Alice sector,
one measured factor, POVM form, a pure-form certificate and no
party-local entry.
\begin{enumerate}[label=(\roman*$'$),leftmargin=2.6em]
\item \emph{(decidability)} (G1) is decidable by Theorem~\ref{thm:validator-soundness} in its
general-format reading. (G2) keeps its five conditions and is decidable in
\emph{one further declared form} (Definition~\ref{def:eligible}): the five
conditions of a pure-form certificate, as above, a certificate declared
in block form being refused at this gate; and, for an ensemble
presentation, (E1), (E2), (E3) on the declared weighted ensemble together with
the entry-for-entry identity of the declaration's Alice slots with the
transposes of the declared blocks and of their declared-order sum --- finitely
many exact identities, comparisons of labels and integers, and finitely many
summation decisions, the number fixed by $|\mathbb{T}|$, $|\Xset|$ and the
declared dimensions. \emph{Purity is decided by nothing there}: a presentation
declares no $\Psi$, and purity of its blocks, (SR1) and (SR2) are theorems of
the presentation (\xref{transmitter/lem:ensemble-presentation}\,(i), (ii)), the
second giving every non-zero block as a positive multiple of a rank-one
projector. \emph{At one Alice sector it is not thereby waived}, the defined
$\Psi$ being that one block, of unit trace, hence a rank-one projector of which
the five conditions hold. \emph{At two or more it is waived}, and the gate
rests there on two statements --- rank-one blocks by theorem, which is the
hypothesis under which \xref{transmitter/thm:certificate-sufficiency}\,(ii) holds in both directions,
and that theorem's (iii), which reads purity nowhere --- the first of which a
general block form does not supply. (G3)(a) is decided by reading the type of the
datum each entry's right-hand-side slot holds and, for the third form, the tag
field it arrives with --- three forms being in the licensed class, and still
\emph{nothing computed and nothing compared numerically}. (G3)(b) is a presence
test, and (SR3) is tested for presence in every form.
\item \emph{(no row is added)} The outcome of every failure of the general
format is fixed by a \emph{row the gate outcome table already has}. A
general-format check
of Definition~\ref{def:validator} that fails is (G1) and is row 1, the refusal
naming the check. A certificate that is \emph{present} and fails a condition
of (G2) is row 3: a certificate that does not verify is a defect in the
supplied data, not a weaker claim. \emph{Row 3's premise is read of the
declared form}, which reads the cell and adds no row: for a pure-form
certificate the five conditions the cell enumerates are the conditions, and
the cell is literal; for a certificate declared in block form the condition
it fails is the form itself, (G2) admitting no block-form certificate
(Definition~\ref{def:eligible}), and the cell's premise --- a certificate
present and failing a condition of (G2) --- covers it, whether or not it also
fails $\Psi^2=\Psi$; for an ensemble presentation the conditions are (E1) to (E3) and the
identity of the declaration's Alice slots with the transposed declared blocks,
and the cell's premise --- a certificate present and failing a condition of
(G2) --- covers them, the enumeration in it being the pure form's list and not
a restriction of the premise. \emph{An ensemble presentation accompanying a
POVM-form declaration is a failure of that kind and is refused} (row 3),
because nothing ties the declared
$\{M_x\}$ to the presentation (Definition~\ref{def:eligible}). A list
carrying an entry declared general in trusted-source mode is row 5 and is
emitted, not key-length-eligible; \emph{a list carrying a party-local
entry is not that row}, that form being licensed in either mode.
\item \emph{(no assertion)} Item (iii) holds as stated, and passing (G3)(b)
records, rather than asserts, one further item: for a list carrying a
party-local entry, what the attack-set licence record attests is a certificate
of the module that produces the chain, on the scoped factor's own channel
(Definition~\ref{def:homog}), and
\emph{no function of the emitted matrices decides it}, exactly as for the
untagged forms. Passing (G2) on an ensemble presentation likewise records that
the transmitter's record is present and that the conditions of \emph{that
form} hold of the emitted data, and asserts nothing about the apparatus; what
it buys for the run is bought under Assumption~\ref{ass:rhom} as in (iii), the
certificate being decided once for the one declaration that describes every
round.
\end{enumerate}
\end{digestclaim}
\end{mainclaim}

% claim protocol-compiler/prop:conservative
\digestlabel{prop:conservative}{9.6}\pdfbookmark[2]{Proposition 9.6 Dropping a trusted marginal enlarges the feasible set}{digest-protocol-compiler/prop:conservative}
\renewcommand*{\RestatedName}{\hyperref{contract.pdf}{}{claim-protocol-compiler/prop:conservative}{Proposition~9.6}}%
\begin{digestclaim}[Dropping a trusted marginal enlarges the feasible set]
Let a specification request trusted-source mode with marginal $\tau_A$
satisfying D(ii), and a licensed list
$\{(\hat O_k,\bowtie_k,c_k)\}_{k=1}^{m}$, $m\ge0$, satisfying D(iii) --- every
$\hat O_k$ a block list against $\Smap$ and Hermitian block by block,
which the proof spends through Lemma~\ref{C-lem:hermitian-constraint-values} ---
with every entry declared homogeneous, so that the trusted-mode
right-hand side it denotes is the constructed $c_k\tau_A$
(Definition~\ref{def:homog}). Let $\Feas^{\mathrm{tr}}$ be the
feasible set of the trusted-mode declaration and $\Feas^{\mathrm{un}}$
that of the untrusted-mode declaration obtained by replacing the marginal
slot with $\nomarg$ and reading each entry's right-hand side as the scalar
$c_k$, all other slots unchanged. Then
\begin{equation*}\tag{\getrefnumber{C-eq:inclusion}}\digesteqlabel{eq:inclusion}
  \Feas^{\mathrm{tr}}\subseteq\Feas^{\mathrm{un}} ,
\end{equation*}
and consequently, for \emph{any} functional $J$ on $\Lin{A\Qsys}$,
$\inf_{\Feas^{\mathrm{un}}}J\le\inf_{\Feas^{\mathrm{tr}}}J$. The inclusion is
\emph{uniform in $\tau_A$}: it holds for every $\tau_A$ satisfying D(ii), and
the untrusted set does not depend on which one the request carried.
\end{digestclaim}

\section*{10\quad Assumptions}\pdfbookmark[1]{10 Assumptions}{digest-group-protocol-compiler/ass:spec}

% claim protocol-compiler/ass:spec
\digestlabel{ass:spec}{10.1}\pdfbookmark[2]{Assumption 10.1 Specification faithful to the apparatus}{digest-protocol-compiler/ass:spec}
\SetShortId{A-SPEC}
\renewcommand*{\RestatedName}{\hyperref{contract.pdf}{}{claim-protocol-compiler/ass:spec}{Assumption~10.1}}%
\begin{digestassumption}[Specification faithful to the apparatus]
For a protocol specification $\Spec$ (Definition~\ref{def:spec}):
(i) the declared dimensions, alphabets and tables of $\Spec$ describe the
apparatus, and so do its declared measurements, namely: the physical
receiver family $\{\Fph_w\}_{w\in\Vph}$ (Definition~\ref{def:spec}\,(b); in
the general format, each measured factor's family of (b$'$)) is the
measurement performed, in every round, on that round's physical receiver
system alone, whatever happens in the other rounds, their outcomes and public
messages included; and
where no source-replacement certificate is demanded of $\Spec$ (in
untrusted-source mode, or at $d_A=1$), $\{M_x\}$
(Definition~\ref{def:spec}\,(a)) is the measurement Alice performs on $A$;
(ii) the key rule $r^{0}_{\hold}$ is valued in $\Zval$.
\end{digestassumption}

% claim protocol-compiler/ass:alice-meas
\digestlabel{ass:alice-meas}{10.2}\pdfbookmark[2]{Assumption 10.2 Alice's declared measurement where no source certificate is demanded}{digest-protocol-compiler/ass:alice-meas}
\SetShortId{A-ALICE-MEAS}
\renewcommand*{\RestatedName}{\hyperref{contract.pdf}{}{claim-protocol-compiler/ass:alice-meas}{Assumption~10.2}}%
\begin{digestassumption}[Alice's declared measurement where no source certificate is demanded]
(i) Where no source-replacement certificate is demanded of the
protocol specification (in untrusted-source mode, or at $d_A=1$), the
letter $x\in\Xset$ Alice obtains in every round is the outcome of the
declared measurement $\{M_x\}$ (Definition~\ref{def:spec}\,(a)) applied to
that round's system $A$ alone, the measurement she performs on $A$ being
$\{M_x\}$ whatever happens in the other rounds, their outcomes and public
messages included.
\end{digestassumption}

% claim protocol-compiler/ass:rhom
\digestlabel{ass:rhom}{10.3}\pdfbookmark[2]{Assumption 10.3 Round homogeneity of the declaration}{digest-protocol-compiler/ass:rhom}
\SetShortId{A-RHOM}
\renewcommand*{\RestatedName}{\hyperref{contract.pdf}{}{claim-protocol-compiler/ass:rhom}{Assumption~10.3}}%
\begin{digestassumption}[Round homogeneity of the declaration]
(i) Every round of the protocol is described by the same specification, hence
by the same declaration, the same source mode and the same operators: the
family of the specification, the declaration, the source mode and the
operators satisfies clause~(i), sameness, of
round homogeneity
(\xref{foundations/def:round-homogeneity}).
\end{digestassumption}

% claim protocol-compiler/ass:coin
\digestlabel{ass:coin}{10.4}\pdfbookmark[2]{Assumption 10.4 Public choices as private coins of declared bias}{digest-protocol-compiler/ass:coin}
\SetShortId{A-COIN}
\renewcommand*{\RestatedName}{\hyperref{contract.pdf}{}{claim-protocol-compiler/ass:coin}{Assumption~10.4}}%
\begin{digestassumption}[Public choices as private coins of declared bias]
(i) The declared public choices $t\in\Tpub$ are drawn from private random
sources with their declared weights $q_t$, mutually independent, whichever
party draws each, and, taken together, independent of the quantum systems and
of the adversary;
(ii) a public choice $f$ declared with a setting-dependent weighting
$P(f\,|\,t)$ --- the setting-dependent form of
Definition~\ref{def:weighted}, which may be declared at any
measured factor --- is drawn in each round by a party whose
readout records the setting component $t$ of the
outcome, from a private random source with the declared conditional
distribution $P(\,\cdot\,|\,t)$ given that round's setting $t$; given $t$,
$f$ is independent of the quantum systems, of the letter and every outcome of
the round (the device outcome $y$ included), of every other declared public
choice and of the adversary, taken together;
(iii) every such choice is announced, or withheld, exactly as the public
structure declares.
\end{digestassumption}

% claim protocol-compiler/ass:link
\digestlabel{ass:link}{10.6}\pdfbookmark[2]{Assumption 10.6 Soundness of every reduction link}{digest-protocol-compiler/ass:link}
\SetShortId{A-LINK}
\renewcommand*{\RestatedName}{\hyperref{contract.pdf}{}{claim-protocol-compiler/ass:link}{Assumption~10.6}}%
\begin{digestassumption}[Soundness of every reduction link]
Each link of either segment of the reduction chain is sound:
(i) for a \emph{receiver} link, (a) where the link supplies no channel as
data, only a theorem that constructs one, its channel exists as claimed;
(b) it satisfies the pullback identity \eqref{eq:pullback}; and (c) security
of the protocol using the emitted measurement implies security of the
protocol using the physical one;
(ii) for a \emph{source} link, (a) the channel of the source identity (S3) of
\xref{transmitter/def:source-segment} exists as claimed; and (b) security of
the protocol preparing the virtual ensemble implies security of the
protocol preparing the specified one;
(iii) for a link of either segment, every reachable state after the reduction
satisfies every constraint entry the link transports;
(iv) a link of either segment that supplies its channel \emph{as data}
supplies a completely positive and trace-preserving map.
\end{digestassumption}

% claim protocol-compiler/ass:atk
\digestlabel{ass:atk}{10.7}\pdfbookmark[2]{Assumption 10.7 Attack-set licence}{digest-protocol-compiler/ass:atk}
\SetShortId{A-ATK}
\renewcommand*{\RestatedName}{\hyperref{contract.pdf}{}{claim-protocol-compiler/ass:atk}{Assumption~10.7}}%
\begin{digestassumption}[Attack-set licence]
For a non-empty constraint list $\CList$ that is
licensed for the declared source mode (Definition~\ref{def:homog})
and backs a key-length claim:
(i) the reduction chain that produced the declaration certifies the
\emph{licensed reduction}: security of the physical protocol reduces to
security of the consuming engine's generated round model
(\xref{meat-core/def:roundmodel}) with attack
sequences restricted to the $\CList$-admissible ones, that is, every attack
available against the physical protocol whose registers are finite-dimensional
and whose first channel receives every transmitted system is
represented, after the reduction, by a sequence whose per-round channel $A_j$ lies, for every round $j$, in the
restricted set $\AttSet_{\CList}(E_{j-1},\Qsys E'_j)$ of that definition, with
no $A$ output, in trusted-source mode, and in its untrusted-source-mode
counterpart, with $A$ supplied by the adversary, in untrusted-source mode.
\end{digestassumption}

% claim protocol-compiler/ass:ideal
\digestlabel{ass:ideal}{10.8}\pdfbookmark[2]{Assumption 10.8 Ideal qubit receiver on the Band-A route}{digest-protocol-compiler/ass:ideal}
\SetShortId{A-IDEAL}
\renewcommand*{\RestatedName}{\hyperref{contract.pdf}{}{claim-protocol-compiler/ass:ideal}{Assumption~10.8}}%
\begin{digestassumption}[Ideal qubit receiver on the Band-A route]
In a Band A specification:
(i) the declared receiver is exactly an ideal projective qubit measurement;
(ii) it has unit detection efficiency, no dark counts, no loss, no
multi-photon events and no misalignment;
(iii) the transmitted system is exactly a qubit.
\end{digestassumption}

\part*{Engine requirements}\pdfbookmark[0]{Engine requirements}{digest-part-2}

\section*{11\quad Requirements: the emitted declaration}\pdfbookmark[1]{11 Requirements: the emitted declaration}{digest-group-protocol-compiler/req:declaration-keys}

% claim protocol-compiler/req:declaration-keys
\digestlabel{req:declaration-keys}{11.1}\pdfbookmark[2]{Requirement 11.1 Machine keys of the emitted declaration}{digest-protocol-compiler/req:declaration-keys}
\renewcommand*{\RestatedName}{\hyperref{contract.pdf}{}{claim-protocol-compiler/req:declaration-keys}{Requirement~11.1}}%
\begin{digestrequirement}[Machine keys of the emitted declaration]
The engine emits the sector map $\Smap$ of a
protocol declaration as three machine keys,
\iface{sector_index_set} for its ordered block identifiers,
\iface{sector_dims} for their dimensions and \iface{sector_block_kind} for
its block kind, so that a consumer comparing sector maps can compare all
three and not two of them. It emits the label-to-index map $\idxm$ of
Definition~\ref{def:compile} alongside the declaration, under
\iface{label_index_map}, and with each constraint entry the tag of its
right-hand side, under \iface{constraint_rhs_tag}, so that the consuming
engine's type check has a datum (Definition~\ref{def:homog}). The sifting
predicate $\varsigma$ and the set $\Chkey$ of key-live values are not
emitted.
\end{digestrequirement}

% claim protocol-compiler/req:operator-form
\digestlabel{req:operator-form}{11.2}\pdfbookmark[2]{Requirement 11.2 Operators on the measured system held as block lists}{digest-protocol-compiler/req:operator-form}
\renewcommand*{\RestatedName}{\hyperref{contract.pdf}{}{claim-protocol-compiler/req:operator-form}{Requirement~11.2}}%
\begin{digestrequirement}[Operators on the measured system held as block lists]
The engine holds and emits every operator on $\Qsys$ as its
block list against the sector map $\Smap$
(Definition~\ref{def:declaration}). Where $|\SecSet|>1$ it does not
materialize the direct sum $\bigoplus_sF^{(s)}$ of such a list, with one
exception: condition (C4) of Definition~\ref{def:chaincond} forms it as the
argument of that one decision, decides it and discards it. The convention
\iface{operator_form} states this rule and its scope.
\end{digestrequirement}

% claim protocol-compiler/req:pass-through
\digestlabel{req:pass-through}{11.3}\pdfbookmark[2]{Requirement 11.3 Declared data passed through unchanged}{digest-protocol-compiler/req:pass-through}
\renewcommand*{\RestatedName}{\hyperref{contract.pdf}{}{claim-protocol-compiler/req:pass-through}{Requirement~11.3}}%
\begin{digestrequirement}[Declared data passed through unchanged]
The engine passes through unchanged the declared data that the compilation map
reads into a declaration: the transmitter artefact's marginal slot, the
value $[1]$ the transmitter filled it with at $d_A=1$ included; Alice's
measurement slot, in the form the artefact declares, which it changes only by
weighting it in that form (Definition~\ref{def:weighted}) where a choice
family is presented; the transcript; and each
constraint entry's right-hand-side datum and, where the entry carries one, its
scope. It fills no slot and demands no certificate of its own.
It constructs no right-hand side of an emitted constraint entry: it forms no
$c_k\tau_A$ for one and reads no $c_k$ off any matrix. It converts neither
form of Alice's measurement slot into the other: it forms neither
\eqref{eq:ensemblepovm} from an ensemble slot nor
$\tau_A^{1/2}M_x\tau_A^{1/2}$ from a POVM slot.
\end{digestrequirement}

% claim protocol-compiler/req:derived-objects
\digestlabel{req:derived-objects}{11.4}\pdfbookmark[2]{Requirement 11.4 Per-factor decisions on the declared blocks}{digest-protocol-compiler/req:derived-objects}
\renewcommand*{\RestatedName}{\hyperref{contract.pdf}{}{claim-protocol-compiler/req:derived-objects}{Requirement~11.4}}%
\begin{digestrequirement}[Per-factor decisions on the declared blocks]
For a declaration in the general format of
Definition~\ref{def:declaration}: no entry of a product family
\eqref{eq:productfamily} is formed as a datum, rounded or decided on, here or
downstream; the lifted operator of a scoped entry is never formed as a
datum; and the family \eqref{eq:ensemblepovm} that an ensemble form
denotes is formed by no one and decided by nothing. Every exact predicate is
decided per measured factor, on that factor's blocks and in that
factor's declared order, and on Alice's data per Alice sector
(Definition~\ref{def:exact}). The engine refuses by type a presentation of a
choice-weighted family indexed by a product alphabet
(Definition~\ref{def:weighted}).
\end{digestrequirement}

% claim protocol-compiler/req:tag-reading
\digestlabel{req:tag-reading}{11.5}\pdfbookmark[2]{Requirement 11.5 Reading the tag of a constraint entry}{digest-protocol-compiler/req:tag-reading}
\renewcommand*{\RestatedName}{\hyperref{contract.pdf}{}{claim-protocol-compiler/req:tag-reading}{Requirement~11.5}}%
\begin{digestrequirement}[Reading the tag of a constraint entry]
The engine reads the tag of an entry of one of the two untagged forms of
Definition~\ref{def:homog} off the type of its right-hand-side datum. Where
such an entry arrives carrying a tag field, check V6 of
Definition~\ref{def:validator} compares that field with the type of the
datum it names and refuses a disagreement; the tag is never taken on the
upstream module's word against the object the entry holds. An entry is read
as party-local exactly when it arrives carrying that tag; check V6
then requires its right-hand-side slot to hold a real scalar and its scope to
name a measured factor of the emitted sector map, and refuses otherwise. No
tag is inferred: an entry arriving without the party-local tag is never read
party-locally, whatever module emitted it.
\end{digestrequirement}

% claim protocol-compiler/req:no-window
\digestlabel{req:no-window}{11.6}\pdfbookmark[2]{Requirement 11.6 Exact decisions without a tolerance}{digest-protocol-compiler/req:no-window}
\renewcommand*{\RestatedName}{\hyperref{contract.pdf}{}{claim-protocol-compiler/req:no-window}{Requirement~11.6}}%
\begin{digestrequirement}[Exact decisions without a tolerance]
The engine declares no numerical tolerance, and there is no window in this
contract at which anything could be decided: no check of the validator
(Definition~\ref{def:validator}) reads a window or is decided at one, and
licensing (Definition~\ref{def:homog}) and the predicates of
Definition~\ref{def:exact} are decided exactly.
\end{digestrequirement}

% claim protocol-compiler/req:source-licence-passed-on
\digestlabel{req:source-licence-passed-on}{11.7}\pdfbookmark[2]{Requirement 11.7 The source-replacement licence passed on with the declaration}{digest-protocol-compiler/req:source-licence-passed-on}
\renewcommand*{\RestatedName}{\hyperref{contract.pdf}{}{claim-protocol-compiler/req:source-licence-passed-on}{Requirement~11.7}}%
\begin{digestrequirement}[The source-replacement licence passed on with the declaration]
The engine passes the transmitter's source-replacement licence record
\iface{source_replacement_license} (\xref{transmitter/def:srcert}) on to the
consuming engine with a declaration exactly when it emits that declaration in
trusted-source mode at $d_A\ge2$, after gate (G2) of
Definition~\ref{def:eligible} has re-decided its conditions on the emitted
data; the certificate and the attestation (SR3) it carries are the ones
supplied, and the engine asserts no attestation of its own
(\xref{transmitter/req:attestation-record}). At $d_A=1$ the record is neither
demanded nor emitted (\xref{transmitter/req:licence-scope}).
\end{digestrequirement}

% claim protocol-compiler/req:da1-source-identity
\digestlabel{req:da1-source-identity}{11.8}\pdfbookmark[2]{Requirement 11.8 Letter-independent emission at a one-dimensional trusted A}{digest-protocol-compiler/req:da1-source-identity}
\renewcommand*{\RestatedName}{\hyperref{contract.pdf}{}{claim-protocol-compiler/req:da1-source-identity}{Requirement~11.8}}%
\begin{digestrequirement}[Letter-independent emission at a one-dimensional trusted A]
For a specification in trusted-source mode at $d_A=1$, the engine
decides, for every $x\in\Xset$ in the declared order and exactly on the
declared numbers, the identity
\begin{equation*}
  \sigma_x \;=\; m_x\sum_{y\in\Xset}\sigma_y ,
\end{equation*}
where $m_x$ is the one entry of Alice's measurement element $M_x$, read as the
declaration carries it (Definition~\ref{def:compile}: for a
choice-weighted family the weighted element the engine forms; in the
ensemble form, $N_x$, which equals $M_x$ at $d_A=1$ by
\eqref{eq:ensemblepovm}), and $\sigma_x=p_x\rvirt_x$ is the weighted member
for $x$ of the virtual ensemble the transmitter artefact emits, which
at an empty source segment is the specified one
(\xref{transmitter/def:source-segment}); where the artefact emits the
ensemble by its weighted members, $\sigma_x$ is the declared member. Products
and sums are formed in exact arithmetic, entry by entry. A specification for
which the identity fails at some letter is refused, before a declaration is
emitted, with \iface{da1_source_identity_failed}, naming the first such letter
in the declared order and the first entry at which the two sides differ.
\end{digestrequirement}

% claim protocol-compiler/req:choice-weighting-record
\digestlabel{req:choice-weighting-record}{11.9}\pdfbookmark[2]{Requirement 11.9 Emission of the weighted families and the weighting record}{digest-protocol-compiler/req:choice-weighting-record}
\renewcommand*{\RestatedName}{\hyperref{contract.pdf}{}{claim-protocol-compiler/req:choice-weighting-record}{Requirement~11.9}}%
\begin{digestrequirement}[Emission of the weighted families and the weighting record]
The engine emits both forms of the choice-weighted family of
Definition~\ref{def:weighted}, each where it is the emitted family: the
setting-independent form \eqref{eq:weighted} on either party's side,
and the setting-dependent form \eqref{eq:weighted-t} at a measured
system, of whichever party, after that system's receiver segment. A
setting-independent presentation of a measuring party's physical receiver
family is the emitted family only where that party's receiver segment
delivers it as its target POVM; in front of any other
segment it is handed to the chain as the physical receiver family of
Definition~\ref{def:spec}\,(b), and the emitted family is the chain's. For the
setting-dependent relabelling it emits no link and no register
for the choice value, and the chain record it consumes and re-emits is
unchanged by the weighting; the weightings are recorded beside the chain in
the emitted artefact, in \iface{choice_weighting_record}, one entry per
measured factor at which a weighting is applied, each entry naming
that factor; like the setting-independent choice-weighted
family, they are this module's own construction.
\end{digestrequirement}

\section*{12\quad Requirements: choice weighting}\pdfbookmark[1]{12 Requirements: choice weighting}{digest-group-protocol-compiler/req:conditional-family-decision}

% claim protocol-compiler/req:conditional-family-decision
\digestlabel{req:conditional-family-decision}{12.1}\pdfbookmark[2]{Requirement 12.1 Decision on each declared conditional family}{digest-protocol-compiler/req:conditional-family-decision}
\renewcommand*{\RestatedName}{\hyperref{contract.pdf}{}{claim-protocol-compiler/req:conditional-family-decision}{Requirement~12.1}}%
\begin{digestrequirement}[Decision on each declared conditional family]
Wherever a choice-weighted presentation is read, on either party's side, the
engine decides, of each choice value $t\in\Tpub$ separately and on the
declared numbers, the condition Definition~\ref{def:weighted} places on the
conditional family $\{N_{u|t}\}_u$ --- in POVM form, its exact
admissibility in the reading of that definition; in ensemble form, the
condition that definition states in its place: (H) and (P) per block, the
ensemble identity in place of (C), and (T) --- before it forms any product and
before any condition of Definition~\ref{def:chaincond} is decided
on the chain: wherever the presented physical family is the weighted one,
that decision and the representability decision on the products it permits
precede the consumption of the chain, so that
no condition on the chain is decided on a product that was not decided
first. A failure is refused with
\iface{conditional_family_not_admissible}, naming the declaring party, and
the choice value, the predicate, the block and the index of the \emph{first}
failure --- for (S), the index of the first inexact addition; in ensemble
form the predicate named is the ensemble identity where that is the one that
failed. The decision stops at the first failure and does not enumerate the
rest.
\end{digestrequirement}

% claim protocol-compiler/req:binary64-weights
\digestlabel{req:binary64-weights}{12.2}\pdfbookmark[2]{Requirement 12.2 Declared weights as binary64 input}{digest-protocol-compiler/req:binary64-weights}
\renewcommand*{\RestatedName}{\hyperref{contract.pdf}{}{claim-protocol-compiler/req:binary64-weights}{Requirement~12.2}}%
\begin{digestrequirement}[Declared weights as binary64 input]
The weights of either form of Definition~\ref{def:weighted} --- each $q_t$ of
the setting-independent form and each entry of $P$ of the
setting-dependent form --- are binary64 numbers. The engine
rejects as malformed input a declared weight, in either form, that is not
exactly representable as a finite binary64 number, before that weight is used
in an accumulation or a product. This is the declared input type: the
rejection is not a failed summation predicate, and it adds no row of the gate
report and no refusal string.
\end{digestrequirement}

% claim protocol-compiler/req:weight-preparation
\digestlabel{req:weight-preparation}{12.3}\pdfbookmark[2]{Requirement 12.3 Preparation of user weights before the module boundary}{digest-protocol-compiler/req:weight-preparation}
\renewcommand*{\RestatedName}{\hyperref{contract.pdf}{}{claim-protocol-compiler/req:weight-preparation}{Requirement~12.3}}%
\begin{digestrequirement}[Preparation of user weights before the module boundary]
Ordinary decimal and rational user inputs for the weights of
Definition~\ref{def:weighted} may be converted and prepared by the declaring
entry point before they cross the module boundary into this module. An
entry point that does so records the supplied and the prepared values and
their displacement,
preserves normalisation and every required bound direction, and leaves this
module's exact tests unchanged. The preparation asserts no bound on a
resulting change of the key rate.
\end{digestrequirement}

% claim protocol-compiler/req:setting-only-weighting
\digestlabel{req:setting-only-weighting}{12.4}\pdfbookmark[2]{Requirement 12.4 Setting-dependent weighting read through the setting only}{digest-protocol-compiler/req:setting-only-weighting}
\renewcommand*{\RestatedName}{\hyperref{contract.pdf}{}{claim-protocol-compiler/req:setting-only-weighting}{Requirement~12.4}}%
\begin{digestrequirement}[Setting-dependent weighting read through the setting only]
The engine refuses by type, with \iface{choice_weight_depends_on_outcome}, a
weighting of the setting-dependent form of
Definition~\ref{def:weighted} that is presented indexed by the outcome
alphabet $\Vset$, by the device outcome, or by any component of the outcome
other than its setting; in the general format, one presented indexed by a
measured factor's whole outcome, by another factor's outcome, or by any
component other than that factor's setting.
\end{digestrequirement}

% claim protocol-compiler/req:weight-list
\digestlabel{req:weight-list}{12.5}\pdfbookmark[2]{Requirement 12.5 Decision on the declared weight list}{digest-protocol-compiler/req:weight-list}
\renewcommand*{\RestatedName}{\hyperref{contract.pdf}{}{claim-protocol-compiler/req:weight-list}{Requirement~12.5}}%
\begin{digestrequirement}[Decision on the declared weight list]
The engine decides the conditions of Definition~\ref{def:weighted} on the
declared weights $(q_t)_{t\in\Tpub}$ of a choice-weighted presentation: each
$q_t>0$ by its sign, and $\sum_tq_t=1$ exactly, with the sum accumulated in
the declared order of $\Tpub$ --- the order in which the declared
public-choice record lists it (\xref{foundations/req:declared-order}) --- and
required to round nowhere. A declared weight that is
zero or negative is refused by type. A weight list whose exact sum is not $1$,
or whose declared-order accumulation rounds at some addition, is refused with
\iface{choice_weights_not_normalized}, and no declaration is emitted.
\end{digestrequirement}

% claim protocol-compiler/req:stochastic-weights
\digestlabel{req:stochastic-weights}{12.6}\pdfbookmark[2]{Requirement 12.6 Stochastic matrix of a setting-dependent weighting}{digest-protocol-compiler/req:stochastic-weights}
\renewcommand*{\RestatedName}{\hyperref{contract.pdf}{}{claim-protocol-compiler/req:stochastic-weights}{Requirement~12.6}}%
\begin{digestrequirement}[Stochastic matrix of a setting-dependent weighting]
The engine decides the conditions of Definition~\ref{def:weighted} on the
declared matrix $P$ of the setting-dependent form: every entry
$P(f\,|\,t)\ge0$ by its sign, and every row sum $\sum_fP(f\,|\,t)=1$ exactly,
accumulated in the declared order of the choice alphabet
(\xref{foundations/req:declared-order}) and required to
round nowhere; the support of the weighting is
decided by the sign of each entry. A negative entry, or a row whose
declared-order sum is not exactly $1$, is refused with
\iface{choice_weight_not_stochastic}.
\end{digestrequirement}

% claim protocol-compiler/req:weighted-family-emission
\digestlabel{req:weighted-family-emission}{12.7}\pdfbookmark[2]{Requirement 12.7 Emission decisions on the weighted products}{digest-protocol-compiler/req:weighted-family-emission}
\renewcommand*{\RestatedName}{\hyperref{contract.pdf}{}{claim-protocol-compiler/req:weighted-family-emission}{Requirement~12.7}}%
\begin{digestrequirement}[Emission decisions on the weighted products]
The engine forms the products $q_tN_{u|t}$ and $P(f\,|\,t(v))F_v$ of
Definition~\ref{def:weighted} exactly, on the declared dyadic numbers, and
takes the two emission decisions of that definition before emitting anything.
The first, that every entry of every product it forms is a binary64 number,
is taken on every product the map forms, whether the product is emitted or
handed on: the intermediate families of an iterated presentation are
products, and a presentation on Bob's side before the chain produces the
physical family that (C2) and (C4) of Definition~\ref{def:chaincond} read and
that, in Band A, (C3) decides, so no condition of the chain is decided on a
product whose entries were not decided first. The second is that the emitted
family satisfies Definition~\ref{def:exact}. If either decision fails, the
engine refuses --- with \iface{emitted_value_not_representable} in the first
case and with the predicate that failed in the second --- and emits no
declaration. \iface{emitted_value_not_representable} covers exactly an exact
product of declared numbers that is not a binary64 number: an entry that is
binary64 but lies off whatever lattice the instance declared its own
data on is emitted, and the five conditions the consumer states are decided
on it exactly and on it alone.
\end{digestrequirement}

% claim protocol-compiler/req:no-lattice-exponent
\digestlabel{req:no-lattice-exponent}{12.8}\pdfbookmark[2]{Requirement 12.8 No lattice exponent and no snapping of the emission}{digest-protocol-compiler/req:no-lattice-exponent}
\renewcommand*{\RestatedName}{\hyperref{contract.pdf}{}{claim-protocol-compiler/req:no-lattice-exponent}{Requirement~12.8}}%
\begin{digestrequirement}[No lattice exponent and no snapping of the emission]
The engine declares no lattice exponent for the families it emits. It
establishes the five conditions of Definition~\ref{def:exact} of those
families directly, on the numbers it is about to emit, and refuses where one
fails (Requirement~\ref{req:weighted-family-emission}); it neither rounds nor
snaps its own emission, and it never declares a nearby family
(\iface{no_snapping}). A specification whose weighting does not land exactly
on the lattice is a specification to be re-declared at the instance.
What the engine owes its consumer is the five conditions, and it decides them
rather than inheriting them from a lattice.
\end{digestrequirement}

% claim protocol-compiler/req:no-fixed-fraction
\digestlabel{req:no-fixed-fraction}{12.9}\pdfbookmark[2]{Requirement 12.9 No fixed fraction outside the measurement operators}{digest-protocol-compiler/req:no-fixed-fraction}
\renewcommand*{\RestatedName}{\hyperref{contract.pdf}{}{claim-protocol-compiler/req:no-fixed-fraction}{Requirement~12.9}}%
\begin{digestrequirement}[No fixed fraction outside the measurement operators]
The emitted declaration carries the output-state test-statistics model
(\xref{reductions/ass:tsm}) and never the fixed-fraction model: no path of the
engine produces an artefact carrying a fixed-fraction model, and no weight
$q_t$ is emitted outside a POVM element. A specification requesting the
fixed-fraction encoding of Theorem~\ref{thm:choice-encoding}\,(iv) --- the conditional
families together with a separate state-independent fraction $q_t$ to be
applied outside the operators --- is refused with
\iface{fixed_fraction_requested}.
\end{digestrequirement}

\section*{13\quad Requirements: validation, feasibility and the chain}\pdfbookmark[1]{13 Requirements: validation, feasibility and the chain}{digest-group-protocol-compiler/req:validator-refusals}

% claim protocol-compiler/req:validator-refusals
\digestlabel{req:validator-refusals}{13.1}\pdfbookmark[2]{Requirement 13.1 Refusals of the validator's checks}{digest-protocol-compiler/req:validator-refusals}
\renewcommand*{\RestatedName}{\hyperref{contract.pdf}{}{claim-protocol-compiler/req:validator-refusals}{Requirement~13.1}}%
\begin{digestrequirement}[Refusals of the validator's checks]
The engine runs checks V1 to V8 of Definition~\ref{def:validator} on every
candidate declaration and refuses a candidate that fails one, emitting no
declaration, with a reason naming the failing check. It repairs nothing and completes no declared table with a default
value. Where a check has a reason of its own, the refusal carries that reason
and what the check names with it:
\begin{enumerate}[label=(\roman*),leftmargin=2.2em]
\item V1: \iface{sector_map_mismatch}, naming the first position at which the
  compared lists differ, and in the general format, for the block kind across
  the factors, the first factor at which the kind differs;
\item V2: \iface{alice_marginal_trace_rounds} where an addition of the
  marginal's trace rounds; in the general format
  \iface{alice_sector_marginal_zero}, naming the sector, where a block of the
  marginal is zero beside a measurement slot in ensemble form;
\item V3: \iface{alice_povm_summation_rounds}, in the general format carrying
  the Alice sector; \iface{bob_povm_summation_rounds}, carrying the sector and
  the outcome index at which the first inexact addition occurred, and in the
  general format the factor as well; in the general format
  \iface{alice_ensemble_marginal_mismatch}, carrying the sector and the first
  differing entry, \iface{alice_ensemble_sum_rounds}, carrying the sector and
  the index of the first inexact addition, and
  \iface{ensemble_form_in_untrusted_mode} for a slot in ensemble form beside a
  marginal slot holding $\nomarg$;
\item V6: in the general format \iface{constraint_scope_not_a_factor},
  naming the entry;
\item V7: \iface{no_feasibility_witness}, naming, where a supplied witness was
  rejected on its map, the first position at which the two lists differ.
\end{enumerate}
\end{digestrequirement}

% claim protocol-compiler/req:validation-transcript
\digestlabel{req:validation-transcript}{13.2}\pdfbookmark[2]{Requirement 13.2 Transcript of a check that stops early}{digest-protocol-compiler/req:validation-transcript}
\renewcommand*{\RestatedName}{\hyperref{contract.pdf}{}{claim-protocol-compiler/req:validation-transcript}{Requirement~13.2}}%
\begin{digestrequirement}[Transcript of a check that stops early]
A check of Definition~\ref{def:validator} may stop after a failed
prerequisite or predicate. Its transcript reports only the decisions the check
actually took, in their execution order, and the engine infers no outcome for
a predicate it did not execute.
\end{digestrequirement}

% claim protocol-compiler/req:feasibility-witness
\digestlabel{req:feasibility-witness}{13.3}\pdfbookmark[2]{Requirement 13.3 Default candidate before a supplied witness}{digest-protocol-compiler/req:feasibility-witness}
\renewcommand*{\RestatedName}{\hyperref{contract.pdf}{}{claim-protocol-compiler/req:feasibility-witness}{Requirement~13.3}}%
\begin{digestrequirement}[Default candidate before a supplied witness]
Where the constraint list is non-empty, on a candidate declaration
satisfying D(i)--D(iii) and D(vi), which checks V1--V3 and V6 of
Definition~\ref{def:validator} decide (a candidate failing one of them is
refused there), the engine decides the
conditions of
\eqref{eq:feas}, in the general format those of \eqref{eq:feas-general},
first on the default candidate \eqref{eq:witness-empty}, in the general format
\eqref{eq:witness-general}, and, if it fails, on a
supplied witness $\omega^{\star}$, read from the input
\iface{feasibility_witness}. It forms the default candidate only to test it,
and it emits neither candidate. It refuses with \iface{no_feasibility_witness}
exactly when the default candidate fails and either no witness is supplied or
the supplied one fails.
\end{digestrequirement}

% claim protocol-compiler/req:witness-refusal-meaning
\digestlabel{req:witness-refusal-meaning}{13.4}\pdfbookmark[2]{Requirement 13.4 Reporting a refusal for want of a feasibility witness}{digest-protocol-compiler/req:witness-refusal-meaning}
\renewcommand*{\RestatedName}{\hyperref{contract.pdf}{}{claim-protocol-compiler/req:witness-refusal-meaning}{Requirement~13.4}}%
\begin{digestrequirement}[Reporting a refusal for want of a feasibility witness]
The engine never reports a refusal with \iface{no_feasibility_witness} to a
user as a defect in the protocol.
\end{digestrequirement}

% claim protocol-compiler/req:physical-family-decision
\digestlabel{req:physical-family-decision}{13.5}\pdfbookmark[2]{Requirement 13.5 Decision on the presented physical receiver family}{digest-protocol-compiler/req:physical-family-decision}
\renewcommand*{\RestatedName}{\hyperref{contract.pdf}{}{claim-protocol-compiler/req:physical-family-decision}{Requirement~13.5}}%
\begin{digestrequirement}[Decision on the presented physical receiver family]
On every route the engine decides, on the declared numbers, the condition
clause~(iv) of Definition~\ref{def:well-formed-spec} places on the physical
receiver family as presented, in the general format on each
measured factor's: presented directly, (S), (H), (P) and (C) of
Definition~\ref{def:exact} on its physical receiver space as one block, in
the declared order of its alphabet; presented as a
choice-weighted family, condition~(S) on its own declared-order sum,
after the decision of Requirement~\ref{req:conditional-family-decision} on
its conditional families and the representability decision on the products
formed from them. It decides it before any condition of
Definition~\ref{def:chaincond} is decided on the chain, and refuses a failure
with \iface{physical_family_not_admissible}, naming the measured factor, the
predicate and the index of the first failure (for (S), the index of the
first inexact addition). The decision stops at the first failure.
\end{digestrequirement}

% claim protocol-compiler/req:chain-refusals
\digestlabel{req:chain-refusals}{13.6}\pdfbookmark[2]{Requirement 13.6 Refusals of the interface conditions on a chain}{digest-protocol-compiler/req:chain-refusals}
\renewcommand*{\RestatedName}{\hyperref{contract.pdf}{}{claim-protocol-compiler/req:chain-refusals}{Requirement~13.6}}%
\begin{digestrequirement}[Refusals of the interface conditions on a chain]
The engine decides the interface conditions on a chain of
Definition~\ref{def:chaincond} where the chain is consumed, before a
declaration exists, and refuses a chain that fails one; no refusal emits a
declaration.
\begin{enumerate}[label=(\roman*),leftmargin=2.2em]
\item A chain failing one of (C1)--(C5) or (S1)--(S4) is refused with
  \iface{chain_interface_condition_failed}, naming the condition and, for
  (C1)--(C5) in the general format, the factor, and, for (C4) and (S3), the outcome or
  signal index at which the first entry disagreed. A non-empty
  source segment failing one of (A1)--(A3) is refused.
\item Beside a non-empty source segment, an entry declared general is refused
  with \iface{general_entry_beside_source_segment}, naming the entry.
\item A declaration carrying a non-empty source segment in
  untrusted-source mode is refused with
  \iface{source_segment_in_untrusted_mode}.
\end{enumerate}
\end{digestrequirement}

% claim protocol-compiler/req:band-a-ideal-receiver
\digestlabel{req:band-a-ideal-receiver}{13.7}\pdfbookmark[2]{Requirement 13.7 Band A only within the ideal qubit receiver}{digest-protocol-compiler/req:band-a-ideal-receiver}
\renewcommand*{\RestatedName}{\hyperref{contract.pdf}{}{claim-protocol-compiler/req:band-a-ideal-receiver}{Requirement~13.7}}%
\begin{digestrequirement}[Band A only within the ideal qubit receiver]
The engine recognises the Band A route on the record, reading the
shape that condition (b) of Definition~\ref{def:banda} gives the chain and
the sector clause of its (a), for a specification whose source segment is
empty, whose emitted sector set (in the general format, the joint one) is a singleton, and
in which, for some measured factor, the receiver segment is one
link with status exact, equal source and
target descriptors, a channel that, where the link
supplies one as data, is exactly the identity map on its system, an empty
constraint slot and an empty licence slot; it decides the link condition on
the chain record of every measured factor. The clauses of (b) this
recognition leaves open --- that the link supplies its channel as data, its
outcome correspondence and its target POVM --- are
decided by Requirement~\ref{req:band-a-identity-link}. It refuses such a
specification, before a declaration is emitted, with
\iface{band_a_outside_ideal_receiver}, naming the first clause below that
fails and, for (iii), the choice value where there is one and the outcome
index of the first element that is not a projector, unless
\begin{enumerate}[label=(\roman*),leftmargin=2.2em]
\item the specification declares exactly one measured factor;
\item that factor's physical receiver system has dimension two and, in
  trusted-source mode, the transmitted system of the
  transmitter artefact (Definition~\ref{def:spec}\,(a)) has dimension two; and
\item that factor's physical receiver family as presented is projective,
  decided exactly on the declared numbers: presented directly, every element
  satisfies $\Fph_w\Fph_w=\Fph_w$; presented as a
  choice-weighted family, every element of every declared conditional
  family satisfies $N_{u|t}N_{u|t}=N_{u|t}$
  (Definition~\ref{def:weighted}).
\end{enumerate}
\end{digestrequirement}

% claim protocol-compiler/req:band-a-identity-link
\digestlabel{req:band-a-identity-link}{13.8}\pdfbookmark[2]{Requirement 13.8 Band A only on an identity link that delivers the declared family}{digest-protocol-compiler/req:band-a-identity-link}
\renewcommand*{\RestatedName}{\hyperref{contract.pdf}{}{claim-protocol-compiler/req:band-a-identity-link}{Requirement~13.8}}%
\begin{digestrequirement}[Band A only on an identity link that delivers the declared family]
For a specification that Requirement~\ref{req:band-a-ideal-receiver}
recognises on the Band A route and does not refuse, the engine
decides on the chain record of its one measured factor, before a
declaration is emitted, that
\begin{enumerate}[label=(\roman*),leftmargin=2.2em]
\item the one link of that factor's receiver segment supplies its
  channel as data, which that recognition has then decided to be exactly the
  identity map on its system; and
\item the outcome correspondence of that segment is the identity: the
  physical receiver alphabet $\Vph$ and the alphabet $\Vset$ of the link's
  target POVM are one set, and $\gamma(w)=w$ for every
  $w\in\Vph$;
\end{enumerate}
and refuses a specification failing either, before a declaration is emitted,
with \iface{band_a_link_not_identity}, naming the first clause that fails. On
a specification passing both, condition (C4) of
Definition~\ref{def:chaincond} is decided, as on every chain whose links
supply their channel as data, and reads $F_v=\Fph_v$ for every $v\in\Vset$,
entry for entry, $\{\Fph_w\}$ being the physical receiver family as presented
and, for a choice-weighted family, the family of products
Definition~\ref{def:weighted} forms; a failure of it is refused as every
failure of (C4) is (Requirement~\ref{req:chain-refusals}). A declaration
emitted from a specification so recognised therefore satisfies
Definition~\ref{def:banda}\,(b) on its record --- its one link the identity,
supplied as data, its outcome correspondence the identity and its target
POVM the declared family --- and the rules that read the route, the ledger's
among them (Requirement~\ref{req:link-row-provenance}), read it of such a
declaration alone.
\end{digestrequirement}

% claim protocol-compiler/req:chain-record-slots
\digestlabel{req:chain-record-slots}{13.9}\pdfbookmark[2]{Requirement 13.9 Emission of both chain segments and every link slot}{digest-protocol-compiler/req:chain-record-slots}
\renewcommand*{\RestatedName}{\hyperref{contract.pdf}{}{claim-protocol-compiler/req:chain-record-slots}{Requirement~13.9}}%
\begin{digestrequirement}[Emission of both chain segments and every link slot]
The engine emits the chain record (\iface{reduction_chain_record}) with both
of its segments, the source segment and the receiver segment, and on
every link with both its constraint slot and its licence slot, each whether
or not it is empty. This includes Band A, where the source segment is
empty and the receiver segment is one identity link with both slots empty.
\end{digestrequirement}

% claim protocol-compiler/req:link-row-provenance
\digestlabel{req:link-row-provenance}{13.10}\pdfbookmark[2]{Requirement 13.10 Omission of the link-soundness row only on a proved chain}{digest-protocol-compiler/req:link-row-provenance}
\renewcommand*{\RestatedName}{\hyperref{contract.pdf}{}{claim-protocol-compiler/req:link-row-provenance}{Requirement~13.10}}%
\begin{digestrequirement}[Omission of the link-soundness row only on a proved chain]
The engine omits the row of Assumption~\ref{ass:link} from the assumption
ledger of a compilation only
(i) on the Band A route, for a declaration whose record satisfies
Definition~\ref{def:banda}\,(b) as Requirements~\ref{req:band-a-ideal-receiver}
and~\ref{req:band-a-identity-link} decide it --- the one link the identity,
supplied as data, its outcome correspondence the identity and its target
POVM, by the decided (C4), the presented physical family --- where
Theorem~\ref{thm:band-a}\,(v) proves every clause of it; or
(ii) where the source segment is empty and, for every
measured factor of the specification, all of the following hold: its
receiver segment consists of one link, the instance link of an
artefact that the reductions module emitted in the same run, with status
\emph{exact} and no channel supplied as data; the outcome correspondence of
the segment is the identity on that factor's target alphabet; and that
factor's physical receiver family as presented is exactly the compression
$\{\iota^{\dagger}\Mdec_v\iota\}_v$ of that artefact's declared model
onto the presented physical receiver space, $\iota$ its isometric inclusion
into the physical system, or the declared model itself. At several measured
factors the omission rests in addition on A-MEAS
(\xref{reductions/ass:meas})\,(iv), that the round's physical measurement is
the product of the factors' measurements. On either route, at a factor
whose delivered family is weighted after its segment, in the
setting-dependent form of Definition~\ref{def:weighted}, the omission
rests in addition on \xref{reductions/lem:linear}\,(i): the weighting of the
family the segment's link delivers and the identically weighted physical
family stand in the squashing relation under the same channel, so the
transfer applies to the weighted protocol as to the unweighted one; on
route~(i) the channel is the identity and the two weighted families are the
same matrices (Theorem~\ref{thm:band-a}\,(v)).
That each link is that artefact's instance link and that each presented
family is that compression is established by the part of the engine that ran
the reductions module and docked its artefacts, only for an instance for which
Remark~\ref{C-rem:obligations-discharged} names the claim showing the docked
family to be that compression, and is passed to the compiler
beside the specification (\iface{chain_provenance}); the compiler never reads
it from the specification. The compiler decides, on the presented chain, the
parts of~(ii) that are data: for every measured factor the length of its
segment, the status tag, the channel slot and the outcome correspondence
(Definition~\ref{def:chaincond}, (C2)), and the length of the source segment.
Where condition~(ii) is not passed, or a part of it decided on the chain
fails, the row is emitted.
\end{digestrequirement}

% claim protocol-compiler/req:reference-companion-compiled
\digestlabel{req:reference-companion-compiled}{13.11}\pdfbookmark[2]{Requirement 13.11 The reference companion compiled beside the declaration}{digest-protocol-compiler/req:reference-companion-compiled}
\renewcommand*{\RestatedName}{\hyperref{contract.pdf}{}{claim-protocol-compiler/req:reference-companion-compiled}{Requirement~13.11}}%
\begin{digestrequirement}[The reference companion compiled beside the declaration]
Where the last link of the receiver segment of one or more
measured factors was emitted with a reference
companion (\xref{reductions/req:reference-companion}), the engine emits,
beside the declaration $\Decl$ it compiles, one \emph{compiled reference
companion} $\Decl^{\mathrm{ref}}$, in \iface{reference_companion}:
\begin{enumerate}[label=(\roman*),leftmargin=2.4em]
\item it is the tuple the compilation map (Definition~\ref{def:compile})
  gives from the same specification with, for every measured factor whose
  last link was so emitted, that link's target family replaced by its
  companion's and that link's entries replaced by its companion's relaxed
  entries (\xref{reductions/lem:row-relaxation}), those of the entries
  requested of the link included where the companion carries them, and none
  in place of a requested entry the companion does not carry
  (\xref{reductions/req:reference-companion}\,(iii)); the engine forms it by
  the code path that forms $\Decl$, every other slot read from the same
  data, the family and the entries of every other measured factor's last
  link among them;
\item where a measured factor's family is weighted after its chain,
  by that factor's weighting (Definition~\ref{def:weighted}), every product
  $P(f\,|\,t(v))\,F_v$ is
  formed exactly from the companion's enclosure: the centre multiplied by
  the weight and every radius by its modulus, so that the result encloses
  the product of the weight and the companion's element. Where the companion
  carries a factor enclosure (\xref{foundations/def:factor-enclosure})
  of a block of $F_v$ (\xref{reductions/req:reference-companion}\,(vii)),
  the block of the product carries, beside its entrywise enclosure, the one of
  Lemma~\ref{C-lem:weighted-factor}: the same box, its centres and radii
  unaltered, and the core multiplied by the weight in exact rational
  arithmetic. Where the family is not weighted after its chain, every factor
  enclosure the companion carries is passed to $\Decl^{\mathrm{ref}}$
  unaltered;
\item the engine decides on $\Decl^{\mathrm{ref}}$ none of the predicates
  it decides on the binary64 numbers of $\Decl$: no Hermiticity, positivity
  or completeness of the family, no completeness sum in a declared order
  (Definition~\ref{def:exact}), no chain condition and no check of the
  validator; it passes the companion's entries through as
  Requirement~\ref{req:pass-through} passes every entry;
\item it emits no bundle of a run of its own from $\Decl^{\mathrm{ref}}$
  and marks it evaluation-only, a declaration with no run of its own, with
  a reference to the declaration $\Decl$ it is the
  companion of and the companion's record carried unchanged.
\end{enumerate}
Where no measured factor's last link carries a companion, nothing of this
requirement applies, and the run forms its first-order number on $\Decl$
alone, its blocks resting as the records of the artefact say.
\end{digestrequirement}

\section*{14\quad Requirements: the emission gates}\pdfbookmark[1]{14 Requirements: the emission gates}{digest-group-protocol-compiler/req:eligibility-flag}

% claim protocol-compiler/req:eligibility-flag
\digestlabel{req:eligibility-flag}{14.1}\pdfbookmark[2]{Requirement 14.1 The key-length eligibility flag}{digest-protocol-compiler/req:eligibility-flag}
\renewcommand*{\RestatedName}{\hyperref{contract.pdf}{}{claim-protocol-compiler/req:eligibility-flag}{Requirement~14.1}}%
\begin{digestrequirement}[The key-length eligibility flag]
The engine emits the flag \iface{key_length_eligible} true for a declaration
only if that declaration is key-length-eligible
(Definition~\ref{def:eligible}).
\end{digestrequirement}

% claim protocol-compiler/req:gate-outcomes
\digestlabel{req:gate-outcomes}{14.2}\pdfbookmark[2]{Requirement 14.2 Outcome of each gate failure}{digest-protocol-compiler/req:gate-outcomes}
\renewcommand*{\RestatedName}{\hyperref{contract.pdf}{}{claim-protocol-compiler/req:gate-outcomes}{Requirement~14.2}}%
\begin{digestrequirement}[Outcome of each gate failure]
For each gate failure the engine produces the outcome the gate outcome table
(Definition~\ref{def:gate-outcome-table}) fixes for it, and no other; where
the failures of one declaration meet the premises of a refusal row and of an
emission row both, it produces the refusal:
\begin{enumerate}[label=(\roman*),leftmargin=2.2em]
\item row~1, a check that fails: a refusal naming the check, which in the
  general format covers \iface{ensemble_form_in_untrusted_mode},
  \iface{alice_ensemble_marginal_mismatch}, \iface{alice_ensemble_sum_rounds},
  \iface{alice_sector_marginal_zero} and
  \iface{constraint_scope_not_a_factor};
\item row~3, trusted-source mode at any $d_A$ with a certificate present and
  failing a condition of gate (G2) for its form: a refusal naming the
  condition, an ensemble
  presentation beside a POVM-form declaration being refused with
  \iface{ensemble_presentation_without_ensemble_form};
\item rows~5 and~6, a constraint list not in the licensed class or the
  attack-set licence record \iface{attack_set_license} absent: an emission
  with \iface{key_length_eligible} false.
\end{enumerate}
\emph{A trusted-mode specification at $d_A\ge2$ whose certificate is absent
does not reach this module.} The mode is read off the transmitter artefact's
marginal slot (Definition~\ref{def:compile}), and the transmitter
refuses a trusted-mode request at $d_A\ge2$ that carries no certificate before
any artefact is emitted, with \iface{trusted_mode_without_certificate}
(\xref{transmitter/req:licence-scope}; the modes are defined in
\xref{transmitter/def:source-spec}). So an artefact in trusted mode at
$d_A\ge2$ carries the record gate (G2) tests for, and the two rows of the
table that describe a certificate-absent case --- rows 2 and 4, at $m=0$ and
at $m\ge1$ --- are \emph{unreachable through the transmitter artefact}. At
$d_A=1$ the gate demands no certificate at all, the slot holds the $[1]$ the
transmitter filled it with (\xref{transmitter/req:marginal-fill}), and one
supplied there and failing is refused by row 3, its scope being every $d_A$.
Rows 2 and 4 are unreachable in the general format for the same reason, which
reads the transmitter's refusal and no property of the general format. The
engine never supplies a missing attestation and never proceeds as though it
had one.
\end{digestrequirement}

% claim protocol-compiler/req:attestation-presence
\digestlabel{req:attestation-presence}{14.3}\pdfbookmark[2]{Requirement 14.3 Attestations tested for presence only}{digest-protocol-compiler/req:attestation-presence}
\renewcommand*{\RestatedName}{\hyperref{contract.pdf}{}{claim-protocol-compiler/req:attestation-presence}{Requirement~14.3}}%
\begin{digestrequirement}[Attestations tested for presence only]
The engine tests the attestations of Definition~\ref{def:eligible} for
presence only, never for truth: at gate (G2) the clause (SR3) of the
source-replacement licence record \iface{source_replacement_license}
(\xref{transmitter/def:srcert}), in every form of the certificate, and at gate
(G3)(b) the attack-set licence record \iface{attack_set_license}, the
attestation of Assumption~\ref{ass:atk}. It never supplies an attestation and
never verifies one.
\end{digestrequirement}

% claim protocol-compiler/req:vacuous-gate-report
\digestlabel{req:vacuous-gate-report}{14.4}\pdfbookmark[2]{Requirement 14.4 A vacuous pass recorded as vacuous}{digest-protocol-compiler/req:vacuous-gate-report}
\renewcommand*{\RestatedName}{\hyperref{contract.pdf}{}{claim-protocol-compiler/req:vacuous-gate-report}{Requirement~14.4}}%
\begin{digestrequirement}[A vacuous pass recorded as vacuous]
Where gate (G3)(a) of Definition~\ref{def:eligible} holds vacuously, in
untrusted-source mode at any $d_A$, the engine records it in the gate report
\iface{gate_report} as \emph{passed vacuously} rather than \emph{passed}, and
it never reports the vacuous pass as evidence of anything.
\end{digestrequirement}

\bibliographystyle{plain}
\bibliography{../../../refs/references}

\end{document}
